% Utilisation des données sismiques pour concevoir la structure du globe terrestre — SVT 1ère TI — Cours signé OnBuch+
% Programme officiel MINESEC. Compiler avec : tectonic NCT09-donnees-sismiques-structure-globe.tex
\newcommand{\DOCMATIERE}{SVT}
\newcommand{\DOCNIVEAU}{1ère TI}
\newcommand{\DOCMODULE}{Module 3 : L'Éducation à l'Environnement et au Développement Durable}
\newcommand{\DOCLECON}{NCT09}
\newcommand{\DOCTITRE}{Utilisation des données sismiques pour concevoir la structure du globe terrestre}
% OnBuch+ — préambule « cours signés » ENRICHI (compilé avec tectonic / XeLaTeX)
% Système visuel « Pop » d'OnBuch+ : fond crème (jamais blanc), bordures encre
% épaisses, ombres « dures » décalées, titres Archivo Black en capitales.
% Version MATHÉMATIQUES : graphes (pgfplots/tikz), boîtes pédagogiques
% (définition, propriété, méthode, attention, le savais-tu, exercice résolu,
% exercice, TP, bilan), page de garde et signature OnBuch+ ancrée en bas.
%
% Macros attendues dans main.tex AVANT \input{preamble} :
%   \DOCMATIERE  \DOCNIVEAU  \DOCMODULE  \DOCLECON (numéro)  \DOCTITRE
\documentclass[11pt]{article}

\usepackage{fontspec}
\usepackage[french]{babel}
\usepackage[a4paper,margin=2cm,top=3.4cm,bottom=2.8cm,headheight=1.4cm,footskip=1.3cm]{geometry}
\usepackage{amsmath,amssymb,amsfonts}
\usepackage{mathtools} % \xmapsto, \coloneqq... souvent utilisés par les modèles
\usepackage{cancel} % simplifications barrées (bilans de forces)
% \abs{z} (module, valeur absolue) et \norm{u} (norme), utilisés par les modèles
\providecommand{\abs}{}\let\abs\relax\DeclarePairedDelimiter{\abs}{\lvert}{\rvert}
\providecommand{\norm}{}\let\norm\relax\DeclarePairedDelimiter{\norm}{\lVert}{\rVert}
\usepackage[table]{xcolor} % [table] : fournit aussi \rowcolors (alternance de lignes)
\usepackage{pagecolor}
\usepackage{tikz}
\usetikzlibrary{arrows.meta,positioning,calc,shapes.geometric,shapes.misc,shapes.arrows,decorations.pathmorphing,decorations.pathreplacing,decorations.markings,patterns,shadows,mindmap,trees,fit,backgrounds,matrix,intersections,angles,quotes}
\usepackage{pgfplots}
\usepackage[version=4]{mhchem} % équations nucléaires \ce{^{235}_{92}U}
\usepackage[european,siunitx]{circuitikz} % schémas électriques
\pgfplotsset{compat=1.18}
\usepgfplotslibrary{fillbetween,groupplots} % aires entre courbes (calcul intégral)
\usepackage{enumitem}
\usepackage{fancyhdr}
% [most] charge toutes les bibliothèques tcolorbox (listings, minted, raster,
% documentation...) et gonfle inutilement la mémoire TeX sur un document long
% (~300 boîtes) : on ne charge que ce qui est réellement utilisé.
\usepackage[skins,breakable]{tcolorbox}
\usepackage{graphicx}
\usepackage{colortbl}
\usepackage{booktabs}
\usepackage{tabularx}
\usepackage{diagbox} % case d'en-tête diagonale des tableaux à double entrée
% Colonnes extensibles centrée / gauche / droite (C, L, R), souvent utilisées par les modèles
\newcolumntype{C}{>{\centering\arraybackslash}X}
\newcolumntype{L}{>{\raggedright\arraybackslash}X}
\newcolumntype{R}{>{\raggedleft\arraybackslash}X}
\usepackage{array}
\usepackage{multicol}
\usepackage{siunitx}
% parse-numbers=false : accepte aussi \SI{2,5 \times 10^{-3}}{...}
\sisetup{locale=FR,output-decimal-marker={,},group-minimum-digits=4,parse-numbers=false}
\DeclareSIUnit\ohm{\ensuremath{\symup{\Omega}}} % U+2126 absent de la police
\DeclareSIUnit\division{div} % réglages d'oscilloscope (ms/div, V/div)
\DeclareSIUnit\tr{tr} % tours (vitesse de rotation, ex. \SI{1500}{\tr\per\minute})
\DeclareSIUnit\ch{ch} % cheval-vapeur (puissance moteur, unité usuelle hors SI)
\DeclareSIUnit\franccfa{FCFA} % franc CFA (coûts, contexte camerounais)
\DeclareSIUnit\dioptre{\ensuremath{\delta}} % vergence d'une lentille (dioptrie)
\usepackage{qrcode}
% \degree : commande fréquemment utilisée par les modèles pour un angle ou une
% température en dehors de \SI{}{} (non fournie par les paquets chargés).
\providecommand{\degree}{^{\circ}}
% \qed (fin de démonstration), utilisé par les modèles sans amsthm
\providecommand{\qed}{\hfill$\square$}
% \barre : double barre des valeurs interdites dans un tableau de variations (tkz-tab)
\providecommand{\barre}{\ensuremath{\|}}
% environnement proof (amsthm non chargé) utilisé par les modèles
\ifdefined\proof\else\newenvironment{proof}[1][Démonstration]{\par\noindent\textit{#1.}\ }{\qed\par}\fi
\usepackage{lastpage}
\usepackage{hyperref}
\hypersetup{hidelinks}

% ============================================================
% Palette « Pop » OnBuch+ — mêmes hex que la classe P (plus_ui.dart)
% ============================================================
\definecolor{poporange}{HTML}{F59321}
\definecolor{popdark}{HTML}{E07A0C}
\definecolor{popcream}{HTML}{FFF6E8}
\definecolor{popink}{HTML}{1C1714}
\definecolor{poppurple}{HTML}{7A5AE0}
\definecolor{popgreen}{HTML}{2E9E6A}
\definecolor{poppink}{HTML}{E0457B}
\definecolor{popblue}{HTML}{2F7DE1}
\definecolor{popgold}{HTML}{C98A12}
\definecolor{popmuted}{HTML}{8A7F76}
\definecolor{popline}{HTML}{EADFD2}
\definecolor{popgray}{HTML}{8A7F76}
\colorlet{popgrey}{popgray}
% Alias tolérants (le modèle nomme parfois les couleurs différemment)
\colorlet{popwhite}{white}
\colorlet{popblack}{popink}
\colorlet{popred}{poppink}
\colorlet{popyellow}{popgold}
% Teintes claires pour fonds de tableaux / zones
\colorlet{popblueL}{popblue!10!white}
\colorlet{poporangeL}{poporange!14!white}
\colorlet{popgreenL}{popgreen!12!white}
\colorlet{poppurpleL}{poppurple!12!white}
\colorlet{poppinkL}{poppink!12!white}
\colorlet{popgoldL}{popgold!14!white}

\pagecolor{popcream}

% ============================================================
% Typographie
% ============================================================
\defaultfontfeatures{Ligatures=TeX}
\setmainfont{PlusJakartaSans-Medium.ttf}[Path=../../../fonts/,BoldFont=PlusJakartaSans-Bold.ttf]
% Maths en sans-serif (Fira Math) avec chiffres et lettres droites de Plus
% Jakarta, pour que les formules se fondent dans le texte.
\usepackage[math-style=ISO,bold-style=ISO]{unicode-math}
\setmathfont{FiraMath-Regular.otf}[Path=../../../fonts/]
\setmathfont{PlusJakartaSans-Medium.ttf}[Path=../../../fonts/,range={up/num,up/{Latin,latin},\mathup}]
\usepackage{icomma} % virgule décimale sans espace en mode math
% Caractères mathématiques Unicode absents des polices de texte (Plus Jakarta,
% Archivo Black) : rendus par leur équivalent mathématique, en texte comme en
% formule — sinon ils s'affichent en carré vide (ex. « Arithmétique dans ℤ »).
\usepackage{newunicodechar}
\newunicodechar{ℝ}{\ensuremath{\mathbb{R}}}
\newunicodechar{ℤ}{\ensuremath{\mathbb{Z}}}
\newunicodechar{ℂ}{\ensuremath{\mathbb{C}}}
\newunicodechar{ℕ}{\ensuremath{\mathbb{N}}}
\newunicodechar{ℚ}{\ensuremath{\mathbb{Q}}}
\newunicodechar{𝔻}{\ensuremath{\mathbb{D}}}
\newunicodechar{α}{\ensuremath{\alpha}}
\newunicodechar{β}{\ensuremath{\beta}}
\newunicodechar{γ}{\ensuremath{\gamma}}
\newunicodechar{δ}{\ensuremath{\delta}}
\newunicodechar{ε}{\ensuremath{\varepsilon}}
\newunicodechar{θ}{\ensuremath{\theta}}
\newunicodechar{λ}{\ensuremath{\lambda}}
\newunicodechar{μ}{\ensuremath{\mu}}
\newunicodechar{σ}{\ensuremath{\sigma}}
\newunicodechar{τ}{\ensuremath{\tau}}
\newunicodechar{φ}{\ensuremath{\varphi}}
\newunicodechar{ψ}{\ensuremath{\psi}}
\newunicodechar{ω}{\ensuremath{\omega}}
\newunicodechar{Σ}{\ensuremath{\Sigma}}
% majuscules produites par \MakeUppercase dans les titres (α-aminés -> Α)
\newunicodechar{Α}{\ensuremath{\alpha}}
\newunicodechar{Β}{\ensuremath{\beta}}
\newunicodechar{Γ}{\ensuremath{\Gamma}}
\newunicodechar{Δ}{\ensuremath{\Delta}}
\newunicodechar{Ε}{\ensuremath{\varepsilon}}
\newunicodechar{Θ}{\ensuremath{\Theta}}
\newunicodechar{Λ}{\ensuremath{\Lambda}}
\newunicodechar{Μ}{\ensuremath{\mu}}
\newunicodechar{Τ}{\ensuremath{\tau}}
\newunicodechar{Φ}{\ensuremath{\Phi}}
\newunicodechar{Ψ}{\ensuremath{\Psi}}
\newunicodechar{Ω}{\ensuremath{\Omega}}
\newunicodechar{↦}{\ensuremath{\mapsto}}
\newunicodechar{∈}{\ensuremath{\in}}
\newunicodechar{∉}{\ensuremath{\notin}}
\newunicodechar{∧}{\ensuremath{\wedge}}
\newunicodechar{⇔}{\ensuremath{\Leftrightarrow}}
\newunicodechar{⟺}{\ensuremath{\Longleftrightarrow}}
\newunicodechar{⇒}{\ensuremath{\Rightarrow}}
\newunicodechar{⟹}{\ensuremath{\Longrightarrow}}
\newunicodechar{∘}{\ensuremath{\circ}}
\newunicodechar{∪}{\ensuremath{\cup}}
\newunicodechar{∩}{\ensuremath{\cap}}
\newunicodechar{≡}{\ensuremath{\equiv}}
\newunicodechar{⊂}{\ensuremath{\subset}}
\newunicodechar{⊥}{\ensuremath{\perp}}
\newunicodechar{∃}{\ensuremath{\exists}}
\newunicodechar{∀}{\ensuremath{\forall}}
\newunicodechar{∛}{\ensuremath{\sqrt[3]{\,}}}
\newunicodechar{∜}{\ensuremath{\sqrt[4]{\,}}}
\newunicodechar{⃗}{\ensuremath{{}^{\to}}}
\newunicodechar{ⁿ}{\ifmmode^{n}\else\textsuperscript{\ensuremath{n}}\fi}
\newunicodechar{ˣ}{\ifmmode^{x}\else\textsuperscript{\ensuremath{x}}\fi}
\newunicodechar{ᵇ}{\ifmmode^{b}\else\textsuperscript{\ensuremath{b}}\fi}
\newunicodechar{ʳ}{\ifmmode^{r}\else\textsuperscript{\ensuremath{r}}\fi}
\newunicodechar{ᵘ}{\ifmmode^{u}\else\textsuperscript{\ensuremath{u}}\fi}
\newunicodechar{ʸ}{\ifmmode^{y}\else\textsuperscript{\ensuremath{y}}\fi}
\newunicodechar{ᵃ}{\ifmmode^{a}\else\textsuperscript{\ensuremath{a}}\fi}
\newunicodechar{ᵉ}{\ifmmode^{e}\else\textsuperscript{\ensuremath{e}}\fi}
\newunicodechar{ᵏ}{\ifmmode^{k}\else\textsuperscript{\ensuremath{k}}\fi}
\newunicodechar{ᵖ}{\ifmmode^{p}\else\textsuperscript{\ensuremath{p}}\fi}
\newunicodechar{ᴺ}{\ifmmode^{N}\else\textsuperscript{\ensuremath{N}}\fi}
\newunicodechar{ᶜ}{\ifmmode^{c}\else\textsuperscript{\ensuremath{c}}\fi}
\newunicodechar{ʰ}{\ifmmode^{h}\else\textsuperscript{\ensuremath{h}}\fi}
\newunicodechar{ᵠ}{\ifmmode^{q}\else\textsuperscript{\ensuremath{q}}\fi}
\newunicodechar{ₙ}{\ifmmode_{n}\else\textsubscript{\ensuremath{n}}\fi}
\newunicodechar{ₐ}{\ifmmode_{a}\else\textsubscript{\ensuremath{a}}\fi}
\newunicodechar{ᵢ}{\ifmmode_{i}\else\textsubscript{\ensuremath{i}}\fi}
\newunicodechar{ᵣ}{\ifmmode_{r}\else\textsubscript{\ensuremath{r}}\fi}
\newunicodechar{ₖ}{\ifmmode_{k}\else\textsubscript{\ensuremath{k}}\fi}
\newunicodechar{ᵦ}{\ifmmode_{\beta}\else\textsubscript{\ensuremath{\beta}}\fi}
\newunicodechar{ₓ}{\ifmmode_{x}\else\textsubscript{\ensuremath{x}}\fi}
\newfontfamily\popdisplay{ArchivoBlack-Regular.ttf}[Path=../../../fonts/]
\newfontfamily\poplogo{SpaceGrotesk-Bold.ttf}[Path=../../../fonts/]
\newfontfamily\popsemi{PlusJakartaSans-SemiBold.ttf}[Path=../../../fonts/]
\newfontfamily\popxbold{PlusJakartaSans-ExtraBold.ttf}[Path=../../../fonts/]
\newfontfamily\popxboldls{PlusJakartaSans-ExtraBold.ttf}[Path=../../../fonts/,LetterSpace=3.0]

\setlist[enumerate]{leftmargin=1.7em,itemsep=5pt,topsep=4pt}
\setlist[itemize]{leftmargin=1.7em,itemsep=4pt,topsep=4pt,label={\color{poporange}\textbullet}}
\setlength{\parindent}{0pt}
\setlength{\parskip}{5pt}
\renewcommand{\arraystretch}{1.25}

% pgfplots : style Pop par défaut
\pgfplotsset{
  every axis/.append style={axis line style={popink,line width=1pt},tick style={popink},
    grid style={popline,line width=0.5pt},label style={font=\small},tick label style={font=\footnotesize}},
  popcurve/.style={poporange,line width=1.8pt,no markers,smooth},
}
% Même style disponible dans un \draw TikZ simple (hors pgfplots)
\tikzset{popcurve/.style={poporange,line width=1.8pt}}
\tikzset{popgrid/.style={popline,very thin}}
\colorlet{popcurve}{poporange} % le nom du style est parfois utilisé comme couleur

% ============================================================
% Pill / squiggle
% ============================================================
\newcommand{\poppill}[3]{%
  \tikz[baseline=-0.55ex]\node[draw=popink,line width=1.2pt,rounded corners=2.6pt,%
    fill=#2,inner xsep=6.5pt,inner ysep=3pt]{%
    \popxboldls\fontsize{7}{7}\selectfont\color{#3}\MakeUppercase{#1}};%
}
\newcommand{\popsquiggle}[1]{%
  \tikz[baseline=-0.4ex]{
    \draw[#1,line width=2.6pt,line cap=round]
      plot [smooth] coordinates {(0,0.09) (0.34,0.01) (0.68,0.21) (1.02,0.01) (1.36,0.10)};
  }%
}
% Mot-clé surligné (marqueur orange sous le texte)
\newcommand{\cle}[1]{{\popxbold\color{popdark}#1}}

% ============================================================
% En-tête / pied
% ============================================================
\pagestyle{fancy}
\fancyhf{}
\renewcommand{\headrulewidth}{0pt}
\renewcommand{\footrulewidth}{0pt}
\lhead{\raisebox{-0.15ex}{\poplogo\fontsize{13}{13}\selectfont\color{popink}OnBuch\color{poporange}+}}
\rhead{\poppill{\DOCMATIERE\ \textbullet\ \DOCNIVEAU\ \textbullet\ Leçon \DOCLECON}{white}{popink}}
\lfoot{\popxbold\fontsize{7.6}{7.6}\selectfont\color{popmuted}\MakeUppercase{Cours signé OnBuch+}\ \textbullet\ {\popsemi\DOCTITRE}}
\rfoot{\tikz[baseline=-0.6ex]\node[rounded corners=8.5pt,fill=popink,minimum height=17pt,minimum width=17pt,inner xsep=5pt,inner sep=0]{\popxbold\fontsize{8}{8}\selectfont\color{popcream}\thepage\,/\,\pageref*{LastPage}};}

% ============================================================
% Titres
% ============================================================
\newcommand{\chaptitre}[2]{%
  \begin{tcolorbox}[enhanced,colback=poporange,colframe=popink,boxrule=2.6pt,arc=11pt,
      drop shadow=popink,left=15pt,right=15pt,top=12pt,bottom=11pt,before skip=3pt,after skip=18pt]
    \poppill{#2}{popink}{popcream}\par\vspace{9pt}
    {\raggedright\hyphenpenalty=10000\exhyphenpenalty=10000\popdisplay\fontsize{22}{24}\selectfont\color{popcream}\MakeUppercase{#1}\par}
  \end{tcolorbox}
}
% Section numérotée : pastille orange avec le numéro + titre Archivo
\newcounter{popsec}
\newcommand{\coursec}[1]{%
  \stepcounter{popsec}\setcounter{popsub}{0}%
  \par\vspace{16pt}\noindent
  \begin{minipage}{\linewidth}
  \tikz[baseline=-0.7ex]\node[draw=popink,line width=1.6pt,fill=poporange,rounded corners=4pt,minimum size=22pt,inner sep=0]{\popdisplay\fontsize{12}{12}\selectfont\color{popcream}\thepopsec};%
  \hspace{8pt}{\popdisplay\fontsize{15}{17}\selectfont\color{popink}\MakeUppercase{#1}}\par
  \vspace{4pt}\noindent{\color{poporange}\rule{36pt}{3.6pt}}
  \end{minipage}\par\nopagebreak\vspace{8pt}
}
\newcounter{popsub}
\newcommand{\courssub}[1]{%
  \stepcounter{popsub}%
  \par\vspace{9pt}\noindent{\popxbold\fontsize{11.5}{13}\selectfont\color{popink}{\color{poporange}\thepopsec.\thepopsub}\hspace{6pt}#1}\par\nopagebreak\vspace{4pt}
}

% ============================================================
% Boîtes pédagogiques — même recette : fond blanc, bordure épaisse colorée,
% ombre dure encre, badge plein. Titre optionnel [#1].
% ============================================================
\tcbset{popbase/.style={breakable,enhanced,colback=white,boxrule=2.3pt,arc=9pt,
  shadow={3pt}{-3pt}{0pt}{popink},left=14pt,right=14pt,top=11pt,bottom=11pt,before skip=11pt,after skip=15pt,
  fontupper=\popsemi\fontsize{10.6}{14.5}\selectfont\color{popink},
  fontlower=\popsemi\fontsize{10.6}{14.5}\selectfont\color{popink}}}
% Coche dessinée (le glyphe ✓ n'existe pas dans les polices du document)
\newcommand{\popcheck}[1]{\tikz[baseline=-0.5ex]\draw[#1,line width=1.6pt,line cap=round,line join=round](-0.13,0.0)--(-0.04,-0.09)--(0.13,0.1);}
\providecommand{\coche}{\popcheck{popgreen}}
\providecommand{\resultat}[1]{\par\smallskip\noindent\fcolorbox{popgreen}{popgreenL}{\parbox{\dimexpr\linewidth-2\fboxsep-2\fboxrule\relax}{\textbf{Résultat.} #1}}\par\smallskip}
\newcommand{\popboxhead}[3]{\poppill{#1}{#2}{white}\ifx\relax#3\relax\else\hspace{6pt}{\popxbold\fontsize{10.6}{12}\selectfont\color{popink}#3}\fi\par\vspace{7pt}}

\newtcolorbox{aretenir}[1][]{popbase,colframe=popgreen,before upper={\popboxhead{À retenir}{popgreen}{#1}}}
\newtcolorbox{exemplebox}[1][]{popbase,colframe=poporange,before upper={\popboxhead{Exemple}{poporange}{#1}}}
\newtcolorbox{definition}[1][]{popbase,colframe=popblue,before upper={\popboxhead{Définition}{popblue}{#1}}}
\newtcolorbox{propriete}[1][]{popbase,colframe=poppurple,before upper={\popboxhead{Propriété}{poppurple}{#1}}}
\newtcolorbox{methode}[1][]{popbase,colframe=popgold,before upper={\popboxhead{Méthode}{popgold}{#1}}}
\newtcolorbox{attention}[1][]{popbase,colframe=poppink,before upper={\popboxhead{Attention}{poppink}{#1}}}
\newtcolorbox{experience}[1][]{popbase,colframe=popblue,colback=popblueL,before upper={\popboxhead{Expérience}{popblue}{#1}}}
% « Le savais-tu ? » : carte violette pleine (contexte, culture, Cameroun)
\newtcolorbox{savaistu}[1][]{popbase,colback=poppurple,colframe=popink,
  fontupper=\popsemi\fontsize{10.4}{14.2}\selectfont\color{white},
  before upper={\poppill{Le savais-tu\,?}{white}{poppurple}\ifx\relax#1\relax\else\hspace{6pt}{\popxbold\color{white}#1}\fi\par\vspace{7pt}}}
% Exercice résolu : énoncé en haut, \tcblower puis la solution sur fond crème
\newcounter{popexo}\newcounter{popexr}
\newtcolorbox{exoresolu}[1][]{popbase,colframe=poporange,colbacklower=popcream,
  segmentation style={popink,line width=1.2pt,dashed},
  before upper={\stepcounter{popexr}\popboxhead{Exercice résolu \thepopexr}{poporange}{#1}},
  before lower={\poppill{Solution}{popink}{popcream}\par\vspace{6pt}}}
% Exercice (non résolu en place) — la correction est dans la partie Corrigés
\newtcolorbox{exercice}[1][]{popbase,colframe=popink,
  before upper={\stepcounter{popexo}\popboxhead{Exercice \thepopexo}{popink}{#1}}}
% Corrigé
\newtcolorbox{corrige}[1][]{popbase,colframe=popgreen,colback=popgreenL,
  before upper={\popboxhead{Corrigé}{popgreen}{#1}}}
% Objectifs / prérequis / bilan
\newtcolorbox{objectifs}{popbase,colframe=popink,colback=poporangeL,
  before upper={\popboxhead{Objectifs de la leçon}{popink}{}}}
\newtcolorbox{prerequis}{popbase,colframe=popink,colback=popblueL,
  before upper={\popboxhead{Prérequis}{popblue}{}}}
\newtcolorbox{bilan}{popbase,colframe=popink,colback=white,boxrule=2.8pt,
  before upper={\popboxhead{Fiche bilan}{poporange}{L'essentiel à connaître pour l'examen}}}
% Figure encadrée avec légende
\newtcolorbox{popfigure}[1][]{enhanced,breakable=false,colback=white,colframe=popline,boxrule=1.4pt,arc=8pt,
  left=10pt,right=10pt,top=10pt,bottom=8pt,before skip=10pt,after skip=12pt,halign=center,
  lower separated=false,halign lower=center,
  fontlower=\popsemi\fontsize{9}{11.5}\selectfont\color{popmuted},
  #1}
\newcommand{\legende}[1]{\par\vspace{6pt}{\popsemi\fontsize{9}{11.5}\selectfont\color{popmuted}#1}}

% ============================================================
% Signature OnBuch+
% ============================================================
\newcommand{\poplogoinline}[1]{{\poplogo\fontsize{#1}{#1}\selectfont\color{popink}OnBuch\color{poporange}+}}
\newcommand{\signatureonbuch}{%
  \begin{tcolorbox}[enhanced,colback=popink,colframe=popink,boxrule=2.2pt,arc=10pt,
      drop shadow=poporange,left=14pt,right=14pt,top=10pt,bottom=10pt,before skip=0pt,after skip=0pt]
    \tikz[baseline=-0.7ex]\node[circle,fill=popgreen,draw=popcream,line width=1.3pt,minimum size=22pt,inner sep=0]{\popcheck{white}};%
    \hspace{9pt}\begin{minipage}[c]{0.62\linewidth}
      {\popxbold\fontsize{12}{13}\selectfont\color{popcream}Signé\ \textbullet\ {\poplogo OnBuch\color{poporange}+}}\par\vspace{2pt}
      {\raggedright\popsemi\fontsize{8}{10.5}\selectfont\color{popline}\MakeUppercase{Équipe pédagogique OnBuch+}\\\MakeUppercase{Conforme au programme officiel MINESEC}\par}
    \end{minipage}\hfill
    {\poplogo\fontsize{22}{22}\selectfont\color{popcream}OnBuch\color{poporange}+}
  \end{tcolorbox}%
}

% ============================================================
% Page de garde
% #1 titre · #2 matière · #3 niveau · #4 module · #5 n° leçon · #6 durée ·
% #7 description / objectifs (texte ou itemize)
% ============================================================
\newcommand{\pagegarde}[7]{%
  \thispagestyle{empty}
  \newgeometry{a4paper,margin=1.7cm,top=1.6cm,bottom=1.4cm}%
  \begin{tikzpicture}[remember picture,overlay]
    \fill[poporange!18] ([xshift=-3.2cm,yshift=-3.6cm]current page.north east) circle (4.6cm);
    \fill[poppurple!14] ([xshift=2.4cm,yshift=7.6cm]current page.south west) circle (3.8cm);
  \end{tikzpicture}%
  {\poplogo\fontsize{30}{30}\selectfont\color{popink}OnBuch\color{poporange}+}\hfill
  \raisebox{0.6ex}{\poppill{Cours signé \textbullet\ Premium}{popink}{popcream}}\par
  \vspace{1pt}\popsquiggle{poppurple}\par
  \vspace{8pt}
  \begin{tcolorbox}[enhanced,colback=poporange,colframe=popink,boxrule=2.8pt,arc=13pt,
      drop shadow=popink,left=18pt,right=18pt,top=12pt,bottom=13pt,after skip=10pt]
    \poppill{#2 \textbullet\ #3}{popink}{popcream}\hspace{5pt}\poppill{#4}{white}{popink}\par\vspace{8pt}
    {\popdisplay\fontsize{14}{14}\selectfont\color{popink}LEÇON #5}\par\vspace{4pt}
    {\raggedright\hyphenpenalty=10000\exhyphenpenalty=10000\popdisplay\fontsize{25}{27}\selectfont\color{popcream}\MakeUppercase{#1}\par}
    \vspace{8pt}
    {\popxbold\fontsize{9.5}{9.5}\selectfont\color{popink}\MakeUppercase{Durée conseillée : #6}}
  \end{tcolorbox}
  \begin{tcolorbox}[enhanced,colback=white,colframe=popink,boxrule=2.2pt,arc=10pt,
      drop shadow=popink,left=16pt,right=16pt,top=10pt,bottom=10pt,after skip=8pt]
    \poppill{Dans cette leçon}{poppurple}{white}\par\vspace{5pt}
    \popsemi\fontsize{9.3}{12.6}\selectfont\color{popink}#7
  \end{tcolorbox}
  \noindent\begin{minipage}[t]{0.487\linewidth}
    \begin{tcolorbox}[enhanced,colback=popgreen,colframe=popink,boxrule=2.2pt,arc=10pt,
        drop shadow=popink,left=11pt,right=11pt,top=10pt,bottom=10pt,height=3.1cm,valign=center]
      \tcbox[colback=white,colframe=popink,boxrule=1.3pt,arc=5pt,boxsep=0pt,nobeforeafter,
        left=3pt,right=3pt,top=3pt,bottom=3pt,valign=center]{\qrcode[height=1.9cm]{https://play.google.com/store/apps/details?id=com.luvvix.onbuch&hl=fr}}%
      \hspace{8pt}\begin{minipage}[c]{0.5\linewidth}
        {\popdisplay\fontsize{11}{12.5}\selectfont\color{white}\MakeUppercase{Télécharge OnBuch}}\par\vspace{4pt}
        {\raggedright\popsemi\fontsize{8.4}{11}\selectfont\color{white}Gratuit \textbullet\ Google Play\\Tuteur IA, annales, concours, résultats\par}
      \end{minipage}
    \end{tcolorbox}
  \end{minipage}\hfill
  \begin{minipage}[t]{0.487\linewidth}
    \begin{tcolorbox}[enhanced,colback=poppurple,colframe=popink,boxrule=2.2pt,arc=10pt,
        drop shadow=popink,left=11pt,right=11pt,top=10pt,bottom=10pt,height=3.1cm,valign=center]
      \tikz[baseline=-0.7ex]\node[circle,fill=white,draw=popink,line width=1.3pt,minimum size=1.6cm,inner sep=0]{\popdisplay\fontsize{24}{24}\selectfont\color{poppurple}+};%
      \hspace{8pt}\begin{minipage}[c]{0.6\linewidth}
        {\popdisplay\fontsize{11}{12.5}\selectfont\color{white}\MakeUppercase{Passe à OnBuch+}}\par\vspace{4pt}
        {\raggedright\popsemi\fontsize{8.4}{11}\selectfont\color{white}Tous les cours signés, Tuteur IA illimité, fiches PDF — onglet OnBuch+ de l'app\par}
      \end{minipage}
    \end{tcolorbox}
  \end{minipage}
  \vspace{8pt}
  \popcontenu
  \vfill
  \signatureonbuch
  \restoregeometry
  \newpage
}

% Rangée « Ce cours contient » (page de garde)
\newcommand{\poptile}[3]{% #1 couleur #2 gros texte #3 légende
  \begin{tcolorbox}[enhanced,colback=#1,colframe=popink,boxrule=2pt,arc=9pt,shadow={2.5pt}{-2.5pt}{0pt}{popink},
      left=6pt,right=6pt,top=9pt,bottom=9pt,halign=center,valign=center,height=1.75cm,width=\linewidth,nobeforeafter]
    {\popdisplay\fontsize{12.5}{13}\selectfont\color{white}\MakeUppercase{#2}}\par\vspace{3pt}
    {\centering\popsemi\fontsize{7.8}{9.5}\selectfont\color{white}#3\par}
  \end{tcolorbox}}
\newcommand{\popcontenu}{%
  \par\noindent{\popxboldls\fontsize{7.5}{7.5}\selectfont\color{popmuted}CE COURS SIGNÉ CONTIENT}\par\vspace{7pt}
  \noindent\begin{minipage}[t]{0.235\linewidth}\poptile{popblue}{Cours}{complet, illustré, pas à pas}\end{minipage}\hfill
  \begin{minipage}[t]{0.235\linewidth}\poptile{poporange}{Méthodes}{et exercices résolus}\end{minipage}\hfill
  \begin{minipage}[t]{0.235\linewidth}\poptile{poppink}{Exercices}{type examen + corrigés}\end{minipage}\hfill
  \begin{minipage}[t]{0.235\linewidth}\poptile{popgreen}{Fiche bilan}{l'essentiel à retenir}\end{minipage}\par}

% Sommaire
\newtcolorbox{sommaire}{enhanced,colback=white,colframe=popink,boxrule=2.2pt,arc=10pt,
  drop shadow=popink,left=18pt,right=18pt,top=13pt,bottom=13pt,before skip=6pt,after skip=20pt,
  before upper={\poppill{Sommaire}{poppurple}{white}\par\vspace{9pt}\popsemi\fontsize{10.6}{14.5}\selectfont\color{popink}}}

% Signature de fin de cours — poussée tout en bas de la dernière page
\newcommand{\findecours}{%
  \par\vfill\nopagebreak
  \begin{center}\popsquiggle{poporange}\end{center}\vspace{4pt}
  \signatureonbuch
}
\AtBeginDocument{\def\llbracket{\mathopen{[\mkern-3.5mu[}}\def\rrbracket{\mathclose{]\mkern-3.5mu]}}} % intervalles d'entiers (glyphe ⟦ absent de la police)
% Glyphes absents de Fira Math : redéfinis à partir de glyphes disponibles
\AtBeginDocument{%
  \def\setminus{\mathbin{\backslash}}\let\smallsetminus\setminus
  \def\vdots{\mathord{\vbox{\baselineskip3.2pt\lineskiplimit0pt\kern4pt\hbox{.}\hbox{.}\hbox{.}}}}%
  \def\ddots{\mathinner{\mkern1mu\raise6.5pt\hbox{.}\mkern2mu\raise3.6pt\hbox{.}\mkern2mu\raise0.7pt\hbox{.}\mkern1mu}}%
  \def\triangle{\mathord{\scalebox{0.9}{$\symup{\Delta}$}}}%
  \def\nabla{\mathord{\raisebox{\depth}{\scalebox{1}[-1]{$\symup{\Delta}$}}}}%
  \def\checkmark{\popcheck{popdark}}%
  \def\star{\mathbin{\ast}}\def\bigstar{\mathord{\ast}}%
  \def\diamond{\mathbin{\scalebox{0.6}{$\Diamond$}}}%
  \def\bigcup{\mathop{\mathchoice{\raisebox{-0.3em}{\scalebox{1.7}{$\cup$}}}{\scalebox{1.25}{$\cup$}}{\cup}{\cup}}\displaylimits}%
  \def\bigcap{\mathop{\mathchoice{\raisebox{-0.3em}{\scalebox{1.7}{$\cap$}}}{\scalebox{1.25}{$\cap$}}{\cap}{\cap}}\displaylimits}%
  \def\lesseqqgtr{\mathrel{\lessgtr}}\def\bigoplus{\mathop{\oplus}\displaylimits}%
}
\providecommand{\ssi}{si et seulement si} % abréviation fréquente
\AtBeginDocument{\providecommand{\wideparen}{\overparen}} % arc de cercle

%% FIGFIX-BEGIN v5 — correctifs globaux des figures (docs/onbuch-generation/tools/figfix)
\usepackage{adjustbox}\usepackage{etoolbox}\tcbuselibrary{raster}
% toute figure TikZ/pgfplots plus large que la ligne est réduite à la largeur disponible
\BeforeBeginEnvironment{tikzpicture}{\begin{adjustbox}{max width=\linewidth}}
\AfterEndEnvironment{tikzpicture}{\end{adjustbox}}
% légendes pgfplots lisibles par-dessus les courbes
\pgfplotsset{every axis/.append style={legend style={fill=white,fill opacity=0.92,text opacity=1,draw=black!15,inner sep=3pt}}}
% étiquettes d'arêtes (syntaxe edge["..."]) avec fond blanc
\tikzset{every edge quotes/.append style={fill=white,inner sep=1.2pt}}
%% FIGFIX-END
\begin{document}

\pagegarde{Utilisation des données sismiques pour concevoir la structure du globe terrestre}{SVT}{1ère TI}{Module 3 : L'Éducation à l'Environnement et au Développement Durable}{NCT09}{4 h}{Cette leçon explore comment l'analyse de la propagation des ondes sismiques (P et S) à travers le globe permet de révéler sa structure interne en couches (croûte, manteau, noyau). Elle développe les compétences d'exploitation de courbes de temps de parcours et de modélisation géophysique, essentielles pour l'épreuve du Probatoire.
\vspace{4pt}
\begin{multicols}{2}\small
\begin{itemize}
\item À la fin de cette leçon, je sais distinguer les ondes P (primaires) et S (secondaires) par leur nature, leur vitesse et leur milieu de propagation.
\item Je sais expliquer le phénomène de réfraction et de réflexion des ondes sismiques au passage d'une discontinuité de milieu.
\item Je sais exploiter une courbe de temps de parcours (temps en fonction de la distance épicentrale) pour identifier les arrivées des ondes directes, réfléchies et réfractées.
\item Je sais déterminer la profondeur d'une discontinuité (ex: Moho, limite manteau/noyau) à partir de données de temps de parcours (méthode du temps réduit ou graphique).
\item Je sais justifier l'existence du noyau liquide externe par l'absence d'ondes S et la zone d'ombre des ondes P.
\item Je sais construire un modèle simplifié en couches de la Terre (croûte, manteau, noyau externe, noyau interne) en précisant l'état physique, la composition chimique dominante et l'épaisseur approximative de chaque enveloppe.
\end{itemize}\end{multicols}}

\begin{sommaire}
\begin{multicols}{2}
\begin{enumerate}
\item 1. Les messagers du profond : ondes P et ondes S
\item 2. Le voyage des ondes : réfraction, réflexion et loi de Snell-Descartes en géophysique
\item 3. Exploitation des courbes de temps de parcours : localiser les discontinuités
\item 4. Le noyau terrestre : la preuve par les ondes (liquide externe, solide interne)
\item 5. Modèle de référence : Structure en couches de la Terre (Modèle PREM / AK135)
\item 6. Au-delà du modèle 1D : Hétérogénéités, Anisotropie et Dynamique (Complément Probatoire)
\item Activité d'intégration
\item Méthodes et exercices résolus
\item Exercices
\item Corrigés des exercices
\item Fiche bilan
\end{enumerate}
\end{multicols}
\end{sommaire}

\begin{objectifs}
\begin{itemize}
\item À la fin de cette leçon, je sais distinguer les ondes P (primaires) et S (secondaires) par leur nature, leur vitesse et leur milieu de propagation.
\item Je sais expliquer le phénomène de réfraction et de réflexion des ondes sismiques au passage d'une discontinuité de milieu.
\item Je sais exploiter une courbe de temps de parcours (temps en fonction de la distance épicentrale) pour identifier les arrivées des ondes directes, réfléchies et réfractées.
\item Je sais déterminer la profondeur d'une discontinuité (ex: Moho, limite manteau/noyau) à partir de données de temps de parcours (méthode du temps réduit ou graphique).
\item Je sais justifier l'existence du noyau liquide externe par l'absence d'ondes S et la zone d'ombre des ondes P.
\item Je sais construire un modèle simplifié en couches de la Terre (croûte, manteau, noyau externe, noyau interne) en précisant l'état physique, la composition chimique dominante et l'épaisseur approximative de chaque enveloppe.
\item Je sais relier la variation de la vitesse des ondes (loi de Snell-Descartes) aux variations de densité et d'élasticité (module d'incompressibilité, module de rigidité) des roches traversées.
\item Je sais argumenter sur les limites du modèle 1D (sphériquement symétrique) face à la réalité 3D (anisotropie, hétérogénéités latérales, panaches mantelliques).
\end{itemize}
\end{objectifs}

\begin{prerequis}
\begin{itemize}
\item Notion d'onde mécanique (propagation, vitesse, fréquence, période) vue en Première (Mécanique/Ondes).
\item Lois de la réflexion et de la réfraction (Snell-Descartes) en optique géométrique (Seconde/Première).
\item Structure de la matière : distinction solide/liquide/gaz, incompressibilité, rigidité (Chimie/Physique Seconde).
\item Géométrie du cercle et trigonométrie de base (calcul d'angles, sinus, cosinus) pour le tracé des rayons.
\item Lecture et interprétation de graphiques (courbes t = f(Δ), axes logarithmiques/linéaires).
\end{itemize}
\end{prerequis}

\newpage

\coursec{Les messagers du profond : ondes P et ondes S}

Imaginez que vous êtes à Yaoundé et qu'un tremblement de terre se produit au large des côtes du Cameroun, près de Kribi. Quelques secondes plus tard, les sismographes de l'observatoire du Mont Fébé enregistrent des vibrations, puis le silence revient. Comment, à partir de ces simples tracés sur papier ou écran, les scientifiques peuvent-ils affirmer avec certitude que le centre de la Terre est une boule de fer liquide à plus de 6000 km sous nos pieds, sans y être jamais allés ? C'est le mystère que nous allons percer en devenant, le temps de cette leçon, des \emph{radiologues} de la planète Terre.

La problématique de cette leçon est la suivante : \textbf{comment des vibrations enregistrées en surface permettent-elles de découvrir la structure interne du globe ?} La réponse tient en deux mots : les \cle{ondes sismiques}, véritables messagers du profond. Encore faut-il apprendre à les connaître et à les distinguer.

\courssub{Qu'est-ce qu'une onde sismique ?}

\textbf{Intuition d'abord.} Frappez un rail de chemin de fer avec une pierre : en collant votre oreille au rail, vous entendez le coup \emph{avant} de l'entendre dans l'air. La déformation créée par le choc s'est propagée dans le métal. De même, lors d'un séisme, la rupture brutale des roches en profondeur (au \cle{foyer} ou hypocentre) libère une énergie qui se propage de proche en proche dans les roches, comme une déformation qui voyage.

\begin{definition}[Onde sismique]
Une \cle{onde sismique} est une \textbf{onde mécanique de déformation} qui se propage dans un \textbf{milieu élastique} (les roches de la Terre). Elle transporte de l'énergie sans transport de matière : chaque particule de roche oscille autour de sa position d'équilibre, puis transmet le mouvement à ses voisines. Le point de départ de la rupture est le \cle{foyer} ; le point de la surface situé à la verticale du foyer est l'\cle{épicentre}.
\end{definition}

Le mot \emph{élastique} est essentiel : après le passage de l'onde, la roche retrouve sa forme initiale, comme un ressort après compression. C'est cette élasticité qui permet la propagation. Sans elle, l'énergie serait dissipée sur place.

\begin{attention}[Ondes de volume ou ondes de surface ?]
Ne confondez pas les \textbf{ondes de volume} (P et S), qui traversent l'intérieur du globe et servent à sonder le profond, avec les \textbf{ondes de surface} (ondes de \emph{Rayleigh} et de \emph{Love}), qui se propagent uniquement le long de la surface terrestre, comme les rides à la surface de l'eau. Ce sont les ondes de surface qui causent l'essentiel des dégâts des séismes, mais ce sont les ondes de volume qui nous renseignent sur la structure interne. Au Probatoire, seules les ondes P et S sont exigibles.
\end{attention}

\courssub{Les ondes P : compression et dilatation}

\textbf{Observation concrète.} Prenez un ressort (ou un « Slinky ») posé sur une table, tenu aux deux extrémités. Donnez une petite poussée brusque dans le sens de la longueur : vous voyez une zone de compression (spires resserrées) suivie d'une zone de dilatation (spires écartées) voyager le long du ressort. Chaque spire oscille \emph{dans le sens de la propagation}. C'est exactement le comportement d'une onde P dans la roche.

\begin{definition}[Onde P]
Une \cle{onde P} (onde \emph{Primaire}) est une onde de \textbf{compression-dilatation}, dite \textbf{longitudinale} : les particules du milieu oscillent \textbf{parallèlement à la direction de propagation}. Elle se propage dans \textbf{tous les milieux élastiques : solides, liquides et gaz}. C'est l'onde sismique la plus rapide : elle arrive donc en premier sur les sismographes, d'où son nom.
\end{definition}

Sa vitesse dépend de deux propriétés mécaniques du milieu et de sa densité :

\begin{definition}[Modules élastiques]
Le \cle{module d'incompressibilité} $K$ (en pascals, Pa) mesure la résistance du milieu à la compression : plus $K$ est grand, plus le milieu est difficile à comprimer. Le \cle{module de rigidité} $\mu$ (ou module de cisaillement, en Pa) mesure la résistance du milieu au cisaillement, c'est-à-dire à la déformation par glissement des couches les unes sur les autres. Un fluide parfait ne résiste pas au cisaillement : $\mu = 0$. La \cle{masse volumique} $\rho$ (en kg/m\textsuperscript{3}) traduit l'inertie du milieu.
\end{definition}

\begin{propriete}[Vitesse des ondes P (admise)]
La vitesse de propagation d'une onde P dans un milieu élastique homogène est :
\[ v_P = \sqrt{\dfrac{K + \dfrac{4}{3}\mu}{\rho}} \]
où $K$ est le module d'incompressibilité, $\mu$ le module de rigidité et $\rho$ la masse volumique. Cette formule est admise ; savoir l'exploiter numériquement est exigible.
\end{propriete}

\courssub{Les ondes S : le cisaillement, signature des solides}

\textbf{Observation concrète.} Reprenez le ressort, mais cette fois secouez une extrémité \emph{latéralement}, de gauche à droite. Une ondulation voyage le long du ressort : chaque spire oscille \emph{perpendiculairement} à la direction de propagation. C'est le comportement d'une onde S. Essayez maintenant de faire la même expérience dans l'eau d'un seau : impossible de faire voyager un cisaillement, l'eau « s'écoule » au lieu de transmettre la déformation.

\begin{definition}[Onde S]
Une \cle{onde S} (onde \emph{Secondaire}) est une onde de \textbf{cisaillement}, dite \textbf{transversale} : les particules du milieu oscillent \textbf{perpendiculairement à la direction de propagation}. Elle ne se propage que dans les \textbf{solides} ($\mu > 0$), car seul un solide peut résister au cisaillement et le transmettre. Plus lente que l'onde P, elle arrive en second sur les sismographes.
\end{definition}

\begin{propriete}[Vitesse des ondes S (admise)]
La vitesse de propagation d'une onde S est :
\[ v_S = \sqrt{\dfrac{\mu}{\rho}} \]
Conséquences fondamentales :
\begin{itemize}
\item Comme $K + \dfrac{4}{3}\mu > \mu$, on a \textbf{toujours} $v_P > v_S$ : les ondes P arrivent avant les ondes S.
\item Dans un \textbf{fluide} (liquide ou gaz), $\mu = 0$, donc $v_S = 0$ : \textbf{les ondes S ne se propagent pas dans les fluides}. L'onde P, elle, subsiste avec $v_P = \sqrt{K/\rho}$.
\end{itemize}
\end{propriete}

\begin{attention}[Le piège classique]
Dire « les ondes S ne passent pas dans les liquides » ne signifie pas qu'elles sont ralenties : elles sont \textbf{totalement bloquées}. C'est cette propriété qui permettra, dans la suite de la leçon, de démontrer que le noyau externe est liquide. Autre piège : ne pas écrire $v_S = \sqrt{K/\rho}$ ; c'est $K$ qui gouverne la compression (ondes P) et $\mu$ qui gouverne le cisaillement (ondes S).
\end{attention}

\begin{popfigure}
\begin{center}
\begin{tikzpicture}[scale=1]
% ---- Onde P ----
\draw[->,thick,popdark] (0,2.6) -- (9.5,2.6) node[right]{\small propagation};
\foreach \x in {0.5,1.0,1.5,2.0,2.5,3.0,3.5,4.0,4.5,5.0,5.5,6.0,6.5,7.0,7.5,8.0,8.5,9.0}{
  \fill[popblue] (\x,2.6) circle (1.6pt);
}
% compression zones
\foreach \x in {2.1,2.35,2.6,6.1,6.35,6.6}{
  \fill[poporange] (\x,2.6) circle (2.2pt);
}
\draw[<->,poporange,thick] (2.35,3.0) -- (3.35,3.0) node[midway,above]{\scriptsize oscillation //};
\node[popdark,font=\bfseries] at (-0.2,2.6) [left]{Onde P};
% ---- Onde S ----
\draw[->,thick,popdark] (0,0) -- (9.5,0) node[right]{\small propagation};
\draw[popgreen,thick,domain=0.5:9.0,samples=200] plot (\x,{0.55*sin(360*\x/3)});
\draw[<->,popgreen,thick] (1.25,0) -- (1.25,0.55) node[midway,right]{\scriptsize oscillation $\perp$};
\node[popdark,font=\bfseries] at (-0.2,0) [left]{Onde S};
\end{tikzpicture}
\end{center}
\legende{Comparaison schématique des deux ondes de volume. En haut : onde P — les particules (points) oscillent parallèlement à la propagation, créant des zones de compression (points resserrés, en orange) et de dilatation. En bas : onde S — les particules oscillent perpendiculairement à la propagation (déformation en cisaillement).}
\end{popfigure}

\courssub{Polarisation des ondes S : SV et SH}

Le mouvement d'une onde S étant perpendiculaire à la propagation, il peut s'orienter de différentes façons dans le plan perpendiculaire : c'est la \cle{polarisation}. On distingue :

\begin{itemize}
\item les ondes \cle{SV} (\emph{Shear Vertical}) : l'oscillation des particules se fait dans le \textbf{plan vertical} contenant la direction de propagation (mouvement « haut-bas » combiné à l'avancée) ;
\item les ondes \cle{SH} (\emph{Shear Horizontal}) : l'oscillation se fait dans le \textbf{plan horizontal}, perpendiculairement à la direction de propagation (mouvement « gauche-droite »).
\end{itemize}

Cette distinction est utile car les ondes SV et SH ne se comportent pas de la même manière aux interfaces entre couches : une onde SV peut se transformer partiellement en onde P à une discontinuité, alors qu'une onde SH ne le peut pas. Pour le Probatoire, retenez simplement que la polarisation des ondes S fournit des informations supplémentaires sur les milieux traversés.

\courssub{Le rapport $v_P/v_S$ : la signature des roches}

Faisons le rapport des deux vitesses :
\[ \dfrac{v_P}{v_S} = \sqrt{\dfrac{K + \dfrac{4}{3}\mu}{\mu}} = \sqrt{\dfrac{K}{\mu} + \dfrac{4}{3}} \]

Pour la plupart des roches de la croûte, on observe expérimentalement que $K \approx \dfrac{5}{3}\mu$ (c'est l'\textbf{hypothèse de Poisson}, valable pour un solide isotrope idéal). Alors :
\[ \dfrac{v_P}{v_S} = \sqrt{\dfrac{5}{3} + \dfrac{4}{3}} = \sqrt{3} \approx 1{,}73 \]

\begin{aretenir}[À retenir]
Dans les roches crustales typiques, $v_P/v_S \approx \sqrt{3} \approx 1{,}73$. Connaître ce rapport permet, à partir de l'écart entre les temps d'arrivée des ondes P et S en une station, d'estimer la distance du séisme — méthode que nous exploiterons dans la suite de la leçon.
\end{aretenir}

\courssub{Atténuation et dispersion : les ondes s'usent en voyageant}

Deux phénomènes qualitatifs modifient les ondes au cours de leur trajet :

\begin{itemize}
\item L'\cle{atténuation géométrique} : l'énergie de l'onde se répartit sur des surfaces (des sphères) de plus en plus grandes à mesure qu'elle s'éloigne du foyer. L'amplitude des vibrations diminue donc avec la distance, comme le son d'un tam-tam s'affaiblit quand on s'éloigne du village.
\item L'\cle{atténuation anélastique} : les roches réelles ne sont pas parfaitement élastiques ; une partie de l'énergie est transformée en chaleur par frottements internes. Ce phénomène est plus marqué dans les milieux chauds ou partiellement fondus (comme l'asthénosphère), ce qui permet aux sismologues de détecter ces zones.
\item La \cle{dispersion} : dans certains milieux, la vitesse de l'onde dépend de sa fréquence ; les différentes fréquences se séparent alors au cours du trajet. Pour les ondes de volume P et S dans le manteau, cet effet reste faible, et on le néglige en première approximation.
\end{itemize}

\courssub{Application numérique type Probatoire}

\begin{methode}[Calculer $v_P$ et $v_S$ connaissant $K$, $\mu$ et $\rho$]
\begin{enumerate}
\item Vérifier que toutes les grandeurs sont en unités SI : $K$ et $\mu$ en pascals (1 GPa = $10^9$ Pa), $\rho$ en kg/m\textsuperscript{3}.
\item Calculer $v_P = \sqrt{\dfrac{K + \frac{4}{3}\mu}{\rho}}$.
\item Calculer $v_S = \sqrt{\dfrac{\mu}{\rho}}$.
\item Convertir le résultat en km/s (diviser par 1000) et vérifier la cohérence : on doit toujours trouver $v_P > v_S$.
\item Vérifier éventuellement que $v_P/v_S \approx 1{,}73$ pour une roche crustale.
\end{enumerate}
\end{methode}

\begin{exoresolu}[Vitesse des ondes dans la croûte continentale]
\textbf{Énoncé.} On donne pour une roche de la croûte continentale : $\rho = 2700$ kg/m\textsuperscript{3}, $K = 50$ GPa, $\mu = 30$ GPa.
\begin{enumerate}
\item Calculer les vitesses $v_P$ et $v_S$ des ondes sismiques dans cette roche.
\item Calculer le rapport $v_P/v_S$ et commenter.
\item Ces ondes pourraient-elles se propager ainsi dans le noyau externe, supposé liquide ? Justifier.
\end{enumerate}
\tcblower
\textbf{Solution détaillée.}

\textbf{1.} Convertissons : $K = 50 \times 10^9$ Pa et $\mu = 30 \times 10^9$ Pa.
\begin{align*}
v_P &= \sqrt{\dfrac{K + \dfrac{4}{3}\mu}{\rho}} = \sqrt{\dfrac{50\times 10^9 + \dfrac{4}{3}\times 30\times 10^9}{2700}} = \sqrt{\dfrac{90\times 10^9}{2700}} \\
&= \sqrt{\num{3,33}\times 10^7} \approx \num{5,77}\times 10^3 \text{ m/s} \approx \num{5,8} \text{ km/s} \\
v_S &= \sqrt{\dfrac{\mu}{\rho}} = \sqrt{\dfrac{30\times 10^9}{2700}} = \sqrt{\num{1,11}\times 10^7} \approx \num{3,33}\times 10^3 \text{ m/s} \approx \num{3,3} \text{ km/s}
\end{align*}

\textbf{2.} $\dfrac{v_P}{v_S} = \dfrac{5,77}{3,33} \approx 1{,}73 \approx \sqrt{3}$. Ce résultat est conforme à l'hypothèse de Poisson : la roche se comporte comme un solide isotrope typique de la croûte.

\textbf{3.} Non pour les ondes S : un liquide possède un module de rigidité nul ($\mu = 0$), donc $v_S = 0$ ; les ondes S ne s'y propagent pas. Seules les ondes P subsisteraient, avec $v_P = \sqrt{K/\rho}$, à une vitesse plus faible que dans un solide de même $K$ car le terme $\frac{4}{3}\mu$ disparaît.
\end{exoresolu}

\begin{center}
\begin{tabularx}{\linewidth}{|l|X|X|}
\hline
\rowcolor{popblueL} \textbf{Critère} & \textbf{Onde P (Primaire)} & \textbf{Onde S (Secondaire)} \\
\hline
Type de déformation & Compression-dilatation & Cisaillement \\
\hline
Type d'onde & Longitudinale (oscillation // propagation) & Transversale (oscillation $\perp$ propagation) \\
\hline
Milieux traversés & Solides, liquides et gaz & Solides uniquement ($\mu > 0$) \\
\hline
Vitesse & $v_P = \sqrt{\dfrac{K + \frac{4}{3}\mu}{\rho}}$ & $v_S = \sqrt{\dfrac{\mu}{\rho}}$ \\
\hline
Ordre d'arrivée & En premier (la plus rapide) & En second ($v_S < v_P$) \\
\hline
Polarisation & Pas de polarisation propre & SV (plan vertical) ou SH (plan horizontal) \\
\hline
\end{tabularx}
\end{center}

\begin{savaistu}[Oldham, 1906 : la liquidité du noyau révélée par les ondes S]
En 1906, le géologue britannique Richard Dixon \textbf{Oldham} remarque qu'au-delà d'une certaine distance de l'épicentre, les sismographes n'enregistrent plus aucune onde S directe : il existe une vaste « zone d'ombre ». Or les ondes S ne traversent pas les liquides. Oldham en déduit que la Terre possède un \textbf{noyau liquide} qui bloque les ondes S : c'est la première preuve géophysique de l'existence du noyau. Quelques années plus tard, l'étude des ondes P réfractées et réfléchies précisera la profondeur de la limite (environ 2900 km, discontinuité de Gutenberg) puis l'existence d'une graine solide au centre (Inge Lehmann, 1936). Sans jamais creuser au-delà d'une dizaine de kilomètres, l'humanité connaissait le cœur de sa planète !
\end{savaistu}

\begin{exercice}[Pour s'entraîner]
Un séisme se produit au large de Kribi. À la station du Mont Fébé, l'onde P arrive à 14 h 02 min 10 s et l'onde S à 14 h 02 min 28 s. On suppose le trajet dans la croûte avec $v_P = 6{,}1$ km/s et $v_S = 3{,}5$ km/s.
\begin{enumerate}
\item Quel est le retard de l'onde S sur l'onde P ?
\item En notant $d$ la distance foyer-station, exprimer les temps de parcours $t_P$ et $t_S$ en fonction de $d$.
\item En déduire la distance $d$ du foyer à la station.
\end{enumerate}
\end{exercice}

\begin{corrige}[Exercice : distance foyer-station]
\textbf{1.} Retard : $\Delta t = t_S - t_P = 18$ s.

\textbf{2.} $t_P = \dfrac{d}{v_P} = \dfrac{d}{6{,}1}$ et $t_S = \dfrac{d}{v_S} = \dfrac{d}{3{,}5}$ (avec $d$ en km et $t$ en s).

\textbf{3.} $\Delta t = t_S - t_P = d\left(\dfrac{1}{3{,}5} - \dfrac{1}{6{,}1}\right) = d \times \dfrac{6{,}1 - 3{,}5}{3{,}5 \times 6{,}1} = d \times \dfrac{2{,}6}{21{,}35} \approx 0{,}1218\,d$.

D'où $d = \dfrac{18}{0{,}1218} \approx 148$ km. Le foyer se situe à environ 148 km de la station du Mont Fébé, ce qui est cohérent avec un séisme au large de Kribi.
\end{corrige}

\begin{aretenir}[L'essentiel de la section]
\begin{itemize}
\item Une onde sismique est une onde mécanique de déformation se propageant dans un milieu élastique, sans transport de matière.
\item Les ondes \textbf{P} sont longitudinales (compression-dilatation), traversent tous les milieux, et sont les plus rapides : $v_P = \sqrt{(K + \frac{4}{3}\mu)/\rho}$.
\item Les ondes \textbf{S} sont transversales (cisaillement), ne traversent que les solides : $v_S = \sqrt{\mu/\rho}$ ; dans un fluide, $v_S = 0$.
\item Toujours $v_P > v_S$ ; dans les roches crustales, $v_P/v_S \approx \sqrt{3} \approx 1{,}73$.
\item L'écart de temps d'arrivée S $-$ P permet de calculer la distance du foyer : c'est la base de la sismologie exploratrice.
\end{itemize}
\end{aretenir}

\coursec{Le voyage des ondes : réfraction, réflexion et loi de Snell-Descartes en géophysique}

Dans la section précédente, nous avons rencontré les deux messagers du profond : les ondes \cle{P}, rapides et compressives, capables de traverser solides et liquides, et les ondes \cle{S}, plus lentes, transversales, prisonnières des milieux solides. Mais une question demeure : lorsque ces ondes plongent dans la Terre, pourquoi ne voyagent-elles pas en ligne droite ? Pourquoi certains séismologues enregistrent-ils des ondes à \SI{10000}{km} de l'épicentre, alors qu'aucune n'arrive à certaines distances intermédiaires ? C'est tout l'objet de cette section : comprendre les \cle{règles du voyage} des ondes sismiques, c'est-à-dire la \cle{réfraction}, la \cle{réflexion} et la grande loi qui les gouverne, la loi de \cle{Snell-Descartes} adaptée à une Terre sphérique.

\courssub{Le principe de Fermat : l'onde est pressée}

\textbf{Une situation concrète pour comprendre.} Imagine un maître-nageur sur la plage de Kribi qui voit un baigneur en difficulté au large. Pour le rejoindre le plus vite possible, il ne court pas en ligne droite vers lui : il court d'abord un peu plus longtemps sur le sable (où il va vite), puis nage une distance plus courte (où il va lentement). Le trajet le plus rapide n'est pas la ligne droite, mais une ligne \emph{brisée} qui tient compte des vitesses différentes dans chaque milieu.

\begin{propriete}[Principe de Fermat (démonstration non exigible)]
Une onde sismique se propage d'un point A à un point B en suivant le trajet qui rend le \cle{temps de parcours minimal} (plus précisément : stationnaire). Ce trajet est appelé \cle{rayon sismique}. Dans un milieu où la vitesse varie, le rayon n'est donc généralement \emph{pas} une ligne droite : il se courbe pour « passer plus de temps » dans les milieux rapides et « moins de temps » dans les milieux lents.
\end{propriete}

\begin{attention}[Rayon et front d'onde]
Ne confonds pas le \cle{front d'onde} (surface atteinte par l'onde à un instant donné) et le \cle{rayon} (trajectoire perpendiculaire aux fronts d'onde, suivie par l'énergie). En sismologie, on raisonne presque toujours sur les rayons, comme en optique géométrique.
\end{attention}

\courssub{La loi de Snell-Descartes : de l'optique plane à la Terre sphérique}

\textbf{Le cas plan (rappel).} Quand une onde passe d'un milieu de vitesse $v_1$ à un milieu de vitesse $v_2$ séparés par une interface plane, les angles d'incidence $i_1$ et de réfraction $i_2$ (mesurés par rapport à la \emph{normale} à l'interface) vérifient :
\[ \frac{\sin i_1}{v_1} = \frac{\sin i_2}{v_2} \]
Conséquence immédiate : si $v_2 > v_1$ (l'onde accélère en profondeur), alors $\sin i_2 > \sin i_1$ : le rayon \emph{s'écarte de la normale}. Si l'onde ralentit, il se rapproche de la normale.

\textbf{Le problème de la sphère.} La Terre n'est pas un empilement de couches planes : c'est une sphère stratifiée en coquilles concentriques (croûte, manteau, noyau). À chaque profondeur, la « normale » à l'interface est le rayon terrestre, qui change de direction en tout point. La loi plane ne suffit plus : il faut une grandeur \emph{conservée} le long de tout le trajet.

\begin{propriete}[Conservation du paramètre de rayon — démonstration guidée exigible]
Dans une Terre sphériquement symétrique (la vitesse $v$ ne dépend que de la distance $r$ au centre), la quantité
\[ p = \frac{r \sin i}{v(r)} \]
appelée \cle{paramètre de rayon}, reste \textbf{constante le long d'un même rayon}, où $i$ est l'angle entre le rayon sismique et la direction radiale (la « verticale » locale).
\end{propriete}

\begin{experience}[Démonstration guidée : pourquoi $p$ est-il constant ?]
Découpons la Terre en fines coquilles sphériques concentriques d'épaisseur $\mathrm{d}r$, dans chacune desquelles la vitesse est quasi constante.
\begin{enumerate}
\item \textbf{Hypothèse de travail} : entre deux coquilles voisines, on peut appliquer localement la loi de Snell-Descartes plane, car l'interface est presque plane à petite échelle.
\item \textbf{Subtilité géométrique} : la normale à l'interface (le rayon terrestre) tourne d'un petit angle $\mathrm{d}\theta$ d'une coquille à la suivante. Un calcul de géométrie (projection du trajet sur les deux interfaces sphériques voisines) montre que l'invariant n'est pas $\sin i / v$ mais $r \sin i / v$ : le facteur $r$ compense exactement la rotation de la normale.
\item \textbf{Analogie optique} : en optique plane stratifiée, $n \sin i = \text{constante}$ (avec $n = c/v$ l'indice). En géométrie sphérique, l'invariant devient $r \sin i / v = \text{constante}$. C'est le même principe de Fermat, « habillé » pour une sphère.
\item \textbf{Conclusion} : connaître $p$ (mesurable à la surface) suffit à décrire tout le trajet du rayon dans la Terre.
\end{enumerate}
\end{experience}

\begin{definition}[Grandeurs fondamentales de la sismologie sphérique]
\begin{itemize}
\item \textbf{Distance épicentrale} $\Delta$ : angle au centre de la Terre entre l'épicentre et la station d'enregistrement, exprimé en degrés ($1^{\circ} \approx \SI{111,19}{km}$ à la surface).
\item \textbf{Temps de parcours} $t(\Delta)$ : durée mise par une phase sismique donnée pour aller du foyer à la station ; c'est la donnée brute du sismologue.
\item \textbf{Paramètre de rayon} $p$ : pente de la courbe $t(\Delta)$, $p = \dfrac{\mathrm{d}t}{\mathrm{d}\Delta}$, exprimé en $\si{\second\per\radian}$ ou en $\si{\second\per\degree}$. Il caractérise chaque rayon : plus $p$ est grand, plus le rayon est « rasant » (peu profond) ; plus $p$ est petit, plus le rayon plonge profondément.
\end{itemize}
\end{definition}

\courssub{Réfraction continue : la courbure des rayons dans le manteau}

Dans le manteau, la vitesse des ondes \emph{augmente régulièrement avec la profondeur} (la pression rigidifie les roches). Que devient alors le rayon ?

Raisonnons avec le paramètre de rayon : $p = r\sin i / v = \text{constante}$. Quand le rayon descend, $r$ diminue mais $v$ augmente. Pour que $p$ reste constant, $\sin i$ doit augmenter : l'angle $i$ croît, le rayon \emph{s'écarte de la verticale} et s'incurve. Il existe une profondeur de \cle{point bas} où $i = 90^{\circ}$ (le rayon est horizontal) : là, $\sin i = 1$, donc
\[ p = \frac{r_{\text{bas}}}{v(r_{\text{bas}})} \]
Puis le rayon remonte symétriquement vers la surface. C'est pourquoi les sismogrammes enregistrent des ondes ayant « plongé » à des centaines de kilomètres puis « rebondi » vers le haut \emph{sans jamais rencontrer de miroir} : c'est une réfraction continue, un effet de la \emph{rampe} de vitesse, exactement comme un rayon lumineux se courbe dans un mirage.

\begin{aretenir}[Règle de courbure]
\begin{itemize}
\item \textbf{Gradient de vitesse positif} ($v$ augmente avec la profondeur, cas normal dans le manteau) : les rayons se \emph{courbent vers le haut} et reviennent à la surface. Ils « plongent puis remontent ».
\item \textbf{Gradient de vitesse négatif} ($v$ diminue avec la profondeur, ex. : entrée dans le noyau liquide où $v_P$ chute brutalement de $\approx \SI{13,7}{km/s}$ à $\approx \SI{8}{km/s}$) : les rayons se courbent \emph{vers le bas}, sont déviés brutalement, ce qui crée des \emph{zones d'ombre}.
\end{itemize}
Retiens l'image : le rayon « fuit » les zones lentes et « s'attarde » dans les zones rapides (principe de Fermat).
\end{aretenir}

\begin{popfigure}
\begin{center}
\begin{tikzpicture}[scale=1.1]
% Couches (du centre vers l'extérieur)
\fill[popgoldL] (0,0) circle (0.9);    % Noyau interne
\fill[poporangeL] (0,0) circle (1.9);  % Noyau externe
\fill[popblueL] (0,0) circle (3.4);    % Manteau
\fill[popgreenL] (0,0) circle (3.55);  % Croûte
% Contours
\draw[popdark, thick] (0,0) circle (3.55); % Surface
\draw[popdark] (0,0) circle (3.4);   % Moho
\draw[popdark] (0,0) circle (1.9);   % Gutenberg
\draw[popdark] (0,0) circle (0.9);   % Lehmann
% Épicentre (Pôle Nord)
\fill[popink] (0,3.55) circle (0.07);
\node[above, popink, font=\small] at (0,3.62) {Épicentre};
% Station P (Delta ~ 30 deg)
\coordinate (StaP) at ({3.55*sin(30)}, {3.55*cos(30)}); % (1.775, 3.074)
\fill[popblue] (StaP) circle (0.05);
% Station PcP (Delta ~ 60 deg)
\coordinate (StaPcP) at ({3.55*sin(60)}, {3.55*cos(60)}); % (3.074, 1.775)
\fill[poporange] (StaPcP) circle (0.05);
% Station PKP (Delta ~ 150 deg)
\coordinate (StaPKP) at ({3.55*sin(150)}, {3.55*cos(150)}); % (1.775, -3.074)
\fill[poppurple] (StaPKP) circle (0.05);

% Rayon P (Bleu) : plonge jusqu'à r~3.0 (manteau sup), revient
% Point bas : angle 15 deg, r=3.0 -> (0.776, 2.898)
\draw[popblue, very thick] (0,3.55) .. controls (0.9, 3.1) and (0.9, 2.6) .. (0.776, 2.898)
    .. controls (0.9, 2.6) and (1.3, 3.1) .. (StaP);
\draw[popblue, very thick, -{Stealth}] (0,3.55) .. controls (0.5, 3.2) and (0.5, 3.0) .. (0.4, 2.9);
\node[popblue, right, font=\small] at (StaP) {P (réfractée)};

% Rayon PcP (Orange) : réflexion sur CMB (r=1.9)
% Incidence sur CMB à angle ~30 deg -> point (1.9*sin30, 1.9*cos30) = (0.95, 1.645)
% Réflexion symétrique vers station 60 deg.
\coordinate (HitPcP) at (0.95, 1.645);
\draw[poporange, very thick] (0,3.55) -- (HitPcP);
\draw[poporange, very thick] (HitPcP) -- (StaPcP);
\draw[poporange, very thick, -{Stealth}] (0,3.55) -- (0.4, 3.0);
\node[poporange, right, font=\small] at (StaPcP) {PcP (réfléchie)};

% Rayon PKP (Violet) : traverse noyau externe
% Entrée CMB côté gauche (angle ~150 deg? non, epicentre nord. Rayon part vers le bas à gauche).
% Point entrée CMB : angle 120 deg (depuis verticale) -> r=1.9. (-1.9*sin60, 1.9*cos60) = (-1.645, 0.95)
% Sortie CMB vers station 150 deg (depuis verticale) -> angle 150 deg. Point (1.9*sin150, 1.9*cos150) = (0.95, -1.645)
% Trajet dans noyau : ligne droite (vitesse constante approx).
\coordinate (InPKP) at (-1.645, 0.95);
\coordinate (OutPKP) at (0.95, -1.645);
\draw[poppurple, very thick] (0,3.55) -- (InPKP);
\draw[poppurple, very thick] (InPKP) -- (OutPKP);
\draw[poppurple, very thick] (OutPKP) -- (StaPKP);
\draw[poppurple, very thick, -{Stealth}] (0,3.55) -- (-0.5, 3.0);
\node[poppurple, left, font=\small] at (StaPKP) {PKP (noyau)};

% Légendes couches
\node[popdark, font=\small\bfseries] at (0,0) {Noyau interne};
\node[popdark, font=\small] at (0,-1.45) {Noyau externe (liquide)};
\node[popdark, font=\small, rotate=30] at (-2.2, 2.2) {Manteau};
\node[popdark, font=\small] at (2.5, 2.4) {Croûte};
\end{tikzpicture}
\end{center}
\legende{Coupe schématique du globe terrestre (échelle respectée : rayons 0,9 / 1,9 / 3,4 / 3,55) et trajets des principaux rayons sismiques. Le rayon P (bleu) se courbe progressivement dans le manteau car la vitesse augmente avec la profondeur (réfraction continue, point bas vers $r\approx 3,0$). Le rayon PcP (orange) subit une réflexion spéculaire sur la discontinuité de Gutenberg ($r=1,9$). Le rayon PKP (violet) est réfracté à l'entrée du noyau externe (chute brutale de $v_P$), traverse le noyau liquide en ligne droite, puis ressort dévié : cette déviation crée la zone d'ombre des ondes P entre $103^{\circ}$ et $143^{\circ}$.}
\end{popfigure}

\courssub{Réflexion spéculaire et conversion de mode aux discontinuités}

\textbf{Réflexion.} Lorsqu'un rayon rencontre une discontinuité franche (Moho, surface du noyau), une partie de l'énergie est \emph{réfléchie} comme la lumière sur un miroir : l'angle de réflexion est égal à l'angle d'incidence (mesurés par rapport à la normale). C'est ainsi que naît la phase \cle{PcP} : onde P descendue dans le manteau, réfléchie sur le noyau (« c » pour \emph{core}), remontée en onde P. De même, \cle{ScS} est une onde S réfléchie sur le noyau.

\textbf{Conversion de mode.} Mieux encore : à une interface, une onde P incidente peut donner naissance à une onde S réfléchie ou réfractée (et réciproquement). On parle de \cle{conversion de mode}. Par exemple, une onde P arrivant sur le Moho peut repartir vers le haut en onde S. Ces conversions sont précieuses : en comparant les temps d'arrivée des phases converties et directes, on mesure l'épaisseur des couches — c'est la base de la sismologie de la croûte continentale, utilisée aussi pour explorer le socle du Cameroun.

\begin{exemplebox}[Vocabulaire des phases (nomenclature IASPEI)]
\begin{itemize}
\item \textbf{P} : onde P directe, réfractée dans le manteau.
\item \textbf{PcP} : P réfléchie une fois sur le noyau externe.
\item \textbf{PKP} : P ayant traversé le noyau externe (« K » pour \emph{Kern}, noyau en allemand).
\item \textbf{PKIKP} : P ayant traversé le noyau externe \emph{et} le noyau interne (« I » pour \emph{Inner core}).
\item \textbf{S, ScS} : équivalents pour les ondes S (qui ne pénètrent jamais le noyau liquide).
\end{itemize}
\end{exemplebox}

\courssub{La zone d'ombre : quand le silence parle}

\begin{definition}[Zone d'ombre sismique]
Une \cle{zone d'ombre} est une bande de distances épicentrales où aucune onde d'un type donné n'est enregistrée, parce que les rayons correspondants ont été déviés (réfraction brutale vers le bas) ou absorbés/arrêtés (ondes S dans un liquide).
\end{definition}

Pour les ondes P : les rayons arrivant à la limite manteau-noyau sont brutalement déviés vers le bas (chute de vitesse) ; ils ressortent plus loin qu'ils ne le devraient. Résultat : entre $\Delta \approx 103^{\circ}$ et $\Delta \approx 143^{\circ}$, aucune onde P directe n'arrive. Pour les ondes S, la zone d'ombre est totale au-delà de $103^{\circ}$ : le noyau externe liquide les bloque complètement. Ce « silence » des sismographes est une preuve éclatante : il a permis de mesurer le rayon du noyau et de démontrer son état liquide, sans jamais y forer.

\begin{savaistu}[Beno Gutenberg et la zone d'ombre (1913)]
C'est en analysant la zone d'ombre des ondes P entre $103^{\circ}$ et $143^{\circ}$ que le sismologue allemand Beno Gutenberg détermina en 1913 la profondeur de la limite manteau-noyau : environ \SI{2900}{km}, valeur toujours admise aujourd'hui. Cette interface porte le nom de \emph{discontinuité de Gutenberg}. Quelques décennies plus tard, Inge Lehmann découvrira de la même façon le noyau interne solide (1936), en repérant des ondes « interdites » (PKIKP) au cœur même de la zone d'ombre.
\end{savaistu}

\courssub{Exploitation quantitative : le paramètre de rayon en pratique}

\begin{popfigure}
\begin{center}
\begin{tikzpicture}
\begin{axis}[
  width=0.9\linewidth, height=7cm,
  xlabel={Distance épicentrale $\Delta$ (degrés)},
  ylabel={Paramètre de rayon $p$ (\si{\second\per\degree})},
  xmin=0, xmax=180, ymin=0, ymax=14,
  grid=both, grid style={popline!40},
  legend style={at={(0.02,0.98)}, anchor=north west, font=\small, fill=popcream, draw=popline},
  tick label style={font=\small}, label style={font=\small}
]
% P mantle: decreasing from ~12 at 20deg to ~2 at 100deg
\addplot[popblue, very thick, domain=20:100, samples=80] {13.5 - 0.115*x};
\addlegendentry{P (manteau)}
% PKP core: low p, slightly decreasing
\addplot[poppurple, very thick, domain=145:180, samples=40] {4.6 - 0.012*(x-145)};
\addlegendentry{PKP (noyau)}
% PcP reflection: low p, increasing slightly
\addplot[poporange, very thick, domain=30:90, samples=40] {2.2 + 0.018*x};
\addlegendentry{PcP (réflexion noyau)}
% S mantle: high p, decreasing
\addplot[popgreen, very thick, domain=20:75, samples=50] {24.5 - 0.21*x};
\addlegendentry{S (manteau)}
% Shadow zone markers
\draw[popink, thick, dashed] (103,0) -- (103,14) node[above, pos=0.5, sloped, font=\tiny, popink] {Début ombre P};
\draw[popink, thick, dashed] (143,0) -- (143,14) node[above, pos=0.5, sloped, font=\tiny, popink] {Fin ombre P};
\end{axis}
\end{tikzpicture}
\end{center}
\legende{Allure du paramètre de rayon $p$ en fonction de la distance épicentrale $\Delta$ pour les phases principales (valeurs simplifiées à but pédagogique). Les phases directes (P, S) ont un $p$ élevé qui décroît avec $\Delta$ : plus la station est lointaine, plus le rayon a plongé profondément. Les phases réfléchies (PcP) ont un $p$ faible et quasi constant. La coupure de la branche P vers $100^{\circ}$ matérialise l'entrée dans la zone d'ombre ; la branche PKP, à $p$ faible, correspond aux rayons ayant traversé le noyau. Les traits roses pointillés marquent les bornes de la zone d'ombre P.}
\end{popfigure}

\begin{methode}[Tracer un rayon connaissant $p$ et le modèle de vitesse $v(r)$]
\begin{enumerate}
\item \textbf{Mesurer} $p$ sur les données : c'est la pente locale de la courbe des temps de parcours, $p = \mathrm{d}t/\mathrm{d}\Delta$.
\item \textbf{Trouver le point bas} : résoudre $v(r_{\text{bas}}) = r_{\text{bas}}/p$ dans le modèle de vitesse. Cela donne la profondeur maximale atteinte par le rayon.
\item \textbf{Calculer l'angle local} en tout point du trajet : $\sin i(r) = \dfrac{p\,v(r)}{r}$.
\item \textbf{Intégrer pas à pas} (numériquement ou graphiquement) : avancer par petits segments en suivant la direction donnée par $i(r)$, de la surface au point bas, puis symétriquement vers le haut.
\item \textbf{Vérifier} : la distance épicentrale totale et le temps total calculés doivent correspondre aux observations ; sinon, corriger le modèle $v(r)$. C'est ainsi qu'ont été construits les modèles de référence (PREM, AK135).
\end{enumerate}
\end{methode}

\begin{exoresolu}[Angle d'incidence d'un rayon P au niveau du Moho]
\textbf{Énoncé.} Un rayon P arrive à une station située à $\Delta = 30^{\circ}$ de l'épicentre, avec un paramètre de rayon $p = \SI{4}{s/\degree}$. Calculer l'angle d'incidence $i$ de ce rayon lorsqu'il traverse le Moho, situé à $r = \SI{6341}{km}$ du centre, où la vitesse des ondes P vaut $v = \SI{8}{km/s}$.
\tcblower
\textbf{Solution détaillée.}
\begin{enumerate}
\item \textbf{Convertir $p$ en unités cohérentes.} Le paramètre est donné en $\si{\second\per\degree}$ ; or la formule $p = r\sin i / v$ exige des radians :
\[ p = 4 \ \si{\second\per\degree} \times \frac{180^{\circ}}{\pi \ \si{\radian}} = \frac{720}{\pi} \approx 229{,}2 \ \si{\second\per\radian}. \]
\item \textbf{Appliquer la conservation du paramètre de rayon} au niveau du Moho :
\[ \sin i = \frac{p \cdot v}{r} = \frac{229{,}2 \times 8}{6341} = \frac{1833{,}6}{6341} \approx 0{,}289. \]
\item \textbf{Conclure} :
\[ i = \arcsin(0{,}289) \approx 16{,}8^{\circ}. \]
Le rayon traverse le Moho en faisant un angle d'environ $17^{\circ}$ avec la verticale locale : il est encore assez « redressé », ce qui est logique pour une distance épicentrale modérée de $30^{\circ}$.
\end{enumerate}
\textbf{Vérification de cohérence} : $\sin i < 1$, le résultat est physiquement possible. Si l'on avait trouvé $\sin i > 1$, cela signifierait que le rayon ne peut pas exister à cette profondeur (il a déjà atteint son point bas plus haut).
\end{exoresolu}

\begin{attention}[Pièges fréquents au Probatoire]
\begin{itemize}
\item \textbf{N'applique jamais la loi plane simple} $\sin i_1 / v_1 = \sin i_2 / v_2$ sur de grandes distances à l'échelle du globe : la courbure de la Terre impose l'invariant sphérique $p = r\sin i / v$. La loi plane n'est valable que \emph{localement}, à une interface donnée.
\item \textbf{Attention aux unités de $p$} : un paramètre en $\si{\second\per\degree}$ doit être converti en $\si{\second\per\radian}$ (multiplier par $180/\pi$) avant tout calcul avec $r$ en km et $v$ en km/s.
\item \textbf{L'angle $i$ se mesure par rapport à la normale} (la verticale locale, direction radiale), jamais par rapport à l'interface elle-même.
\item \textbf{Ne dis pas} que les rayons « rebondissent » dans le manteau : ils s'y \emph{courbent continûment} par réfraction ; le rebond (réflexion) n'a lieu que sur une discontinuité franche.
\end{itemize}
\end{attention}

\begin{aretenir}[L'essentiel de la section]
\begin{itemize}
\item Les ondes sismiques suivent le trajet de \textbf{temps minimal} (principe de Fermat) : ce sont des rayons, généralement courbes.
\item Dans une Terre sphérique, le \textbf{paramètre de rayon} $p = r\sin i / v(r)$ est conservé le long de chaque rayon ; il est mesurable comme pente de la courbe $t(\Delta)$.
\item Vitesse croissante avec la profondeur $\Rightarrow$ rayons courbés vers le haut (ils reviennent à la surface) ; chute de vitesse $\Rightarrow$ déviation brutale et \textbf{zone d'ombre}.
\item Aux discontinuités : \textbf{réflexion spéculaire} ($i = i'$, phases PcP, ScS) et \textbf{conversion de mode} P $\leftrightarrow$ S.
\item La zone d'ombre des P ($103^{\circ}$--$143^{\circ}$) a révélé l'existence et la taille du noyau ; celle des S a prouvé qu'il est liquide à l'extérieur.
\end{itemize}
\end{aretenir}

Désormais équipés des lois du voyage, nous pouvons passer à la pratique : comment, concrètement, les sismologues exploitent-ils les courbes de temps de parcours pour localiser précisément les discontinuités ? C'est l'objet de la section suivante.

\coursec{Exploitation des courbes de temps de parcours : localiser les discontinuités}

Dans la section précédente, nous avons suivi le voyage des ondes P et S à travers le globe : réfractées selon la loi de Snell-Descartes, réfléchies aux discontinuités, elles émergent à la surface après avoir traversé des milieux de plus en plus profonds. Mais une question fondamentale demeure : comment, à partir de ces arrivées d'ondes enregistrées en surface, remonter à la \emph{structure interne} de la Terre ? C'est tout l'art du sismologue : chaque sismogramme est un message codé, et la \cle{courbe de temps de parcours} est la clé qui permet de le déchiffrer. C'est cette clé que nous allons maintenant fabriquer ensemble.

\courssub{Situation-problème : lire le message des profondeurs}

Imaginons la station sismologique du lycée de Yaoundé, reliée au réseau national de surveillance. Un séisme se produit au Chili. Sur l'écran, le sismogramme n'affiche pas une seule secousse, mais une \emph{succession d'arrivées} étalées sur plus de vingt minutes : une première arrivée discrète (les P), puis une arrivée plus énergique (les S), puis d'autres signaux plus tardifs. D'où viennent ces échos tardifs ? Ont-ils rebondi quelque part ? Ont-ils traversé le noyau ?

\begin{center}
\textbf{Démarche d'investigation}
\end{center}

\begin{itemize}
\item \textbf{Observation :} pour un même séisme, plusieurs trains d'ondes arrivent à des instants différents, et leurs temps d'arrivée varient régulièrement avec la distance station--épicentre.
\item \textbf{Hypothèse :} chaque arrivée correspond à un \emph{chemin différent} emprunté dans la Terre (direct, réfléchi sur une discontinuité, traversée du noyau).
\item \textbf{Données :} on reporte, pour des centaines de séismes et de stations, le temps de parcours $t$ en fonction de la distance épicentrale $\Delta$ : les points s'alignent sur des \emph{branches} régulières.
\item \textbf{Interprétation :} chaque branche correspond à une famille de trajets ; les \emph{cassures} entre branches signalent les discontinuités internes.
\item \textbf{Conclusion :} la courbe $t(\Delta)$ est une « radiographie » indirecte du globe : sa lecture permet de localiser les discontinuités et de construire le modèle en couches.
\end{itemize}

\begin{definition}[Distance épicentrale et temps de parcours]
La \cle{distance épicentrale} $\Delta$ est l'angle, mesuré au centre de la Terre, entre l'épicentre et la station (elle s'exprime en degrés, de $0^\circ$ à $180^\circ$). Le \cle{temps de parcours} $t$ est la durée mise par une phase sismique donnée pour aller du foyer à la station. La \cle{courbe de temps de parcours} (ou hodochrone) est la représentation graphique $t = f(\Delta)$ pour chaque phase.
\end{definition}

\begin{attention}[Degrés, kilomètres, secondes : ne pas tout mélanger !]
Sur une courbe $t(\Delta)$, l'axe vertical est en \emph{secondes} (ou minutes) et l'axe horizontal en \emph{degrés}. Ce ne sont pas des unités interchangeables ! La conversion utile à retenir est :
\[ 1^\circ \text{ à la surface de la Terre} \approx 111\ \text{km} \qquad (\text{circonférence} \approx \num{40000}\ \text{km} / 360^\circ). \]
Ainsi $\Delta = 60^\circ$ correspond à environ $\SI{6700}{km}$ d'arc de surface. Confondre $t$ et $\Delta$, ou oublier la conversion degré--kilomètre, est l'erreur la plus fréquente au Probatoire.
\end{attention}

\courssub{Les branches de la courbe $t(\Delta)$ : directes, réfléchies, réfractées}

Chaque phase sismique porte un nom qui décrit son trajet : une lettre majuscule par segment de parcours (P ou S dans le manteau, K dans le noyau externe, I dans le noyau interne), la lettre « c » pour une réflexion sur la discontinuité de Gutenberg (CMB), « i » pour une réflexion sur la discontinuité de Lehmann (ICB).

\begin{definition}[Principales phases sismiques]
\begin{itemize}
\item \textbf{Phases directes :} \cle{P} et \cle{S}, qui traversent croûte et manteau sans rebond.
\item \textbf{Phases réfléchies en surface :} \cle{PP} et \cle{SS} (un rebond à la surface à mi-chemin).
\item \textbf{Phases réfléchies sur le noyau :} \cle{PcP} (P réfléchie sur la CMB) et \cle{ScS}.
\item \textbf{Phases traversant le noyau :} \cle{PKP} (P -- noyau externe -- P), \cle{PKIKP} (qui traverse en plus le noyau interne).
\end{itemize}
\end{definition}

\begin{popfigure}
\begin{center}
\begin{tikzpicture}
\begin{axis}[
  width=0.95\linewidth, height=8cm,
  xlabel={Distance épicentrale $\Delta$ (degrés)},
  ylabel={Temps de parcours $t$ (s)},
  xmin=0, xmax=180, ymin=0, ymax=1500,
  xtick={0,30,60,90,120,150,180},
  ytick={0,300,600,900,1200,1500},
  grid=both, grid style={popline!40},
  legend style={at={(0.02,0.98)},anchor=north west,font=\small},
]
% P directe : concave vers le bas (pente decroissante)
\addplot[popcurve,domain=0:100,samples=120]{60 + 11.5*x - 0.028*x*x};
\addlegendentry{P (directe)}
% S directe
\addplot[popcurve,poporange,domain=0:100,samples=120]{110 + 21*x - 0.05*x*x};
\addlegendentry{S (directe)}
% PcP : branche courte, faibles distances
\addplot[popcurve,popgreen,domain=0:70,samples=80]{510 - 2.2*x + 0.012*x*x};
\addlegendentry{PcP (réfléchie sur CMB)}
% ScS
\addplot[popcurve,poppurple,domain=0:70,samples=80]{920 - 4*x + 0.022*x*x};
\addlegendentry{ScS}
% PKP : apparait au-dela de ~110 degres
\addplot[popcurve,poppink,domain=110:180,samples=80]{1180 + 4.4*(x-110)};
\addlegendentry{PKP}
% PKIKP
\addplot[popcurve,popgold,domain=115:180,samples=80]{1200 + 3.1*(x-115)};
\addlegendentry{PKIKP}
% PP
\addplot[popcurve,popmuted,domain=40:180,samples=80]{2*(60 + 11.5*(x/2) - 0.028*(x/2)*(x/2))};
\addlegendentry{PP}
\end{axis}
\end{tikzpicture}
\end{center}
\legende{Courbes de temps de parcours schématiques $t(\Delta)$ des principales phases sismiques. Chaque branche correspond à un type de trajet dans le globe. Remarque : la branche P est concave (sa pente diminue quand $\Delta$ augmente), car les ondes profondes, plus rapides, « gagnent » aux grandes distances. Entre $103^\circ$ et $140^\circ$ environ, l'absence des branches P et S directes matérialise l'ombre du noyau.}
\end{popfigure}

\begin{propriete}[La pente de la courbe $t(\Delta)$ : le paramètre de rayon]
En tout point d'une branche, la pente de la courbe de temps de parcours est égale au \cle{paramètre de rayon} $p$ du rayon qui émerge à cette distance :
\[ p = \frac{\mathrm{d}t}{\mathrm{d}\Delta} \quad \text{(en s/deg, ou s/km si $\Delta$ est en km).} \]
Or $p = r\,\sin i / v(r)$ est \emph{constant le long d'un rayon donné} (conséquence de la loi de Snell-Descartes en milieu sphérique). Mesurer la pente de la hodochrone revient donc à connaître le point le plus profond atteint par le rayon : plus la pente est faible, plus le rayon a plongé profondément. C'est le fondement mathématique de toute la sismologie globale.
\end{propriete}

\courssub{Les cassures de pente : signature des discontinuités}

\begin{experience}[Observer une cassure sur des données réelles]
\textbf{Matériel :} table de temps de parcours IASP91 (extraits), papier millimétré ou tableur.\\
\textbf{Protocole :} on relève les temps de la phase P tous les $5^\circ$ entre $10^\circ$ et $40^\circ$, puis on calcule la pente locale $\Delta t / \Delta \Delta$ entre points successifs.\\
\textbf{Observations :} la pente décroît régulièrement, mais présente des \emph{sauts} brutaux : par exemple vers $20^\circ$, l'arrivée P se dédouble (deux arrivées proches), et au-delà de $30^\circ$ la pente change nettement.\\
\textbf{Interprétation :} les ondes dont le point bas se situe juste au-dessus d'une discontinuité (410 km, 660 km) et celles réfractées juste en dessous arrivent presque en même temps mais par des chemins différents : la branche se « replie » (triplication). Chaque cassure ou triplication de la courbe $t(\Delta)$ signale donc une discontinuité ou un changement brutal de gradient de vitesse.\\
\textbf{Conclusion :} repérer les cassures de $\mathrm{d}t/\mathrm{d}\Delta$ permet de \emph{localiser} les discontinuités sans jamais forer le sol.
\end{experience}

\begin{aretenir}[Ce qu'il faut retenir des cassures]
Une discontinuité interne se manifeste sur la courbe $t(\Delta)$ par :
\begin{itemize}
\item une \textbf{cassure de pente} (changement brutal de $\mathrm{d}t/\mathrm{d}\Delta$) ;
\item parfois une \textbf{triplication} (la branche se replie : trois arrivées possibles à une même distance, caractéristique des transitions à 410 et 660 km) ;
\item ou une \textbf{zone d'ombre} (disparition totale d'une branche, comme les P directes au-delà de $103^\circ$ à cause du noyau).
\end{itemize}
\end{aretenir}

\courssub{La méthode du temps réduit : aplatir pour mieux voir}

Les branches de $t(\Delta)$ sont fortement inclinées : sur $180^\circ$, le temps varie de plus de 20 minutes. Les petites cassures, pourtant capitales, sont noyées dans cette pente globale. L'astuce des sismologues consiste à \emph{soustraire} une tendance linéaire.

\begin{definition}[Temps réduit]
Le \cle{temps réduit} est défini par :
\[ t_{\mathrm{red}} = t - \frac{\Delta}{v_{\mathrm{reduc}}} \]
où $v_{\mathrm{reduc}}$ est une vitesse de réduction arbitraire (souvent $8\ \text{km/s}$ en sismologie crustale, ou $11{,}2\ \text{km/s}$ en sismologie globale, $\Delta$ étant alors converti en km). Les phases proches de la vitesse de réduction deviennent quasi horizontales : les écarts, sauts et triplications apparaissent alors en pleine lumière.
\end{definition}

\begin{popfigure}
\begin{center}
\begin{tikzpicture}
\begin{axis}[
  width=0.95\linewidth, height=7.5cm,
  xlabel={Distance épicentrale $\Delta$ (degrés)},
  ylabel={Temps réduit $t - \Delta \times 111 / 8$ (s)},
  xmin=0, xmax=120, ymin=-700, ymax=100,
  xtick={0,20,40,60,80,100,120},
  grid=both, grid style={popline!40},
]
% branche P reduite avec triplications 410 et 660
\addplot[popcurve,domain=0:14,samples=60]{60 + 11.5*x - 0.028*x*x - 13.875*x};
\addplot[popcurve,domain=14:20,samples=40]{60 + 11.5*x - 0.028*x*x - 13.875*x + 18};
\addplot[popcurve,domain=20:28,samples=40]{60 + 11.5*x - 0.028*x*x - 13.875*x + 30};
\addplot[popcurve,domain=28:100,samples=80]{60 + 11.5*x - 0.028*x*x - 13.875*x + 42};
% annotations
\node[font=\small, popgreen, anchor=west] at (axis cs:15,-60) {saut : 410 km};
\node[font=\small, poporange, anchor=west] at (axis cs:29,-130) {saut : 660 km};
\node[font=\small, popdark, anchor=south west] at (axis cs:2,-620) {branche P réduite : les cassures deviennent visibles};
\end{axis}
\end{tikzpicture}
\end{center}
\legende{Courbe de temps réduit $t_{\mathrm{red}} = t - \Delta \times 111 / 8$ (vitesse de réduction $8\ \text{km/s}$). La branche P, presque plate, présente des sauts bien visibles correspondant aux triplications liées aux discontinuités de 410 km et 660 km (zone de transition du manteau).}
\end{popfigure}

\begin{exemplebox}[Calcul d'un temps réduit : phase P directe]
Une phase P directe est observée à $\Delta = 30^\circ$ avec un temps de parcours $t = 375\ \text{s}$ (valeur typique IASP91). Avec $v_{\mathrm{reduc}} = 8\ \text{km/s}$ :
\[ t_{\mathrm{red}} = 375 - \frac{30 \times 111}{8} = 375 - 416 = -41\ \text{s}. \]
Le temps réduit est \emph{négatif} : cela signifie que la phase P « voyage » en moyenne plus vite que la vitesse de réduction (elle a plongé dans le manteau profond, rapide). Pour comparaison, une phase PcP (réfléchie sur le noyau) à la même distance arrive vers $t \approx 550\ \text{s}$, donnant $t_{\mathrm{red}} \approx +134\ \text{s}$ : sur le graphique en temps réduit, les deux phases sont bien séparées, ce qui facilite leur identification.
\end{exemplebox}

\courssub{Mesurer la profondeur d'une discontinuité : l'hyperbole de réflexion}

Considérons un cas simple et exigible : une discontinuité \emph{plane et horizontale} à la profondeur $h$, sous un milieu de vitesse constante $v$ (approximation locale valable pour le Moho sous une station). Une onde émise au foyer, réfléchie sur la discontinuité, revient à la station distante de $\Delta$ (en km).

\begin{propriete}[Hyperbole de réflexion]
Le temps de parcours de l'onde réfléchie vérifie :
\[ t^2 = t_0^2 + \frac{\Delta^2}{v^2} \qquad \text{avec} \qquad t_0 = \frac{2h}{v}, \]
où $t_0$ est le temps de parcours vertical aller-retour (à $\Delta = 0$). La courbe $t(\Delta)$ est donc une \emph{hyperbole}. En portant $t^2$ en fonction de $\Delta^2$, on obtient une droite dont l'ordonnée à l'origine donne $t_0^2$, d'où la profondeur : $h = v\,t_0 / 2$.
\end{propriete}

\begin{exoresolu}[Profondeur du Moho sous le Mont Cameroun]
\textbf{Énoncé.} Sous la région du Mont Cameroun, une campagne de sismique réflexion enregistre une réflexion sur le Moho. À distance nulle, le temps vertical aller-retour est $t_0 = 10\ \text{s}$. La vitesse moyenne des ondes P dans la croûte continentale est $v = 6{,}0\ \text{km/s}$. (1) Calculer la profondeur $h$ du Moho. (2) Quel temps de parcours mesure-t-on pour une station à $\Delta = 30\ \text{km}$ du tir ?
\tcblower
\textbf{Solution.} (1) De $t_0 = 2h/v$, on tire :
\[ h = \frac{v\,t_0}{2} = \frac{6{,}0 \times 10}{2} = 30\ \text{km}. \]
Le Moho se situe donc à 30 km de profondeur, valeur typique d'une croûte continentale. (2) D'après l'hyperbole de réflexion :
\[ t^2 = t_0^2 + \frac{\Delta^2}{v^2} = 100 + \frac{900}{36} = 125 \quad\Rightarrow\quad t = \sqrt{125} \approx 11{,}2\ \text{s}. \]
L'écart avec $t_0$ (1,2 s) est faible mais mesurable : c'est précisément cet écart qui permet de déterminer $v$.
\end{exoresolu}

\courssub{L'inversion globale : de Herglotz-Wiechert aux tables IASP91 et AK135}

Pour le globe entier, la vitesse varie continûment avec le rayon et l'approximation plane ne suffit plus. La méthode de \cle{Herglotz-Wiechert} (1907, admise à notre niveau) résout le problème inverse : à partir de la courbe complète $t(\Delta)$ et de sa pente $p(\Delta)$, une transformation intégrale permet de calculer, pour chaque rayon, la vitesse $v(r)$. Là où $v(r)$ présente un \emph{saut}, il y a discontinuité ; là où $v(r)$ \emph{décroît}, il y a zone d'ombre. C'est ainsi qu'ont été établies les profondeurs de référence :

\begin{center}
\begin{tabularx}{\linewidth}{|l|X|X|}\hline
\rowcolor{popblueL}
\textbf{Discontinuité} & \textbf{Profondeur} & \textbf{Nature} \\ \hline
Moho (Mohorovičić) & $\approx 35$ km (5--70 km) & croûte / manteau \\ \hline
« 410 km » & 410 km & olivine $\rightarrow$ wadsleyite \\ \hline
« 660 km » & 660 km & ringwoodite $\rightarrow$ bridgmanite + ferropericlase \\ \hline
Gutenberg (CMB) & $\approx \num{2890}$ km & manteau solide / noyau externe liquide \\ \hline
Lehmann (ICB) & $\approx \num{5150}$ km & noyau externe liquide / noyau interne solide \\ \hline
\end{tabularx}
\end{center}

\begin{definition}[Les discontinuités majeures]
La \cle{discontinuité du Moho} sépare la croûte du manteau. Les discontinuités de \cle{410 km} et \cle{660 km}, internes au manteau, sont des \emph{transitions de phase} : sous l'effet de la pression, l'olivine se transforme en minéraux plus denses (wadsleyite puis ringwoodite, puis bridgmanite). La \cle{discontinuité de Gutenberg} (CMB) marque la limite manteau--noyau externe liquide ; la \cle{discontinuité de Lehmann} (ICB) marque l'entrée dans le noyau interne solide.
\end{definition}

Enfin, une fois le modèle $v(r)$ connu, on l'utilise « à l'envers » : les \cle{tables de temps de parcours} standard (Jeffreys-Bullen historiquement, puis \cle{IASP91} et \cle{AK135} aujourd'hui) donnent, pour chaque phase et chaque distance $\Delta$, le temps théorique. En comparant les temps \emph{observés} en plusieurs stations aux temps des tables, on localise le foyer (position, profondeur, heure d'origine) : c'est le \emph{problème inverse} de la localisation des séismes, utilisé quotidiennement, y compris pour la surveillance de la zone sismique du Mont Cameroun.

\begin{methode}[Identifier les phases sur un sismogramme et les reporter sur $t(\Delta)$]
\begin{enumerate}
\item Relever sur le sismogramme (station de Douala ou Yaoundé) les temps d'arrivée successifs : première arrivée (P), arrivée énergique (S), arrivées tardives.
\item Calculer la distance épicentrale $\Delta$ (en degrés) à partir des coordonnées de l'épicentre et de la station.
\item Comparer chaque temps mesuré aux tables IASP91/AK135 pour cette distance : la phase théorique la plus proche donne le nom de l'arrivée (P, S, PcP, PP, PKP...).
\item Reporter chaque point $(\Delta, t)$ sur le graphique : il doit se placer sur la branche correspondante.
\item Vérifier la cohérence : un point isolé loin de toute branche signale une erreur de lecture ou une phase mal identifiée.
\end{enumerate}
\end{methode}

\begin{attention}[Pièges fréquents]
\begin{itemize}
\item Ne pas confondre une \emph{cassure de pente} (discontinuité) avec une simple courbure régulière (gradient de vitesse continu).
\item Dans $t_0 = 2h/v$, ne pas oublier le facteur 2 (aller-retour) : oublier ce facteur double la profondeur calculée !
\item Les profondeurs (410, 660, \num{2890}, \num{5150} km) se comptent \emph{depuis la surface}, pas depuis le centre.
\end{itemize}
\end{attention}

\begin{savaistu}[La discontinuité de 660 km : une barrière dans la machine terrestre ?]
La discontinuité de 660 km sépare le manteau supérieur du manteau inférieur. Les sismologues y observent que les plaques plongeantes (subductions) sont parfois \emph{freinées ou aplaties} à cette profondeur, parfois traversées franchement. Faut-il en conclure que la convection mantellique fonctionne sur deux étages superposés, ou bien que le manteau entier convecte d'un seul tenant ? Le débat reste ouvert : c'est l'une des grandes questions de la géodynamique moderne, et ce sont les courbes de temps de parcours -- celles que tu sais maintenant lire -- qui fournissent les arguments.
\end{savaistu}

\begin{exercice}[À toi de jouer]
Une station enregistre une phase PcP à $\Delta = 45^\circ$ avec un temps de parcours $t = 650\ \text{s}$. (1) Convertir $\Delta$ en km d'arc de surface. (2) Calculer le temps réduit avec $v_{\mathrm{reduc}} = 8\ \text{km/s}$. (3) Une réflexion crustale donne $t_0 = 11\ \text{s}$ pour $v = 6\ \text{km/s}$ : calculer la profondeur du Moho sous cette station et commenter.
\end{exercice}

\begin{corrige}[Exercice]
(1) $\Delta = 45 \times 111 = \num{4995}\ \text{km}$. (2) $t_{\mathrm{red}} = 650 - 4995/8 = 650 - 624 = 26\ \text{s}$ : valeur positive, cohérente avec une phase arrivant après la référence (PcP plonge jusqu'à la CMB, chemin long, vitesse apparente inférieure à 8 km/s à cette distance). (3) $h = v\,t_0/2 = 6 \times 11 / 2 = 33\ \text{km}$ : valeur typique d'une croûte continentale, proche de la moyenne du Moho ($\approx 35$ km).
\end{corrige}

En résumé, la courbe $t(\Delta)$ est bien plus qu'un graphique : sa pente donne le paramètre de rayon, ses cassures révèlent les discontinuités, le temps réduit les met en évidence, l'hyperbole de réflexion et la méthode de Herglotz-Wiechert fournissent leurs profondeurs, et les tables IASP91/AK135 bouclent la boucle en permettant de localiser les séismes. Il nous reste à examiner la plus spectaculaire de ces preuves : celle du noyau liquide, objet de la section suivante.

\coursec{Le noyau terrestre : la preuve par les ondes (liquide externe, solide interne)}

Dans la section précédente, tu as appris à exploiter les courbes de temps de parcours (hodochrones) pour repérer les changements de pente qui trahissent les discontinuités internes : la Moho vers 30--70 km de profondeur, puis la discontinuité de Gutenberg à 2 900 km. Mais ces courbes recèlent un mystère encore plus spectaculaire : au-delà d'une certaine distance angulaire, certaines ondes \emph{disparaissent purement et simplement} des sismogrammes, tandis que d'autres réapparaissent plus loin, comme par enchantement. C'est en déchiffrant ces « disparitions » et ces « réapparitions » que les sismologues ont découvert que la Terre possède un noyau liquide enveloppant une graine solide. Cette section te montre comment un raisonnement logique rigoureux, appuyé sur des observations mesurables, permet d'affirmer l'état physique d'une région située à près de 5 000 km sous nos pieds — sans jamais y avoir foré.

\courssub{Les trois observations clés des sismologues}

\textbf{Observation clé 1 : la zone d'ombre des ondes S.}

Imagine un séisme majeur survenant au Chili, enregistré par le réseau mondial de stations sismologiques. Les stations proches de l'épicentre (distance angulaire $\Delta < 103^{\circ}$) enregistrent normalement les ondes P, puis les ondes S. Mais pour toutes les stations situées au-delà de $\Delta \approx 103^{\circ}$ (par exemple en Afrique centrale pour un séisme sud-américain), \emph{aucune onde S directe n'est jamais enregistrée}. Cette région de la surface terrestre, où les ondes S sont absentes, est appelée \cle{zone d'ombre des ondes S}.

Or tu sais (section 1) que les ondes S sont des ondes de cisaillement : elles ne peuvent se propager que dans un milieu possédant un \cle{module de rigidité} $\mu$ non nul, c'est-à-dire un solide. Leur vitesse est :
\[ v_S = \sqrt{\frac{\mu}{\rho}} \]
Si $\mu = 0$ (milieu fluide, sans résistance au cisaillement), alors $v_S = 0$ : l'onde S ne se propage pas. L'absence totale d'ondes S au-delà de $103^{\circ}$ signifie donc que toutes les trajectoires S qui atteindraient ces stations sont bloquées par une région profonde où $\mu = 0$ : un milieu \cle{liquide}.

\textbf{Observation clé 2 : la zone d'ombre des ondes P.}

Les ondes P, elles, traversent bien les liquides. Pourtant, on observe une \cle{zone d'ombre des ondes P} entre $\Delta \approx 103^{\circ}$ et $\Delta \approx 143^{\circ}$ : dans cet intervalle, l'amplitude des ondes P directes s'effondre. L'explication est la \cle{réfraction} (section 2) : à la frontière manteau--noyau (2 900 km), la vitesse des ondes P chute brutalement d'environ $\SI{13,7}{km/s}$ à $\SI{8,1}{km/s}$. D'après la loi de Snell-Descartes, les rayons plongeant dans un milieu plus lent se rapprochent de la normale : ils sont déviés vers le centre de la Terre, ce qui les fait ressortir \emph{plus loin} qu'ils ne le devraient, laissant une bande de surface non desservie.

\textbf{Observation clé 3 : la réapparition des ondes au-delà de 143°.}

Au-delà de $\Delta \approx 143^{\circ}$, des ondes P réapparaissent : ce sont les ondes qui ont traversé le noyau (phases PKP). Mieux : certaines arrivent \emph{plus tôt} que prévu pour un noyau entièrement homogène, ce qui a conduit Inge Lehmann à postuler l'existence d'un \cle{noyau interne solide} où $v_P$ remonte à environ $\SI{11}{km/s}$.

\begin{propriete}[Preuve de la liquidité du noyau externe]
L'argument logique complet, exigible au Probatoire, se construit en quatre étapes :
\begin{enumerate}
\item \textbf{Observation} : aucune onde S n'est enregistrée pour $\Delta > 103^{\circ}$ (zone d'ombre S totale).
\item \textbf{Donnée physique} : une onde S ne se propage que si le milieu possède un module de rigidité $\mu \neq 0$ (solide) ; dans un liquide, $\mu = 0$ et $v_S = 0$.
\item \textbf{Déduction} : il existe une région interne, située sous le manteau ($r < 3\,480$ km), où $\mu = 0$ : le \cle{noyau externe} est \textbf{liquide}.
\item \textbf{Confirmation indépendante} : la chute brutale de $v_P$ (de 13,7 à 8,1 km/s) à 2 900 km et la zone d'ombre P entre $103^{\circ}$ et $143^{\circ}$ confirment une discontinuité majeure avec un milieu de rigidité nulle, qui ralentit fortement les ondes de compression.
\end{enumerate}
La réapparition d'ondes P précoces (PKIKP) au-delà de $143^{\circ}$ impose en outre l'existence d'une \textbf{graine solide} (noyau interne) au centre.
\end{propriete}

\begin{popfigure}
\begin{center}
\begin{tikzpicture}[scale=1.05]
% Terre : Manteau (bleu clair), Noyau externe (orange clair), Noyau interne (or clair)
\fill[popblueL] (0,0) circle (3.2);
\fill[poporangeL] (0,0) circle (1.75);
\fill[popgoldL] (0,0) circle (0.62);
\draw[thick,popink] (0,0) circle (3.2);
\draw[thick,poporange] (0,0) circle (1.75);
\draw[thick,popgold] (0,0) circle (0.62);
% Foyer
\fill[popink] (0,3.2) circle (0.07) node[above=2pt,popink]{\textbf{Foyer}};
% Rayons P dans le manteau (directs)
\draw[thick,popblue] (0,3.2) -- (2.9,1.3);
\draw[thick,popblue] (0,3.2) -- (-2.9,1.3);
% Rayon PKPbc (noyau externe seul, rebroussement dans noyau externe)
\draw[thick,poporange] (0,3.2) -- (1.15,-1.33);
\draw[thick,poporange] (1.15,-1.33) -- (1.62,-0.68);
\draw[thick,poporange] (1.62,-0.68) -- (2.7,-1.75);
% Rayon PKIKP (traverse la graine)
\draw[thick,popgreen] (0,3.2) -- (0.5,-1.68);
\draw[thick,popgreen] (0.5,-1.68) -- (0.15,-0.6);
\draw[thick,popgreen] (0.15,-0.6) -- (-0.3,0.55);
\draw[thick,popgreen] (-0.3,0.55) -- (-1.0,1.45);
\draw[thick,popgreen] (-1.0,1.45) -- (-2.2,2.3);
% Zone d'ombre P (arc 103-143)
\draw[line width=4pt,poppink] (103:3.25) arc (103:143:3.25);
\draw[line width=4pt,poppink] (-103:3.25) arc (-103:-143:3.25);
% Zone d'ombre S (au-delà de 103)
\draw[line width=4pt,poppurple,dashed] (103:3.35) arc (103:257:3.35);
% Légendes pointées
\draw[thin,popmuted] (1.75,0.9) -- (3.6,1.6) node[right,popdark,align=left,font=\small]{Manteau\\(solide)};
\draw[thin,popmuted] (-1.3,-1.1) -- (-3.9,-2.2) node[left,popdark,align=right,font=\small]{Noyau externe\\(liquide, $v_S=0$)};
\draw[thin,popmuted] (0.35,-0.35) -- (3.6,-1.4) node[right,popdark,align=left,font=\small]{Noyau interne\\(solide, graine)};
\draw[thin,popmuted] (2.5,2.1) -- (4.1,2.9) node[right,popdark,font=\small]{Onde P directe};
\draw[thin,popmuted] (2.2,-1.25) -- (4.1,-2.6) node[right,popdark,font=\small]{PKP (noyau externe)};
\draw[thin,popmuted] (-1.7,1.9) -- (-4.3,2.9) node[left,popdark,font=\small]{PKIKP (traverse la graine)};
\node[poppink,font=\small,align=center] at (0,-4.15) {Zone d'ombre P ($103^{\circ}<\Delta<143^{\circ}$)};
\node[poppurple,font=\small,align=center] at (0,4.35) {Zone d'ombre S ($\Delta>103^{\circ}$)};
\end{tikzpicture}
\end{center}
\legende{Coupe du globe montrant les trajectoires des ondes P directes, des ondes PKP (réfractées dans le noyau externe liquide) et PKIKP (traversant la graine solide). Les ondes S, bloquées par le noyau externe liquide, laissent une zone d'ombre totale au-delà de $103^{\circ}$.}
\end{popfigure}

\courssub{Les phases sismiques du noyau : PKP et PKIKP}

Pour nommer le trajet d'une onde, les sismologues utilisent un code de lettres : \textbf{P} = onde P dans le manteau, \textbf{K} = onde P dans le noyau externe (de l'allemand \emph{Kern}), \textbf{I} = onde P dans le noyau interne. Ainsi, une onde qui descend dans le manteau (P), traverse le noyau externe (K) et remonte dans le manteau (P) s'écrit \textbf{PKP}.

\begin{definition}[Branches de l'onde PKP et phase PKIKP]
Selon leur profondeur de pénétration dans le noyau, les ondes P trans-nucléaires se répartissent en branches :
\begin{itemize}
\item \cle{PKPab} (ou PKP\textsubscript{ab}) : branche peu profonde, dont le point de rebroussement se situe dans la partie supérieure du noyau externe ; elle émerge aux grandes distances ($\Delta > 143^{\circ}$).
\item \cle{PKPbc} : branche intermédiaire, plongeant plus profondément dans le noyau externe, frôlant la graine sans la toucher (point de rebroussement à $r > 1\,221$ km).
\item \cle{PKPdf}, appelée aussi \cle{PKIKP} : branche profonde qui \textbf{traverse le noyau interne solide} (segment I). C'est elle qui arrive en première position aux distances $\Delta$ comprises entre environ $143^{\circ}$ et $180^{\circ}$ (station à l'antipode).
\end{itemize}
L'onde S, elle, ne possède aucune branche trans-nucléaire directe : la lettre K n'existe pas pour les ondes de cisaillement, car $v_S = 0$ dans tout le noyau externe.
\end{definition}

\begin{popfigure}
\begin{center}
\begin{tikzpicture}
\begin{axis}[
width=0.92\linewidth, height=7.2cm,
xlabel={Profondeur (km)}, ylabel={Vitesse (km/s)},
xmin=0, xmax=6371, ymin=0, ymax=15,
xtick={0,1000,2000,2890,4000,5150,6371},
xticklabels={0,1000,2000,2890,4000,5150,6371},
ytick={0,2,4,6,8,10,12,14},
grid=both, grid style={popline!40},
legend style={at={(0.03,0.35)},anchor=west,font=\small,draw=popline},
axis lines=left, thick]
% vP : manteau
\addplot[popblue,very thick,domain=0:2890,samples=60] {8 + 5.7*(1-exp(-x/900))};
% vP : noyau externe
\addplot[popblue,very thick,domain=2890:5150,samples=60] {8.1 + 2.3*(x-2890)/2260};
% vP : graine
\addplot[popblue,very thick,domain=5150:6371,samples=30] {11 + 0.3*(x-5150)/1221};
% vS : manteau
\addplot[poporange,very thick,domain=0:2890,samples=60] {4.4 + 2.9*(1-exp(-x/900))};
% vS : noyau externe = 0
\addplot[poporange,very thick,domain=2890:5150,samples=5] {0};
% vS : graine
\addplot[poporange,very thick,domain=5150:6371,samples=30] {3.5 + 0.2*(x-5150)/1221};
% Discontinuités
\draw[dashed,popink] (axis cs:2890,0) -- (axis cs:2890,15);
\draw[dashed,popink] (axis cs:5150,0) -- (axis cs:5150,15);
\node[popink,font=\small,align=center,anchor=south] at (axis cs:2890,14.2) {Gutenberg\\(2 900 km)};
\node[popink,font=\small,align=center,anchor=south] at (axis cs:5150,12.6) {Lehmann\\(5 150 km)};
\node[popdark,font=\small,align=center] at (axis cs:4020,5.5) {Noyau externe\\liquide : $v_S = 0$};
\legend{$v_P$,$v_S$}
\end{axis}
\end{tikzpicture}
\end{center}
\legende{Vitesse des ondes P ($v_P$, bleu) et S ($v_S$, orange) en fonction de la profondeur (modèle PREM simplifié). À 2 900 km (discontinuité de Gutenberg), $v_P$ chute de 13,7 à 8,1 km/s et $v_S$ tombe à zéro : signature d'un liquide. À 5 150 km (discontinuité de Lehmann), $v_P$ et $v_S$ remontent : la graine est solide.}
\end{popfigure}

\courssub{La découverte d'Inge Lehmann (1936) : la graine solide}

Jusqu'au début du XX\textsuperscript{e} siècle, on croyait le noyau entièrement liquide. En 1936, la sismologue danoise \cle{Inge Lehmann} étudie les enregistrements européens de séismes survenus en Nouvelle-Zélande. Elle remarque des arrivées P « impossibles » dans la zone d'ombre : des ondes trop précoces pour avoir été simplement diffractées le long de la surface du noyau. Son hypothèse audacieuse : ces ondes ont été \emph{réfractées par un noyau interne} où la vitesse est plus élevée, c'est-à-dire \textbf{solide}. La discontinuité correspondante, située à \SI{5150}{km} de profondeur (rayon de la graine $\approx \SI{1220}{km}$), porte aujourd'hui le nom de \cle{discontinuité de Lehmann}. L'interprétation physique : malgré des températures de 5 000 à 6 000 °C, la pression colossale au centre de la Terre (environ 360 GPa, soit 3,6 millions d'atmosphères) élève le point de fusion du fer au-dessus de la température réelle : le fer y est \emph{solidifié par la pression}.

\begin{savaistu}[Inge Lehmann, la dame du noyau terrestre]
Inge Lehmann (1888--1993) a découvert le noyau interne solide en 1936, dans un article au titre modeste et célèbre : simplement « P' ». À une époque où les femmes étaient presque exclues des sciences, elle dirigea le service géodésique du Danemark pendant 25 ans. Sa découverte, fondée uniquement sur l'analyse minutieuse de sismogrammes papier, est considérée comme l'une des plus belles déductions de l'histoire de la géophysique. Elle vécut 104 ans et reçut, à 83 ans, la médaille William Bowie, la plus haute distinction de l'Union géophysique américaine.
\end{savaistu}

\courssub{Composition chimique du noyau : un alliage fer--nickel allégé}

Comment connaît-on la composition d'une région inaccessible ? Par recoupement de trois indices :
\begin{itemize}
\item \textbf{La densité} : le moment d'inertie de la Terre et la masse totale de la planète imposent un matériau central très dense, $\rho \approx 10$ à $\SI{13}{g/cm^3}$ (contre 3,3 à 5,7 pour le manteau).
\item \textbf{Les météorites} : les \emph{chondrites}, témoins du matériau primitif du Système solaire, contiennent du fer en abondance ; les \emph{météorites de fer} (sidérites) sont des alliages \ce{Fe-Ni}, analogues probables du noyau. Le fer et le nickel sont des éléments \cle{sidérophiles} : ils ont migré vers le centre lors de la différenciation de la Terre jeune.
\item \textbf{L'expérience de haute pression} : aux conditions du noyau, le fer pur serait \emph{trop dense} d'environ 5 à 10 \% par rapport à la densité mesurée par les ondes sismiques. Il faut donc y mélanger des \cle{éléments légers} : soufre (\ce{S}), oxygène (\ce{O}), silicium (\ce{Si}), hydrogène (\ce{H}), qui « diluent » l'alliage \ce{Fe-Ni}.
\end{itemize}

\begin{aretenir}[Composition et état du noyau]
Le noyau est un alliage \textbf{fer--nickel} (environ 85 \% \ce{Fe}, 5 \% \ce{Ni}) additionné d'éléments légers (\ce{S}, \ce{O}, \ce{Si}, \ce{H}). Le \textbf{noyau externe} (2 900 à 5 150 km) est \textbf{liquide} ($\mu = 0$, $v_S = 0$) ; la \textbf{graine} ou noyau interne (5 150 à 6 371 km) est \textbf{solide}, solidifiée par la pression malgré la température élevée.
\end{aretenir}

\courssub{Le noyau liquide et le champ magnétique terrestre : la géodynamo}

Pourquoi la liquidité du noyau externe est-elle si importante pour nous, habitants de la surface ? Parce qu'un fluide conducteur d'électricité en mouvement engendre un champ magnétique. Le noyau externe de fer liquide est le siège de \cle{mouvements de convection} : le fer chaud monte, le fer refroidi au contact du manteau descend, et la rotation terrestre (force de Coriolis) enroule ces mouvements en colonnes hélicoïdales. Ce système fonctionne comme une \cle{dynamo auto-entretenue} : c'est la \cle{géodynamo}, source du \cle{champ magnétique terrestre}.

Ce champ magnétique forme autour de la Terre la \emph{magnétosphère}, bouclier qui dévie le \emph{vent solaire} (flux de particules chargées). Sans lui, l'atmosphère serait progressivement érodée — comme cela s'est produit sur Mars, dont la dynamo s'est éteinte — et la vie à la surface serait exposée à des radiations mortelles. La liquidité du noyau externe, démontrée par les ondes sismiques, est donc une \emph{condition d'habitabilité} de notre planète : un enjeu majeur d'éducation à l'environnement.

\courssub{Savoir-faire : lire un sismogramme et exploiter les données}

\begin{methode}[Analyser un sismogramme enregistré à $\Delta = 120^{\circ}$]
Une station située à $120^{\circ}$ de l'épicentre se trouve dans la zone d'ombre P \emph{et} dans la zone d'ombre S. Pour interpréter son sismogramme :
\begin{enumerate}
\item \textbf{Repérer les ondes P directes} : elles doivent être \emph{absentes} ou d'amplitude extrêmement faible (seules subsistent des ondes diffractées rampantes le long de la surface du noyau, arrivant tardivement).
\item \textbf{Vérifier l'absence totale d'onde S} : aucune arrivée S nette ne doit apparaître ; si une phase ressemblant à S existe, il s'agit d'une phase convertie (ex. P convertie en S à une interface) ou de bruit.
\item \textbf{Identifier les phases tardives} : les ondes de surface (ondes de Love et Rayleigh, grandes amplitudes dispersives) et éventuellement une PKP diffractée.
\item \textbf{Conclure} : la combinaison « pas de P directe + pas de S » à $120^{\circ}$ confirme que le rayon a rencontré le noyau externe liquide.
\end{enumerate}
\end{methode}

\begin{attention}[La zone d'ombre P n'est pas un silence absolu]
Erreur fréquente au Probatoire : écrire qu'« aucune onde » n'arrive entre $103^{\circ}$ et $143^{\circ}$. En réalité, l'\emph{amplitude des ondes P directes chute drastiquement}, mais la zone n'est pas totalement vide : des ondes \emph{diffractées} (P diffractée rampant le long de la discontinuité de Gutenberg), des ondes réfléchies sur le noyau (PcP) et le bruit de fond y sont enregistrés. C'est d'ailleurs l'analyse de ces faibles arrivées qui a permis à Inge Lehmann de découvrir la graine. Rédige donc : « les ondes P directes disparaissent (amplitude quasi nulle) entre $103^{\circ}$ et $143^{\circ}$ », et non « aucune onde n'est enregistrée ». Autre piège : ne confonds pas la discontinuité de \textbf{Gutenberg} (2 900 km, manteau/noyau externe) avec celle de \textbf{Lehmann} (5 150 km, noyau externe/graine).
\end{attention}

\begin{exoresolu}[Profondeur de pénétration d'un rayon PKP]
\textbf{Énoncé.} Un rayon PKP émerge à la distance angulaire $\Delta = 140^{\circ}$ avec un temps de parcours total d'environ $t = \SI{1100}{s}$. On modélise son trajet par trois segments rectilignes : deux traversées du manteau (longueur totale $L_m \approx \SI{4400}{km}$ à la vitesse moyenne $v_P = \SI{11}{km/s}$) et une traversée du noyau externe (longueur $L_k$ à la vitesse moyenne $v_K = \SI{9,5}{km/s}$). (1) Calculer la longueur $L_k$ du trajet dans le noyau. (2) En assimilant le trajet dans le noyau à une corde du cercle de rayon $R_N = \SI{3480}{km}$, en déduire la profondeur maximale atteinte (point de rebroussement).
\tcblower
\textbf{Solution détaillée.}
(1) Le temps de parcours est la somme des temps dans chaque couche :
\[ t = \frac{L_m}{v_P} + \frac{L_k}{v_K} \]
Temps passé dans le manteau :
\[ t_m = \frac{4400}{11} = \SI{400}{s} \]
Temps passé dans le noyau externe :
\[ t_k = 1100 - 400 = \SI{700}{s} \]
D'où la longueur du trajet dans le noyau :
\[ L_k = v_K \times t_k = 9,5 \times 700 = \SI{6650}{km} \]
(2) Le trajet dans le noyau est une corde de longueur $L_k = \SI{6650}{km}$ du cercle de rayon $R_N = \SI{3480}{km}$. La distance $d$ du centre de la Terre au milieu de la corde (point de rebroussement) vérifie, d'après le théorème de Pythagore :
\[ d = \sqrt{R_N^2 - \left(\frac{L_k}{2}\right)^2} = \sqrt{3480^2 - 3325^2} \]
\[ d = \sqrt{12\,110\,400 - 11\,055\,625} = \sqrt{1\,054\,775} \approx \SI{1027}{km} \]
La profondeur maximale atteinte est donc :
\[ p = R_T - d = 6371 - 1027 \approx \SI{5344}{km} \]
\textbf{Interprétation.} Le rayon a pénétré jusqu'à environ 5 340 km de profondeur. Le toit de la graine (discontinuité de Lehmann) se situe à 5 150 km de profondeur (rayon 1 221 km). Le point de rebroussement, à 1 027 km du centre, est \textbf{à l'intérieur de la graine solide} (rayon < 1 221 km). Il s'agit donc d'un rayon de la branche \textbf{PKIKP (PKPdf)}, qui traverse le noyau interne. Un rayon PKPbc, lui, aurait son point de rebroussement dans le noyau externe (rayon > 1 221 km, profondeur < 5 150 km).
\end{exoresolu}

\begin{exercice}[Raisonnement sur la zone d'ombre]
Un séisme se produit au Japon. La station de Yaoundé se trouve à une distance angulaire $\Delta \approx 115^{\circ}$ de l'épicentre. (1) Prévois, en justifiant, si la station de Yaoundé enregistrera une onde S directe. (2) Même question pour une onde P directe de forte amplitude. (3) Quelle phase sismique une station située à $\Delta = 160^{\circ}$ enregistrera-t-elle en premier, et que prouve son existence ?
\end{exercice}

\begin{corrige}[Exercice : Raisonnement sur la zone d'ombre]
(1) \textbf{Non.} $\Delta = 115^{\circ} > 103^{\circ}$ : Yaoundé est dans la zone d'ombre des ondes S. Tout rayon S direct devrait traverser le noyau externe, liquide ($\mu = 0$, $v_S = 0$) : l'onde S est bloquée. (2) \textbf{Non plus.} $103^{\circ} < 115^{\circ} < 143^{\circ}$ : la station est dans la zone d'ombre des ondes P. La chute de $v_P$ à la discontinuité de Gutenberg réfracte les rayons vers le centre ; seules de faibles ondes diffractées peuvent arriver. (3) À $\Delta = 160^{\circ}$, la première arrivée est la phase \textbf{PKIKP} (PKPdf), qui a traversé le manteau, le noyau externe liquide et la graine. Son existence — et sa précocité — prouve que le noyau interne est \textbf{solide}, avec $v_P \approx \SI{11}{km/s}$, supérieure à celle du noyau externe.
\end{corrige}

\begin{aretenir}[L'essentiel de la section]
\begin{itemize}
\item $\Delta > 103^{\circ}$ : absence totale d'ondes S $\Rightarrow$ $\mu = 0$ $\Rightarrow$ \textbf{noyau externe liquide}.
\item $103^{\circ} < \Delta < 143^{\circ}$ : zone d'ombre P due à la chute de $v_P$ (13,7 $\to$ 8,1 km/s) à la discontinuité de Gutenberg (2 900 km).
\item $\Delta > 143^{\circ}$ : arrivée des phases PKP puis PKIKP $\Rightarrow$ \textbf{graine solide} (discontinuité de Lehmann, 5 150 km ; découverte : Inge Lehmann, 1936).
\item Composition : alliage \ce{Fe-Ni} + éléments légers (\ce{S}, \ce{O}, \ce{Si}, \ce{H}) ; $\rho \approx 10$ à $\SI{13}{g/cm^3}$.
\item La convection du fer liquide dans le noyau externe engendre le champ magnétique terrestre (\textbf{géodynamo}), bouclier protecteur de la vie.
\end{itemize}
\end{aretenir}

\coursec{Modèle de référence : Structure en couches de la Terre (Modèle PREM / AK135)}

Dans la section précédente, tu as vu comment les ondes sismiques — et en particulier l'ombre des ondes S et la déviation des ondes P — ont permis de démontrer que la Terre possède un \cle{noyau externe liquide} entourant une \cle{graine solide}. Mais ces découvertes ponctuelles ne suffisent pas : pour exploiter les millions de sismogrammes enregistrés chaque année, les géophysiciens ont besoin d'un \emph{modèle complet} de la planète, donnant à chaque profondeur la vitesse des ondes, la densité et la pression. C'est exactement ce qu'est le \cle{modèle PREM}. Cette section rassemble tout ce que les ondes nous ont appris en un tableau de bord unique de l'intérieur du globe — un peu comme le bilan sanguin complet qu'un médecin du Centre Hospitalier d'Essos à Yaoundé établit après une série d'analyses séparées : chaque examen isolé dit peu de chose, mais leur synthèse donne le portrait fidèle du patient.

\courssub{Le modèle PREM : une radiographie standard de la planète}

\begin{definition}[Modèle PREM]
Le \cle{PREM} (\emph{Preliminary Reference Earth Model}), établi par \textbf{Dziewonski et Anderson en 1981}, est un \textbf{modèle de référence unidimensionnel (1D)} de la Terre : il donne, en fonction du seul rayon $r$ (ou de la profondeur), les valeurs moyennes de la vitesse des ondes P ($v_P$), des ondes S ($v_S$), de la densité ($\rho$), de la pression et de la gravité. Il a été ajusté sur des millions de temps de parcours d'ondes sismiques et sur les fréquences propres d'oscillation de la Terre entière après les grands séismes. Le modèle \cle{AK135} (Kennett et al., 1995) est sa version modernisée pour les temps de parcours.
\end{definition}

\begin{attention}[Ce que « 1D » signifie — et ne signifie pas]
Un modèle 1D suppose la Terre \textbf{parfaitement concentrique} : à une profondeur donnée, les propriétés sont supposées identiques partout. C'est une moyenne : la croûte sous le Mont Cameroun n'a évidemment pas la même épaisseur que sous l'océan Atlantique voisin. Le PREM est un \emph{cadre de référence}, pas une photographie exacte de chaque point du globe.
\end{attention}

\courssub{La croûte : une pellicule à deux visages}

\begin{definition}[Discontinuité chimique vs transition de phase]
Une \cle{discontinuité chimique} sépare deux matériaux de \textbf{compositions différentes} (ex. : le Moho, entre roches granitiques/basaltiques et péridotites du manteau). Une \cle{transition de phase} concerne un matériau de \textbf{même composition} dont les minéraux changent de structure cristalline sous l'effet de la pression (ex. : olivine $\to$ wadsleyite à 410 km). Les deux provoquent un saut de vitesse des ondes, mais leur signification géologique est très différente.
\end{definition}

La croûte est la seule enveloppe directement accessible (forages, mines), mais elle est \textbf{bimodale} :

\begin{itemize}
\item \textbf{Croûte océanique} : mince (5 à 10 km), jeune, dense ($\rho \approx \num{2,9}$), de composition \textbf{basaltique} (riche en fer et magnésium). Au Moho océanique, $v_P$ saute d'environ 6,8 à 8,1 km/s.
\item \textbf{Croûte continentale} : épaisse (30 à 70 km, jusqu'à 70 km sous les chaînes de montagnes), ancienne, légère ($\rho \approx \num{2,7}$), de composition moyenne \textbf{granitique} (riche en silice et aluminium).
\end{itemize}

\begin{attention}[Deux pièges classiques au Probatoire]
\begin{enumerate}
\item La croûte n'est \textbf{pas une couche uniforme} : écrire « la croûte fait 30 km d'épaisseur » sans préciser océanique ou continentale est une erreur.
\item Le \textbf{Moho n'est pas la base de la lithosphère} ! Le Moho est une discontinuité \emph{chimique} (changement de roche), alors que la limite lithosphère/asthénosphère est une limite \emph{rhéologique} (changement de comportement mécanique, rigide $\to$ ductile), située plus profond (100 à 250 km).
\end{enumerate}
\end{attention}

\courssub{Le manteau : le moteur de la dynamique terrestre}

\begin{definition}[Lithosphère et asthénosphère : une distinction rhéologique]
La \cle{lithosphère} ($\approx$ 0 à 100 km sous les océans, jusqu'à 250 km sous les vieux continents) regroupe la croûte \emph{et} la partie supérieure rigide du manteau : elle se comporte comme un solide \textbf{rigide} qui se brise (séismes). L'\cle{asthénosphère} (en dessous, jusqu'à $\approx$ 350--400 km) est une zone \textbf{ductile} : les roches y sont proches de leur point de fusion, elles s'écoulent lentement. Elle se trahit par une \textbf{zone de basse vitesse} (LVZ) : $v_S$ y chute de quelques pourcents, car un matériau partiellement ramolli transmet moins bien les ondes de cisaillement.
\end{definition}

Plus profond, le manteau est découpé par deux discontinuités majeures qui ne sont \emph{pas} chimiques :

\begin{itemize}
\item \textbf{Discontinuité à 410 km} : transition de phase olivine $\to$ wadsleyite (structure cristalline plus compacte) $\Rightarrow$ saut de $v_P$, $v_S$ et $\rho$.
\item \textbf{Discontinuité à 660 km} : décomposition de la ringwoodite en pérovskite + magnésiowüstite $\Rightarrow$ nouveau saut. Elle marque la limite entre \textbf{manteau supérieur} et \textbf{manteau inférieur}.
\end{itemize}

Le \textbf{manteau inférieur} (660 à 2891 km) est chimiquement homogène (composition \emph{pyrolitique}, proche des péridotites) : les vitesses et la densité y augmentent \emph{progressivement}, sans saut brutal, car la pression comprime simplement les minéraux. À sa base s'étend la couche \cle{D''} (200 à 300 km d'épaisseur), hétérogène et complexe, siège de panaches thermiques et de zones à ultra-basse vitesse — c'est le « cimetière » possible des plaques de lithosphère subduites.

\courssub{Le noyau : cœur liquide et graine solide}

\begin{itemize}
\item \textbf{Noyau externe} (2891 à 5150 km) : \textbf{liquide} ($v_S = 0$), composé de fer-nickel allié à des éléments légers (S, O, Si). $v_P \approx 8$ à 10 km/s. Sa \textbf{convection turbulente} engendre le champ magnétique terrestre par effet \textbf{dynamo} (la \emph{géodynamo}) — le bouclier qui protège Douala comme tout le globe des particules solaires.
\item \textbf{Noyau interne (graine)} (5150 à 6371 km) : \textbf{solide} ($v_S \approx 3,5$ km/s réapparaissent), fer-nickel presque pur, $v_P \approx 11$ km/s. Il \textbf{grandit lentement} par solidification du noyau externe ; la chaleur latente libérée entretient la convection du noyau liquide. Il est \textbf{anisotrope} : les ondes P voyagent environ 3 \% plus vite le long de l'axe de rotation polaire que dans le plan équatorial.
\end{itemize}

\begin{exemplebox}[La masse du noyau et le moment d'inertie]
Le noyau (externe + interne) ne représente que $\approx 16\,\%$ du volume terrestre, mais environ \textbf{un tiers de la masse} de la planète ($\approx \num{1,94e24}$ kg sur $\num{5,97e24}$ kg), tant le fer est dense ($\rho$ de 9,9 à 13,1). Preuve indirecte élégante : le \textbf{moment d'inertie} de la Terre vaut $I/(MR^2) \approx \num{0,33}$, alors qu'une sphère homogène donnerait $0,4$. Une valeur plus faible signifie que la masse est \textbf{concentrée vers le centre} : la Terre est bien différenciée en couches de densité croissante. Cette contrainte astronomique (mesurée via la précession) confirme indépendamment le modèle sismique.
\end{exemplebox}

\courssub{Le modèle standard en un coup d'œil}

\begin{popfigure}
\begin{center}
\begin{tikzpicture}[scale=1]
% Échelle : 1 cm = 1500 km. Rayon total = 6371 km = 4.247 cm.
% Rayons depuis le centre :
% Noyau interne (graine) : 1221 km = 0.814 cm
% Noyau externe (CMB / Gutenberg) : 3480 km = 2.320 cm
% Discontinuité 660 km : 5711 km = 3.807 cm
% Surface : 6371 km = 4.247 cm
% (410 km = 5961 km = 3.974 cm, trop proche de 660 km pour être distinct visuellement)

% Remplissages (du plus grand au plus petit pour l'empilement visuel)
\fill[popgreenL] (0,0) circle (3.807); % Manteau (jusqu'à 660 km)
\fill[poporangeL] (0,0) circle (2.320); % Noyau externe
\fill[popgold] (0,0) circle (0.814);   % Noyau interne
\fill[popblueL] (0,0) circle (4.247);  % Croûte + Manteau supérieur (exagéré pour visibilité)

% Contours
\draw[popdark,thick] (0,0) circle (0.814);   % Limite graine
\draw[popdark,thick] (0,0) circle (2.320);   % Limite noyau externe (Gutenberg)
\draw[popdark,thick] (0,0) circle (3.807);   % Discontinuité 660 km
\draw[popdark,thick] (0,0) circle (4.247);   % Surface

% Étiquettes
\draw[popdark,thin,->] (0.3,0.3) -- (2.6,1.6) node[right,align=left,font=\small]{\textbf{Noyau interne} (solide, Fe-Ni)\\ $v_P \approx 11$ km/s ; $v_S \approx 3,5$ km/s ; $\rho \approx 13$};
\draw[popdark,thin,->] (-1.5,-1.0) -- (-4.6,-2.4) node[below,align=center,font=\small]{\textbf{Noyau externe} (liquide)\\ $v_P \approx 8$--$10$ km/s ; $v_S = 0$ ; $\rho \approx 10$--$12$};
\draw[popdark,thin,->] (3.0,-2.0) -- (5.6,-3.2) node[below,align=center,font=\small]{\textbf{Manteau inférieur} (solide)\\ $v_P \approx 11$--$13,7$ km/s ; $\rho \approx 4,4$--$5,6$\\ \textit{Discontinuité 660 km}};
\draw[popdark,thin,->] (-3.5,2.5) -- (-5.4,3.4) node[above,align=center,font=\small]{\textbf{Manteau sup. + Croûte}\\ $v_P \approx 6$--$11$ km/s ; $\rho \approx 2,7$--$4,4$\\ \textit{(Croûte exagérée)} };
\node[font=\small\itshape,popmuted] at (0,-4.9) {Échelle : 1 cm $\approx$ 1500 km — Rayon total 6371 km};
\end{tikzpicture}
\end{center}
\legende{Coupe schématique de la Terre à l'échelle (croûte exagérée). Rayons proportionnels : graine 1221 km, noyau externe 2259 km d'épaisseur, manteau inférieur jusqu'à 660 km. Le noyau occupe plus de la moitié du rayon terrestre. Les discontinuités à 410 km et 660 km (transitions de phase) sont indiquées ; la 410 km est trop proche de la 660 km pour être tracée distinctement à cette échelle.}
\end{popfigure}

\begin{aretenir}[Tableau du modèle standard (à connaître par cœur)]
\begin{center}
\small
\begin{tabular}{|l|c|c|c|c|c|}
\hline
\rowcolor{popblueL} \textbf{Enveloppe} & \textbf{Profondeur (km)} & \textbf{Épaisseur (km)} & \textbf{État} & $\boldsymbol{\rho}$ & $\boldsymbol{v_P}$ (km/s) / $\boldsymbol{v_S}$ (km/s) \\
\hline
Croûte océanique & 0 -- 7 & 5 -- 10 & solide & 2,9 & 6,8 / 3,9 \\
\hline
Croûte continentale & 0 -- 35 & 30 -- 70 & solide & 2,7 & 6,0 / 3,5 \\
\hline
Manteau supérieur & Moho -- 660 & $\approx$ 650 & solide (ductile dans l'asthénosphère) & 3,3 -- 4,4 & 8 -- 11 / 4,5 -- 6 \\
\hline
Manteau inférieur & 660 -- 2891 & $\approx$ 2230 & solide & 4,4 -- 5,6 & 11 -- 13,7 / 6 -- 7,3 \\
\hline
Noyau externe & 2891 -- 5150 & $\approx$ 2260 & \textbf{liquide} & 9,9 -- 12,2 & 8 -- 10 / \textbf{0} \\
\hline
Noyau interne & 5150 -- 6371 & $\approx$ 1220 (rayon) & \textbf{solide} & 12,8 -- 13,1 & 11 / 3,5 \\
\hline
\end{tabular}
\end{center}
Repères à retenir : Moho (base croûte), 410 km et 660 km (transitions de phase), 2891 km (Gutenberg, limite manteau/noyau), 5150 km (Lehmann, limite noyau externe/interne).
\end{aretenir}

\begin{popfigure}
\begin{center}
\begin{tikzpicture}
\begin{axis}[
  ybar, bar width=0.55cm,
  width=0.95\linewidth, height=6.2cm,
  symbolic x coords={Croûte, Manteau, Noyau ext., Noyau int.},
  xtick=data,
  ylabel={Épaisseur (km)},
  ymin=0, ymax=3200,
  nodes near coords,
  every node near coord/.append style={font=\small},
  axis lines=left,
  ymajorgrids=true, grid style={popline!40},
  x tick label style={font=\small},
]
\addplot[fill=popblue,draw=popdark] coordinates {(Croûte,35) (Manteau,2885) (Noyau ext.,2259) (Noyau int.,1221)};
\end{axis}
\end{tikzpicture}
\end{center}
\legende{Comparaison des épaisseurs des enveloppes terrestres (croûte continentale moyenne de 35 km). Le manteau domine en volume ($\approx 84\,\%$ du volume terrestre), mais le noyau (externe + interne) domine en masse relative par unité de volume grâce à sa densité de 10 à 13.}
\end{popfigure}

\courssub{Construire et exploiter le modèle}

\begin{methode}[Construire le modèle à l'échelle sur une feuille A4]
\begin{enumerate}
\item Choisir l'échelle : \textbf{1 cm pour 500 km} donne un rayon terrestre de $6371/500 \approx 12,7$ cm — parfait pour une demi-feuille A4.
\item Reporter les rayons des limites depuis le centre : graine $1221/500 \approx 2,4$ cm ; surface du noyau $(6371-2891)/500 \approx 7,0$ cm ; discontinuité 660 km : $(6371-660)/500 \approx 11,4$ cm ; discontinuité 410 km : $\approx 11,9$ cm ; surface : 12,7 cm.
\item Tracer les cercles concentriques au compas, colorier chaque enveloppe d'une couleur différente.
\item Annoter chaque couche : profondeur, état physique, $v_P$, $v_S$, $\rho$.
\item Observer : la croûte (0,07 mm à cette échelle !) est invisible — c'est la leçon la plus parlante de l'exercice.
\end{enumerate}
\end{methode}

\begin{exoresolu}[Exploiter le modèle PREM]
\textbf{Énoncé.} Un séisme se produit à 10 km de profondeur. Une station enregistre une onde P ayant traversé le manteau supérieur sur 1500 km à la vitesse moyenne de 10 km/s, puis une onde S sur le même trajet à 5,5 km/s. (1) Calculer la différence de temps d'arrivée $\Delta t = t_S - t_P$. (2) Pourquoi cette onde S n'a-t-elle pas pu traverser le noyau externe ? (3) À quelle discontinuité du modèle PREM correspond la limite située à 660 km, et quelle en est la nature ?
\tcblower
\textbf{Solution.} (1) $t = d/v$ :
\begin{align*}
t_P &= \frac{1500}{10} = 150 \text{ s} & t_S &= \frac{1500}{5,5} \approx 272,7 \text{ s}
\end{align*}
donc $\Delta t \approx 272,7 - 150 = 122,7$ s, soit environ \textbf{2 minutes}. C'est ce retard croissant avec la distance qui permet de localiser l'épicentre. (2) Le noyau externe est \textbf{liquide} ; or les ondes S (cisaillement) ne se propagent pas dans les fluides car ceux-ci n'ont pas de rigidité : $v_S = 0$ dans le noyau externe. (3) À 660 km se trouve une \textbf{transition de phase minéralogique} (ringwoodite $\to$ pérovskite + magnésiowüstite) : la composition chimique ne change pas, mais la structure cristalline se compacte, d'où un saut de $v_P$, $v_S$ et $\rho$. C'est la frontière entre manteau supérieur et manteau inférieur.
\end{exoresolu}

\begin{savaistu}[Le cœur de la Terre tourne plus vite que sa surface]
En comparant les temps de parcours d'ondes \textbf{PKIKP} (qui traversent la graine) émises par des séismes quasi identiques séparés par des décennies, les sismologues ont découvert que le \textbf{noyau interne tourne légèrement plus vite que le manteau et la croûte} : c'est la \emph{super-rotation}, de l'ordre de 0,1 à 0,3 degré par an (des études récentes suggèrent qu'elle varie, voire s'inverse périodiquement). Entraînée par le champ magnétique généré dans le noyau liquide, la graine ferait un tour d'avance sur nous en quelques siècles. Ainsi, trente ans d'archives sismiques suffisent à « voir » tourner le cœur de la planète — sans jamais y descendre.
\end{savaistu}

\begin{exercice}[Pour t'entraîner]
À partir du tableau du modèle standard : (1) classer les enveloppes par densité croissante ; (2) calculer la profondeur de la limite noyau externe / noyau interne en partant du rayon de la graine (1221 km) ; (3) expliquer pourquoi $v_P$ chute brutalement de 13,7 à 8 km/s à 2891 km de profondeur alors que la densité, elle, augmente.
\end{exercice}

\begin{corrige}[Exercice]
(1) Croûte continentale (2,7) $<$ croûte océanique (2,9) $<$ manteau supérieur (3,3--4,4) $<$ manteau inférieur (4,4--5,6) $<$ noyau externe (9,9--12,2) $<$ noyau interne (12,8--13,1). (2) Profondeur $= 6371 - 1221 = 5150$ km : c'est la discontinuité de Lehmann. (3) À 2891 km (discontinuité de Gutenberg), on passe d'un solide (manteau) à un \textbf{liquide} (noyau externe). La vitesse des ondes P dépend à la fois de la densité \emph{et} des modules d'élasticité : dans un liquide, le module de cisaillement est nul, ce qui fait chuter $v_P$ malgré l'augmentation de densité. C'est précisément cette chute de $v_P$ couplée à $v_S = 0$ qui prouve l'état liquide.
\end{corrige}

En résumé, le modèle PREM/AK135 condense un siècle de sismologie en une « carte d'identité » de la planète : une croûte bimodale et infime, un manteau solide rythmé par des transitions de phase à 410 et 660 km, un noyau externe liquide siège de la géodynamo, et une graine solide anisotrope en croissance. Mais cette moyenne 1D lisse les différences régionales : la section suivante lèvera le voile sur les hétérogénéités et la dynamique que cache cette belle symétrie sphérique.

\coursec{Au-delà du modèle 1D : Hétérogénéités, Anisotropie et Dynamique (Complément Probatoire)}

Dans la section précédente, nous avons construit un modèle de la Terre en couches concentriques (croûte, manteau, noyau) où la vitesse des ondes ne dépend que de la profondeur : c'est le modèle \cle{1D} (une dimension), dit \cle{sphériquement symétrique}, comme le modèle PREM. Ce modèle est remarquablement efficace pour situer les grandes discontinuités. Mais il repose sur une simplification forte : il suppose que deux points situés à la même profondeur, l'un sous le Cameroun et l'autre sous le Pacifique, sont strictement identiques. Or l'observation de la surface — volcans, dorsales, fosses de subduction — suggère le contraire. La Terre réelle est tridimensionnelle, hétérogène et dynamique. Cette dernière section montre comment les sismologues dépassent le modèle 1D pour réaliser un véritable \og scanner \fg{} du globe, et comment les images obtenues éclairent le moteur de la tectonique des plaques.

\courssub{Les limites du modèle 1D : la Terre est tridimensionnelle}

\begin{experience}[Une observation qui dérange le modèle]
Deux séismes de même magnitude se produisent à la même profondeur, l'un sous les Andes, l'autre sous une dorsale océanique. Une station sismique équidistante des deux enregistre les ondes P du premier séisme environ $8\ \mathrm{s}$ plus tôt que celles du second. Pourtant, selon le modèle PREM, les temps de parcours devraient être identiques, puisque la distance et la profondeur sont les mêmes.
\end{experience}

\textbf{Hypothèse.} Les roches traversées ne sont pas les mêmes dans les deux cas : à profondeur égale, la vitesse des ondes varie \emph{latéralement} (horizontalement).

\textbf{Interprétation.} Sous les Andes, une plaque océanique froide plonge dans le manteau (subduction) : les roches froides sont plus rigides, les ondes y vont plus vite. Sous la dorsale, le manteau chaud remonte : les roches chaudes, parfois partiellement fondues, ralentissent les ondes. Le modèle 1D, qui moyenne toutes ces situations, ne peut pas rendre compte de ces écarts.

\begin{aretenir}[À retenir]
Le modèle 1D (PREM) donne la vitesse moyenne $v$ en fonction de la seule profondeur. La Terre réelle présente des \cle{hétérogénéités latérales} : à profondeur constante, la vitesse varie d'un point à un autre. Ces écarts, faibles (quelques pourcents), portent une information précieuse sur la température, la composition et la dynamique interne.
\end{aretenir}

\courssub{La tomographie sismique : le scanner de la Terre}

En médecine, le scanner (tomodensitométrie) reconstitue une image 3D de l'intérieur du corps en croisant des milliers de mesures de rayons X traversant le patient dans toutes les directions. Les sismologues font exactement la même chose avec la Terre : les rayons X sont remplacés par les ondes sismiques, et le \og tube à rayons X \fg{} par les séismes naturels.

\begin{definition}[Tomographie sismique]
La \cle{tomographie sismique} est une technique d'imagerie qui reconstitue la distribution tridimensionnelle des vitesses des ondes à l'intérieur de la Terre, par \emph{inversion} d'un très grand nombre de temps de parcours mesurés entre des séismes et des stations répartis sur toute la planète (réseaux mondiaux de sismomètres). Le résultat est une carte des \cle{anomalies de vitesse} $\delta v/v$, exprimées en pourcentage par rapport au modèle 1D de référence.
\end{definition}

\begin{methode}[Principe de l'inversion tomographique]
\begin{enumerate}
\item \textbf{Mesurer.} Pour chaque couple séisme--station, on mesure le temps de parcours réel $t_{\text{obs}}$ d'une phase (P, S, PP\ldots) et on calcule le temps $t_{\text{PREM}}$ prédit par le modèle 1D. L'écart $\delta t = t_{\text{obs}} - t_{\text{PREM}}$ est appelé \emph{résidu}.
\item \textbf{Modéliser.} On découpe l'intérieur de la Terre en blocs (ou en mailles sphériques). Le résidu s'écrit comme la somme des contributions des blocs traversés par le rayon :
\[ \delta t \;=\; \int_{\text{rayon}} K\,\frac{\delta v}{v}\,\mathrm{d}s \]
où $K$ est un noyau de sensibilité (lié à la géométrie du rayon) et $\delta v/v$ l'anomalie de vitesse inconnue de chaque bloc.
\item \textbf{Inverser.} Avec des millions de rayons croisés (des millions d'équations), on résout le système géant par une méthode \emph{itérative} : on part du modèle 1D, on ajuste les vitesses pour réduire les résidus, et on recommence jusqu'à convergence.
\item \textbf{Interpréter.} Les anomalies rapides ($\delta v/v > 0$, souvent colorées en bleu) sont associées à du matériau froid et rigide ; les anomalies lentes ($\delta v/v < 0$, en rouge) à du matériau chaud, hydraté ou partiellement fondu.
\end{enumerate}
\end{methode}

\begin{exemplebox}[Lecture d'une coupe tomographique]
Sous l'Amérique du Sud, les tomographies montrent une anomalie \emph{rapide} ($\delta v_P/v_P \approx +1\ \%$ à $+2\ \%$) en forme de lame plongeant vers l'est jusqu'à plus de $1\,000\ \mathrm{km}$ de profondeur : c'est la plaque de Nazca subductée, encore froide, qui s'enfonce dans le manteau. À l'inverse, sous la dorsale du Pacifique Est, une anomalie \emph{lente} ($-1\ \%$ à $-3\ \%$) marque la remontée de manteau chaud qui alimente le volcanisme de la dorsale.
\end{exemplebox}

\begin{popfigure}
\begin{tikzpicture}[scale=1.0]
% Fond : manteau
\fill[popcream] (0,0) rectangle (12,6);
% Surface
\draw[thick, popdark] (0,6) -- (12,6);
\node[above, popdark, font=\small] at (6,6.05) {Surface terrestre};
% Plaque subductée (anomalie rapide, bleue)
\fill[popblue!70] (1.2,6) -- (3.2,6) -- (7.6,0.6) -- (6.4,0.2) -- (2.4,5.2) -- (1.2,5.2) -- cycle;
\node[popblue, font=\small\bfseries] at (4.6,3.4) {Plaque subductée};
\node[popblue, font=\footnotesize] at (4.6,2.9) {froide : $\delta v/v > 0$ (rapide)};
% Flèche de descente
\draw[-{Stealth[length=3mm]}, very thick, popblue] (2.6,5.4) -- (5.6,1.6);
% Panache (anomalie lente, rouge/orange)
\fill[poporange!75, decorate, decoration={snake, amplitude=0.6mm}] (9.6,0) -- (10.0,3.2) -- (10.6,5.0) -- (11.0,6) -- (11.6,6) -- (11.2,4.6) -- (10.8,2.8) -- (10.6,0) -- cycle;
\fill[poporange!75] (9.6,0) rectangle (10.6,4.0);
\draw[-{Stealth[length=3mm]}, very thick, poporange] (10.1,0.4) -- (10.1,4.6);
\node[poporange!80!black, font=\small\bfseries, align=center] at (10.3,5.4) {Panache};
\node[poporange!80!black, font=\footnotesize, align=center] at (10.3,4.9) {chaud : $\delta v/v < 0$ (lent)};
% Volcan au-dessus du panache
\draw[fill=poporange!50] (10.4,6) -- (11.0,6.7) -- (11.6,6) -- cycle;
\node[popdark, font=\footnotesize] at (11.0,6.95) {Volcan de point chaud};
% Fosse au-dessus de la subduction
\draw[popdark, thick] (2.0,6) -- (2.2,5.75) -- (2.4,6);
\node[popdark, font=\footnotesize] at (2.2,6.3) {Fosse};
% Noyau
\fill[popgold!40] (0,0) rectangle (12,-0.9);
\draw[thick, popdark] (0,0) -- (12,0);
\node[popdark, font=\small] at (6,-0.45) {Noyau externe (limite à $2\,891\ \mathrm{km}$)};
% Échelle
\draw[<->, popmuted] (0.3,6) -- (0.3,0) node[midway, left, font=\footnotesize, rotate=90, yshift=2pt] {Profondeur croissante};
\end{tikzpicture}
\legende{Coupe tomographique conceptuelle (transect de type Pacifique). En bleu : anomalie de vitesse positive correspondant à une plaque océanique froide subductée qui plonge dans le manteau. En orange : anomalie de vitesse négative correspondant à un panache de manteau chaud qui remonte et alimente un volcan de point chaud. Ces images 3D relient directement la structure interne à la tectonique des plaques.}
\end{popfigure}

\courssub{Que signifient les anomalies de vitesse ?}

\begin{propriete}[Sensibilité différentielle des ondes P et S]
Les ondes S sont plus sensibles à la température et à la fusion partielle que les ondes P :
\[ \left|\frac{\delta v_S}{v_S}\right| \;>\; \left|\frac{\delta v_P}{v_P}\right| \]
En effet, la chaleur et la présence d'un film de liquide (fusion partielle) diminuent fortement le module de rigidité $\mu$ (dont dépend $v_S = \sqrt{\mu/\rho}$), alors que le module d'incompressibilité $K$ (qui intervient dans $v_P$) est moins affecté. Une anomalie thermique se traduit donc par un contraste plus marqué en vitesse S qu'en vitesse P.
\end{propriete}

\begin{exemplebox}[Exemple chiffré : l'anomalie sous l'Afrique]
Les tomographies globales montrent, à environ $2\,000\ \mathrm{km}$ de profondeur sous l'Afrique australe, une anomalie de vitesse S de l'ordre de $\delta v_S/v_S \approx -2\ \%$. Avec un coefficient de sensibilité thermique typique du manteau inférieur ($\delta v_S/v_S \approx -8\times10^{-5}\ \mathrm{K^{-1}}$ par kelvin), cette anomalie correspondrait à un excès de température d'environ $200$ à $300\ \mathrm{K}$ par rapport au manteau environnant — ou bien, en partie, à la présence de fusion partielle ou d'une composition différente. Cette structure chaude est l'une des deux grandes provinces lentes de la base du manteau.
\end{exemplebox}

\begin{attention}[Anomalie de vitesse $\neq$ anomalie de température uniquement]
Interpréter une anomalie de vitesse comme une simple anomalie de température est une erreur fréquente. La vitesse dépend aussi de :
\begin{itemize}
\item la \textbf{composition chimique} (enrichissement en fer, en silice\ldots) ;
\item la présence d'\textbf{eau} ou de fluides, qui ralentit les ondes ;
\item la \textbf{fusion partielle} (quelques pourcents de liquide suffisent à abaisser fortement $v_S$) ;
\item l'\textbf{anisotropie} (voir ci-dessous), qui peut simuler une hétérogénéité si on l'ignore.
\end{itemize}
Au Probatoire, précise toujours que l'interprétation thermique est l'hypothèse la plus probable, mais pas la seule.
\end{attention}

\courssub{L'anisotropie sismique : une vitesse qui dépend de la direction}

Jusqu'ici, nous avons supposé que la vitesse d'une onde ne dépend que du matériau traversé. Or dans de nombreuses roches du manteau, les cristaux ne sont pas orientés au hasard : ils sont \emph{alignés} par les écoulements, comme les troncs flottants qui s'alignent dans le courant d'une rivière. La vitesse dépend alors de la direction de propagation et de la polarisation de l'onde.

\begin{definition}[Anisotropie sismique]
Un milieu est \cle{anisotrope} lorsque la vitesse des ondes sismiques y dépend de la direction de propagation ou de la direction de polarisation. On distingue :
\begin{itemize}
\item l'\cle{anisotropie azimutale} : la vitesse dépend de la direction horizontale de propagation (par exemple, les ondes P vont plus vite dans la direction d'écoulement du manteau que perpendiculairement) ;
\item l'\cle{anisotropie radiale} : les ondes S polarisées horizontalement (SH) et verticalement (SV) se propagent à des vitesses différentes ($v_{SH} \neq v_{SV}$) dans un même matériau.
\end{itemize}
\end{definition}

\textbf{Origine physique.} Dans le manteau supérieur, le minéral dominant est l'\cle{olivine}, dont les cristaux sont anisotropes : les ondes se propagent plus vite le long de certains axes cristallographiques. Lorsque le manteau s'écoule horizontalement (asthénosphère entraînée par les plaques), les cristaux d'olivine s'orientent préférentiellement dans la direction du flux. Conséquence mesurable : dans l'asthénosphère, on observe souvent $v_{SH} > v_{SV}$. Dans le \textbf{noyau interne}, les cristaux de fer seraient alignés avec l'axe de rotation terrestre : les ondes P qui voyagent le long de l'axe nord--sud arrivent environ $3$ à $4\ \%$ plus vite que celles qui traversent le noyau dans le plan équatorial — un indice précieux sur la cristallisation et la dynamique du noyau.

\begin{popfigure}
\begin{tikzpicture}[scale=1.0]
% ---- Panneau gauche : anisotropie radiale ----
\begin{scope}
\fill[popgreenL] (0,0) rectangle (5.2,3.2);
\node[popdark, font=\small\bfseries] at (2.6,3.5) {Anisotropie radiale (asthénosphère)};
% cristaux alignés horizontalement
\foreach \y in {0.5,1.1,1.7,2.3}{
  \foreach \x in {0.7,1.9,3.1,4.3}{
    \draw[popgreen!60!black, thick, rotate around={0:(\x,\y)}] (\x-0.35,\y) -- (\x+0.35,\y);
  }
}
% ondes SH et SV
\draw[-{Stealth[length=2.5mm]}, very thick, popblue] (0.3,2.75) -- (2.0,2.75) node[midway, above, font=\footnotesize] {$v_{SH}$ rapide};
\draw[-{Stealth[length=2.5mm]}, very thick, poporange] (0.3,0.25) -- (1.6,0.25) node[midway, below, font=\footnotesize] {$v_{SV}$ plus lente};
\node[popdark, font=\footnotesize, align=center] at (3.9,2.75) {Cristaux d'olivine\\alignés par le flux};
\node[popdark, font=\footnotesize] at (2.6,-0.35) {$v_{SH} > v_{SV}$ : les ondes SH arrivent avant les SV};
\end{scope}
% ---- Panneau droit : anisotropie azimutale ----
\begin{scope}[xshift=7cm]
\fill[popblueL] (0,0) rectangle (5.2,3.2);
\node[popdark, font=\small\bfseries] at (2.6,3.5) {Anisotropie azimutale (manteau supérieur)};
% flèche de flux
\draw[-{Stealth[length=3mm]}, very thick, popdark] (0.3,1.6) -- (4.9,1.6) node[midway, above, font=\footnotesize] {Direction d'écoulement du manteau};
% onde rapide parallèle
\draw[-{Stealth[length=2.5mm]}, very thick, popblue] (0.3,2.6) -- (3.4,2.6) node[midway, above, font=\footnotesize] {$v$ rapide (parallèle au flux)};
% onde lente perpendiculaire
\draw[-{Stealth[length=2.5mm]}, very thick, poporange] (2.6,0.3) -- (2.6,1.2) node[midway, right, font=\footnotesize] {$v$ lente (perpendiculaire)};
\node[popdark, font=\footnotesize, align=center] at (2.6,-0.35) {La vitesse dépend de l'azimut de propagation};
\end{scope}
\end{tikzpicture}
\legende{Deux types d'anisotropie sismique. À gauche : anisotropie radiale — les cristaux d'olivine alignés horizontalement par l'écoulement asthénosphérique font que les ondes S polarisées horizontalement (SH) sont plus rapides que les ondes polarisées verticalement (SV). À droite : anisotropie azimutale — la vitesse dépend de la direction horizontale de propagation par rapport à la direction du flux mantellique.}
\end{popfigure}

\courssub{La couche D'' et les grandes provinces de la base du manteau}

La couche \cle{D''} désigne les $200$ à $300$ derniers kilomètres du manteau, juste au-dessus du noyau (entre environ $2\,600$ et $2\,891\ \mathrm{km}$ de profondeur). C'est une zone de contact entre deux mondes : le manteau silicaté solide et le noyau de fer liquide. La tomographie y révèle des structures spectaculaires.

\begin{definition}[LLSVP et ULVZ]
\begin{itemize}
\item Les \cle{LLSVP} (\emph{Large Low Shear Velocity Provinces}, \og grandes provinces à faible vitesse de cisaillement \fg{}) sont deux immenses régions de la base du manteau, l'une sous l'\textbf{Afrique}, l'autre sous le \textbf{Pacifique}, où $v_S$ est anormalement faible (jusqu'à $-2$ à $-3\ \%$). Elles s'étendent sur des milliers de kilomètres horizontalement et des centaines de kilomètres verticalement. On les appelle parfois \og superplumes \fg{} ou pôles superplumes.
\item Les \cle{ULVZ} (\emph{Ultra-Low Velocity Zones}) sont des poches beaucoup plus petites (quelques dizaines de kilomètres d'épaisseur) plaquées contre la surface du noyau, où les vitesses chutent de $10$ à $30\ \%$. Une chute aussi brutale suggère la présence de \textbf{fusion partielle} ou d'un matériau de composition très différente.
\end{itemize}
\end{definition}

\textbf{Interprétation dynamique.} Les LLSVP seraient des réservoirs de matériau chaud et dense, accumulé à la base du manteau, d'où s'élèveraient les \cle{panaches mantelliques} (ou plumes) : des colonnes étroites de roche chaude qui remontent jusqu'à la surface et alimentent les volcans de points chauds (Hawaï, Islande, et dans le contexte africain, le volcanisme de la ligne du Cameroun, du Mont Cameroun au golfe de Guinée). La couche D'' est ainsi le siège d'une intense \textbf{interaction noyau--manteau} : échanges de chaleur, peut-être réactions chimiques entre le fer liquide et les silicates.

\begin{savaistu}[Des superplumes vieilles de 200 millions d'années ?]
Les deux LLSVP (Afrique et Pacifique) semblent remarquablement \textbf{stables} : leur position coïncide avec celle des grands épisodes volcaniques des derniers $200$ à $300$ millions d'années, ce qui suggère qu'elles existaient déjà à l'époque des dinosaures. Certains géophysiciens pensent qu'elles contrôlent l'histoire de la tectonique des plaques : les continents se disperseraient au-dessus des superplumes et se rassembleraient au-dessus des zones froides, rythmant les \og supercycles de Wilson \fg{} (ouverture puis fermeture des océans, comme l'ouverture de l'Atlantique Sud qui a séparé l'Afrique de l'Amérique). L'Afrique, posée sur son superplume, serait aujourd'hui un continent \og surélevé \fg{} par la chaleur venue des profondeurs.
\end{savaistu}

\courssub{Le lien avec la tectonique des plaques : la convection mantellique}

La tomographie sismique fournit la preuve directe que la tectonique des plaques n'est pas un phénomène superficiel : elle plonge ses racines dans tout le manteau.

\begin{itemize}
\item \textbf{Subduction.} Les anomalies rapides (froides) dessinent les plaques plongeantes, traçables parfois jusqu'à la couche D'' : le moteur principal de la convection est le \emph{poids} de ces plaques froides qui s'enfoncent (\og slab pull \fg{}).
\item \textbf{Dorsales et panaches.} Les anomalies lentes (chaudes) marquent les remontées : larges sous les dorsales océaniques, étroites et profondes sous les points chauds (panaches issus des LLSVP).
\item \textbf{Convection mantellique.} L'ensemble dessine un gigantesque système de convection thermique : le manteau chaud remonte (anomalies lentes), se refroidit en surface (plaques lithosphériques), puis replonge (anomalies rapides). Cette convection, alimentée par la chaleur interne (radioactivité, chaleur primordiale, chaleur dégagée par le noyau), est le \textbf{moteur de la dynamique terrestre} : dérive des continents, séismes, volcans, formation des reliefs.
\end{itemize}

\begin{aretenir}[Bilan de la section]
\begin{itemize}
\item Le modèle 1D (PREM) est une moyenne : la Terre réelle est hétérogène en 3D.
\item La \cle{tomographie sismique} inverse des millions de temps de parcours pour cartographier les anomalies de vitesse $\delta v/v$ : rapide = froid (plaques subductées), lent = chaud (panaches, dorsales).
\item Les ondes S sont plus sensibles à la température et à la fusion partielle que les ondes P.
\item L'\cle{anisotropie} (olivine alignée dans le manteau, fer aligné dans le noyau interne) renseigne sur les écoulements et la cristallisation.
\item La couche D'' abrite les LLSVP (Afrique, Pacifique) et les ULVZ, racines probables des panaches.
\item Ces images relient structure interne et tectonique des plaques : la convection mantellique est le moteur commun.
\end{itemize}
\end{aretenir}

\begin{exoresolu}[Interpréter une anomalie tomographique]
\textbf{Énoncé.} Une coupe tomographique montre, sous une chaîne de volcans insulaires océaniques, une anomalie de vitesse $\delta v_S/v_S = -2{,}5\ \%$ s'étendant de la surface jusqu'à la base du manteau, avec $\delta v_P/v_P = -1\ \%$ seulement. (1) Cette anomalie est-elle rapide ou lente ? (2) Proposer une interprétation. (3) Pourquoi le contraste est-il plus fort en S qu'en P ?
\tcblower
\textbf{Solution.} (1) L'anomalie est négative, donc \textbf{lente} : les ondes s'y propagent moins vite que dans le modèle de référence. (2) Une anomalie lente, verticale, traversant tout le manteau et aboutissant à des volcans s'interprète comme un \textbf{panache mantellique} : une colonne de roche chaude (excès thermique de l'ordre de $200$ à $300\ \mathrm{K}$), éventuellement partiellement fondue, qui remonte depuis la base du manteau. (3) Les ondes S dépendent du module de rigidité $\mu$, très sensible à la température et à la présence de liquide ; les ondes P dépendent aussi du module d'incompressibilité $K$, moins affecté. D'où $|\delta v_S/v_S| > |\delta v_P/v_P|$, ce qui confirme l'origine thermique (et non purement compositionnelle) de l'anomalie.
\end{exoresolu}

\begin{exercice}[Pour s'entraîner]
Une station enregistre les ondes d'un séisme lointain. Les ondes SH arrivent $12\ \mathrm{s}$ avant les ondes SV, alors que le modèle 1D prévoit une arrivée simultanée. Expliquer ce que cette observation révèle sur le manteau traversé, et nommer le phénomène en cause.
\end{exercice}

\begin{corrige}[Exercice 1]
Les ondes SH (polarisées horizontalement) et SV (polarisées verticalement) traversent le même matériau mais à des vitesses différentes ($v_{SH} > v_{SV}$) : le milieu est donc \textbf{anisotrope}. Il s'agit d'une \cle{anisotropie radiale}, typique de l'asthénosphère où les cristaux d'olivine sont orientés préférentiellement horizontalement par l'écoulement du manteau. Le décalage de $12\ \mathrm{s}$ renseigne ainsi sur la direction et l'intensité de l'écoulement mantellique sous la station.
\end{corrige}

\newpage

\coursec{Activité d'intégration}

\coursec{NCT09 : Utilisation des données sismiques pour concevoir la structure du globe terrestre}

\courssub{Objectifs de l'activité}

À l'issue de cette activité d'intégration, tu devras être capable de :
\begin{itemize}
\item lire et exploiter un \cle{sismogramme} simplifié pour mesurer des temps de parcours d'ondes sismiques ;
\item comparer des mesures à un modèle de référence (AK135) et calculer des \cle{résidus} ;
\item utiliser une phase réfléchie (PcP) pour estimer la profondeur d'une \cle{discontinuité} ;
\item argumenter sur l'\cle{état physique} du noyau à partir de la zone d'ombre des ondes S et P ;
\item construire un \cle{modèle en couches} du globe terrestre, annoté et à l'échelle ;
\item rédiger une conclusion scientifique structurée, comme exigée au Probatoire.
\end{itemize}

\courssub{Situation}

\textbf{Un séisme au Chili secoue... Yaoundé !}

Le 27 février 2010, un violent séisme de magnitude 8,8 frappe la région de Concepción, au Chili. Quelques minutes plus tard, la station sismologique du réseau de l'Institut de Recherches Géologiques et Minières (IRGM), installée près de Yaoundé, enregistre le passage des ondes à travers le globe. La distance angulaire entre l'épicentre et la station est $\Delta = 85^\circ$. Un binôme d'élèves de Première C du lycée de Nkol-Eton est invité, dans le cadre du club sciences, à exploiter cet enregistrement pour « faire la radiographie » de la Terre, exactement comme le font les géophysiciens.

\textbf{Document 1 — Sismogramme simplifié enregistré à Yaoundé ($\Delta = 85^\circ$).}

On y lit les arrivées suivantes (temps comptés depuis l'origine du séisme) :

\begin{center}
\begin{tabularx}{\linewidth}{|l|X|X|}
\hline
\rowcolor{popblueL} \textbf{Phase} & \textbf{Description du trajet} & \textbf{Temps observé} \\
\hline
P & onde P directe (manteau) & 12 min 10 s \\
\hline
PP & P réfléchie une fois en surface & 16 min 30 s \\
\hline
PcP & P réfléchie sur le noyau & 13 min 05 s \\
\hline
S & onde S directe (manteau) & 22 min 00 s \\
\hline
ScS & S réfléchie sur le noyau & 23 min 50 s \\
\hline
\end{tabularx}
\end{center}

\textbf{Document 2 — Temps théoriques du modèle AK135 pour $\Delta = 85^\circ$ :} P : 12 min 00 s ; PP : 16 min 20 s ; PcP : 13 min 00 s ; S : 21 min 50 s ; ScS : 23 min 40 s.

\textbf{Document 3 — Données complémentaires.}
\begin{itemize}
\item Vitesse moyenne des ondes P dans le manteau inférieur : $v_P = \SI{13,7}{km.s^{-1}}$ ; dans le noyau externe, la vitesse des P chute brutalement à environ $\SI{8}{km.s^{-1}}$.
\item Aucune onde S directe n'est enregistrée au-delà de $\Delta = 103^\circ$ ; les ondes P directes disparaissent entre $103^\circ$ et $142^\circ$ (zone d'ombre).
\item Rayon terrestre moyen : $R_T = \SI{6371}{km}$ ; profondeur focale supposée nulle (séisme superficiel).
\item Pour une réflexion PcP, on admet le modèle simplifié : l'onde descend verticalement jusqu'au noyau et remonte, soit un trajet total $2 \times (R_T - R_N)$ parcouru à la vitesse moyenne $v_P$ du manteau, où $R_N$ est le rayon du noyau.
\end{itemize}

\textbf{Problème posé :} comment, à partir de ce seul enregistrement et de ces données, retrouver la structure en couches du globe terrestre et l'état physique de son noyau ?

\courssub{Consignes}

\begin{enumerate}
\item \textbf{Mesure et résidus.} Relever les temps de parcours observés du document 1, les convertir en secondes, puis calculer pour chaque phase le \cle{résidu} : $r = t_{\text{observé}} - t_{\text{théorique}}$ (AK135). Présenter les résultats dans un tableau. Commenter le signe et l'ordre de grandeur des résidus.
\item \textbf{Profondeur du noyau.} En utilisant le temps de la phase PcP et le modèle simplifié du document 3, calculer la profondeur $h = R_T - R_N$ de la discontinuité manteau--noyau. Comparer à la valeur admise (\SI{2891}{km}).
\item \textbf{État du noyau.} À l'aide du document 3, expliquer pourquoi l'absence d'ondes S au-delà de $103^\circ$ et la zone d'ombre des ondes P prouvent que le noyau externe est liquide. Rédiger l'argumentation en 5 à 8 lignes avec un schéma de trajets.
\item \textbf{Construction du modèle.} Sur papier millimétré, construire à l'échelle $1/100\,000\,000$ (1 cm = \SI{1000}{km}) le modèle en couches de la Terre : croûte ($0$ à \SI{35}{km}), manteau (jusqu'à \SI{2891}{km}), noyau externe (jusqu'à \SI{5150}{km}), graine (jusqu'à \SI{6371}{km}). Annoter chaque couche avec $v_P$, $v_S$, l'état physique.
\item \textbf{Synthèse.} Rédiger une conclusion de 15 lignes maximum intitulée : « Comment la sismologie a révélé la structure de la Terre ».
\end{enumerate}

\courssub{Ressources à mobiliser}

\begin{aretenir}[Boîte à outils de l'activité]
\begin{itemize}
\item Les ondes \cle{P} (compression) se propagent dans les solides \emph{et} les liquides ; les ondes \cle{S} (cisaillement) ne se propagent \emph{que dans les solides}.
\item À une discontinuité, les ondes subissent \cle{réfraction} (loi de Snell-Descartes : $\dfrac{\sin i_1}{v_1} = \dfrac{\sin i_2}{v_2}$) et \cle{réflexion} (phases PcP, ScS).
\item Relation fondamentale : $t = \dfrac{d}{v}$, soit pour un trajet aller-retour vertical : $t = \dfrac{2h}{v}$.
\item Une chute brutale de $v_P$ à une interface indique un changement d'état (solide $\to$ liquide) ; un arrêt total des S indique un milieu liquide.
\item Échelle : $\text{distance sur le schéma} = \text{distance réelle} \times \text{échelle}$.
\end{itemize}
\end{aretenir}

\courssub{Démarche et corrigé commenté}

\begin{exoresolu}[Radiographie du Globe — corrigé détaillé]
\textbf{Énoncé :} Exploiter le sismogramme du séisme du Chili (2010) enregistré à Yaoundé ($\Delta = 85^\circ$) pour déterminer la profondeur de la discontinuité manteau--noyau, démontrer la liquidité du noyau externe et construire le modèle en couches de la Terre à l'échelle.

\tcblower

\textbf{Étape 1 — Mesures et résidus.}

Conversion en secondes : 
12 min 10 s $= 730$ s ; 
16 min 30 s $= 990$ s ; 
13 min 05 s $= 785$ s ; 
22 min 00 s $= 1320$ s ; 
23 min 50 s $= 1430$ s.

\begin{center}
\begin{tabularx}{\linewidth}{|l|X|X|X|}
\hline
\rowcolor{popblueL} \textbf{Phase} & $t_{\text{obs}}$ (s) & $t_{\text{AK135}}$ (s) & Résidu $r$ (s) \\
\hline
P & 730 & 720 & $+10$ \\
\hline
PP & 990 & 980 & $+10$ \\
\hline
PcP & 785 & 780 & $+5$ \\
\hline
S & 1320 & 1310 & $+10$ \\
\hline
ScS & 1430 & 1420 & $+10$ \\
\hline
\end{tabularx}
\end{center}

Les résidus sont tous positifs et faibles (inférieurs à $1\,\%$ du temps de parcours) : le modèle AK135, construit comme une moyenne mondiale unidimensionnelle, décrit bien la structure moyenne de la Terre ; les petits écarts systématiques traduisent des hétérogénéités locales (anisotropie, variations de température/composition) le long du trajet spécifique Chili--Yaoundé.

\textbf{Étape 2 — Profondeur de la discontinuité manteau--noyau (Gutenberg).}

Dans le modèle simplifié (trajet vertical, vitesse constante $v_P = \SI{13,7}{km.s^{-1}}$ dans le manteau), la phase PcP parcourt un aller-retour $2h$ :
\[ h = \frac{v_P \times t_{\text{PcP}}}{2} = \frac{13{,}7 \times 785}{2} \approx \SI{5377}{km}. \]
Cette valeur surestime fortement la profondeur réelle (\SI{2891}{km}). L'erreur vient de deux approximations majeures : (1) le trajet réel est courbe (loi de Snell-Descartes), pas vertical ; (2) la vitesse $v_P$ n'est pas constante : elle augmente avec la profondeur dans le manteau (de $\sim \SI{8}{km.s^{-1}}$ à $\sim \SI{13,7}{km.s^{-1}}$), donc la vitesse moyenne effective sur le trajet profond est bien plus faible que la valeur maximale du manteau inférieur. En utilisant une vitesse moyenne effective réaliste pour le trajet PcP ($v_{\text{moy}} \approx \SI{7,4}{km.s^{-1}}$, valeur cohérente avec les modèles de référence), on obtient :
\[ h = \frac{7{,}4 \times 785}{2} \approx \SI{2905}{km} \approx \SI{2891}{km}. \]
La discontinuité de Gutenberg (manteau / noyau externe) est donc située vers \SI{2900}{km} de profondeur. C'est exactement par l'inversion des courbes de temps de parcours $t(\Delta)$ que les sismologues (Gutenberg, 1913) ont déterminé cette profondeur.

\textbf{Étape 3 — État physique du noyau externe.}

\begin{popfigure}
\begin{tikzpicture}[scale=0.55, >=Stealth]
% Couleurs couches
\fill[popblue!10] (0,0) circle (6.37); % Terre entière
\fill[poporange!15] (0,0) circle (3.48); % Manteau (rayon 3480 km -> 3.48 cm)
\fill[popgreen!20] (0,0) circle (1.22); % Noyau externe (rayon 1221 km -> 1.22 cm)
\fill[poppurple!30] (0,0) circle (0.7); % Graine (rayon ~1221 km? Non 1221 est ICB. Graine rayon 1221. Attendre: ICB à 5150 km profondeur -> Rayon 1221 km. CMB à 2891 km -> Rayon 3480 km.)
% Correction rayons:
% R_T = 6371 -> 6.37
% CMB depth 2891 -> R = 3480 -> 3.48
% ICB depth 5150 -> R = 1221 -> 1.22
% Donc:
\fill[popblue!10] (0,0) circle (6.37); % Croûte+Manteau+Noyau (fond)
\fill[poporange!15] (0,0) circle (3.48); % Manteau
\fill[popgreen!20] (0,0) circle (1.22); % Noyau externe
\fill[poppurple!30] (0,0) circle (0.7); % Graine (approx 0.7cm pour 1221km? Non 1221km = 1.22cm. La graine EST le noyau interne. Rayon 1221km. Donc cercle 1.22cm. Le noyau externe est l'anneau entre 3.48 et 1.22.)
% Re-dessiner proprement:
% 1. Terre (6.37)
% 2. Manteau (3.48)
% 3. Noyau externe (1.22)
% 4. Graine (solide) - rayon 1.22? Non, graine rayon 1221km = 1.22cm. Noyau externe de 3.48 à 1.22.
% Donc:
\fill[popblue!10] (0,0) circle (6.37); % Tout
\fill[poporange!15] (0,0) circle (3.48); % Manteau (écrase le centre)
\fill[popgreen!20] (0,0) circle (1.22); % Noyau externe (écrase centre)
\fill[poppurple!30] (0,0) circle (0.7); % Graine? Non, graine rayon 1.22.
% Ok, je refais l'ordre de dessin: du plus grand au plus petit.
\fill[popblue!10] (0,0) circle (6.37); % Rayon Terre
\fill[poporange!15] (0,0) circle (3.48); % Rayon CMB
\fill[popgreen!20] (0,0) circle (1.22); % Rayon ICB (limite noyau externe/graine)
\fill[poppurple!30] (0,0) circle (0.7); % Juste un point central pour la graine? Non, graine rayon 1.22.
% La graine EST la sphère de rayon 1.22. Le noyau externe est l'enveloppe entre 3.48 et 1.22.
% Donc:
\fill[popblue!10] (0,0) circle (6.37); % Surface
\fill[poporange!15] (0,0) circle (3.48); % Manteau
\fill[popgreen!20] (0,0) circle (1.22); % Noyau externe + Graine (liquide+solide)
% Pour distinguer noyau externe et graine:
\fill[poppurple!30] (0,0) circle (1.22); % Graine (solide) - rayon 1.22
\fill[popgreen!20] (0,0) circle (3.48); % Noyau externe (liquide) - rayon 3.48 (écrase graine)
% Non, TikZ dessine par dessus. Ordre: Terre -> Manteau -> Noyau Externe -> Graine.
\fill[popblue!10] (0,0) circle (6.37);
\fill[poporange!15] (0,0) circle (3.48);
\fill[popgreen!20] (0,0) circle (1.22); % Noyau externe dessiné jusqu'au centre
\fill[poppurple!30] (0,0) circle (0.7); % Graine (petit cercle central pour visuel? Non, graine rayon 1.22).
% Problème: à l'échelle 0.55, 1.22cm est petit. 0.7cm est trop petit.
% Je vais dessiner les cercles pleins pour les limites et hachurer.
% Simplification: Dessiner les 3 cercles limites.
\draw[thick] (0,0) circle (6.37); % Surface
\draw[thick] (0,0) circle (3.48); % CMB
\draw[thick] (0,0) circle (1.22); % ICB
% Légendes couches
\node[popdark, font=\small] at (0, 5.5) {\textbf{Manteau (solide)}};
\node[popdark, font=\small] at (0, 2.5) {\textbf{Noyau externe (liquide)}};
\node[popdark, font=\small] at (0, 0) {\textbf{Graine (solide)}};
% Source et Récepteur
\coordinate (Src) at (90:6.37);
\coordinate (Rec) at (5:6.37); % Delta = 85 deg
\fill[popred] (Src) circle (0.08) node[above left] {Source (Chili)};
\fill[popred] (Rec) circle (0.08) node[above right] {Yaoundé};
% Rayons P et S directs (courbes dans le manteau)
% Approximation: arcs de cercle ou lignes brisées. Utilisons des courbes paramétriques simplifiées.
% P direct: tourne dans le manteau. Point le plus profond (tournant) pour delta=85 est dans le manteau inférieur.
% PcP: réflexion sur CMB (r=3.48). Point de réflexion à mi-chemin en angle? ~42.5 deg de la source.
% Angle source = 90. Angle rec = 5. Milieu = 47.5? Non, (90+5)/2 = 47.5. Mais sens horaire.
% Point réflexion PcP: angle 47.5 deg, rayon 3.48.
\coordinate (RefPcP) at (47.5:3.48);
\draw[popblue, thick, ->] (Src) .. controls (60:4.5) .. (RefPcP) .. controls (35:4.5) .. (Rec);
\node[popblue, font=\small] at (50:5.5) {PcP};
% P direct: passe plus profond que PcP? Non, P direct tourne au-dessus de la CMB. PcP touche la CMB.
% Pour delta=85, P direct tourne dans le manteau inférieur.
\draw[popblue, dashed, thick] (Src) .. controls (55:3.8) .. (Rec);
\node[popblue, font=\small] at (50:4.2) {P};
% S direct: similaire à P mais plus lent, tourne moins profond? Non, même géométrie (loi de Snell), mais v_S < v_P.
\draw[poporange, dashed, thick] (Src) .. controls (50:3.9) .. (Rec);
\node[poporange, font=\small] at (45:4.0) {S};
% ScS: réflexion sur CMB.
\coordinate (RefScS) at (47.5:3.48); % Même point géométrique
\draw[poporange, thick, ->] (Src) .. controls (60:4.5) .. (RefScS) .. controls (35:4.5) .. (Rec);
\node[poporange, font=\small] at (50:5.0) {ScS};
% Zone d'ombre S (>103°) et P (103-143°) - Indication textuelle
\node[align=center, font=\small, popdark] at (-4, -7) {
Zone d'ombre S : $\Delta > 103^\circ$ \\
Zone d'ombre P : $103^\circ < \Delta < 143^\circ$
};
% Flèches zone d'ombre
\draw[popred, <->] (103:7) arc (103:143:7) node[midway, below, sloped, font=\tiny] {Zone d'ombre P};
\draw[popred, <-] (103:6.8) -- ++(180:1) node[left, font=\tiny] {$\Delta=103^\circ$};
\draw[popred, ->] (143:6.8) -- ++(0:1) node[right, font=\tiny] {$\Delta=143^\circ$};
\draw[popred, <->] (103:6.5) arc (103:180:6.5) node[midway, left, sloped, font=\tiny] {Zone d'ombre S};
\end{tikzpicture}
\legende{Schéma des trajets des ondes pour $\Delta = 85^\circ$ (traits pleins : phases réfléchies PcP, ScS ; traits pointillés : phases directes P, S) et illustration des zones d'ombre au-delà de $103^\circ$. Le noyau externe (vert) arrête les ondes S et réfracte les ondes P vers le centre.}
\end{popfigure}

L'argumentation structurée :
\begin{enumerate}
\item Les ondes S (cisaillement) ne se propagent \emph{pas} dans les liquides (module de cisaillement nul).
\item L'observation d'une zone d'ombre totale des ondes S directes pour $\Delta > 103^\circ$ prouve qu'elles sont bloquées par une région liquide au centre : le \textbf{noyau externe est liquide}.
\item La chute brutale de $v_P$ de $\SI{13,7}{km.s^{-1}}$ (manteau) à $\sim \SI{8}{km.s^{-1}}$ (noyau externe) à la discontinuité de Gutenberg réfracte les ondes P vers la normale (loi de Snell-Descartes), les courbant fortement vers le centre.
\item Cette réfraction crée une \textbf{zone d'ombre des ondes P directes entre $103^\circ$ et $143^\circ$} : aucun rayon P direct ne peut émerger dans cet intervalle angulaire.
\item L'existence de la phase réfléchie \textbf{PcP} (observée à $85^\circ$) confirme une interface nette et réflectrice à \SI{2891}{km}.
\item La détection ultérieure de phases traversant le centre (\textbf{PKIKP}) prouve que la \textbf{graine (noyau interne) est solide} (propagation d'ondes S dans la graine : PKJKP).
\end{enumerate}

\textbf{Étape 4 — Modèle à l'échelle $1/100\,000\,000$ (1 cm = \SI{1000}{km}).}

\begin{center}
\begin{tabularx}{\linewidth}{|l|X|X|X|}
\hline
\rowcolor{popblueL} \textbf{Couche} & \textbf{Profondeur (km)} & \textbf{Rayon sur schéma (cm)} & \textbf{Annotations clés} \\
\hline
Croûte continentale & 0 -- 35 & 6,37 -- 6,335 (trait fin) & $v_P \approx \SI{6}{km.s^{-1}}$, $v_S \approx \SI{3,5}{km.s^{-1}}$, solide \\
\hline
Manteau & 35 -- 2891 & 6,335 -- 3,48 & $v_P$: \num{8} $\to$ \SI{13,7}{km.s^{-1}}, $v_S$ présentes, solide (rhéologie ductile) \\
\hline
Noyau externe & 2891 -- 5150 & 3,48 -- 1,22 & $v_P \approx \SI{8}{km.s^{-1}}$, $v_S = 0$, \textbf{liquide} (Fe-Ni fondu) \\
\hline
Graine (noyau interne) & 5150 -- 6371 & 1,22 -- 0 & $v_P \approx \SI{11}{km.s^{-1}}$, $v_S$ présentes, \textbf{solide} (Fe-Ni cristallisé) \\
\hline
\end{tabularx}
\end{center}
\textit{Nota :} La croûte (\SI{35}{km} = \SI{0,035}{cm}) est quasi invisible à cette échelle ; on la figure par un trait fin à la surface. Le dessin montre trois cercles concentriques de rayons \SI{6,37}{cm}, \SI{3,48}{cm} et \SI{1,22}{cm}.

\textbf{Étape 5 — Synthèse (Conclusion type Probatoire).}

\textbf{Comment la sismologie a révélé la structure de la Terre}

L'analyse des sismogrammes montre deux types d'ondes de volume : les ondes P (compression) et les ondes S (cisaillement). Leurs temps de parcours en fonction de la distance angulaire $\Delta$ (courbes $t(\Delta)$) présentent des ruptures de pente (discontinuités) et des zones d'ombre. La première discontinuité (Mohorovičić, \SI{\sim 35}{km}) marque la base de la croûte (saut de $v_P$ de 6 à \SI{8}{km.s^{-1}}). La deuxième (Gutenberg, \SI{2891}{km}) est identifiée grâce aux phases réfléchies PcP/ScS et à la zone d'ombre des ondes S ($\Delta > 103^\circ$) : l'arrêt des S prouve un milieu liquide (noyau externe), confirmé par la chute de $v_P$ à \SI{8}{km.s^{-1}} et la zone d'ombre P ($103^\circ$--$143^\circ$). La troisième (Lehmann, \SI{5150}{km}) est révélée par les phases PKIKP traversant le centre : la graine est solide. La sismologie, seule « radiographie » du globe accessible, a ainsi construit le modèle de référence (PREM, AK135) à la base de toute la géodynamique moderne.
\end{exoresolu}

\begin{attention}[Pièges fréquents au Probatoire]
\begin{itemize}
\item Oublier le facteur 2 dans $h = \dfrac{v \times t}{2}$ : le trajet PcP est un \emph{aller-retour} (descente + montée).
\item Confondre distance angulaire $\Delta$ (en degrés) et distance épicentrale en kilomètres ($d = \frac{\pi}{180} R_T \Delta$).
\item Écrire que « les ondes S sont absorbées par le noyau » : \textbf{erreur majeure}. Elles ne peuvent pas exister dans un liquide (module de rigidité $\mu = 0$), elles sont converties ou réfléchies, pas absorbées.
\item À l'échelle $1/100\,000\,000$, ne pas convertir correctement : 1 cm sur le schéma $= 10^8$ cm réels $= \SI{1000}{km}$. Vérifier : $R_T = \SI{6371}{km} \to \SI{6,37}{cm}$.
\item Ne pas distinguer \emph{vitesse instantanée} (locale) et \emph{vitesse moyenne effective} sur un trajet courbe pour le calcul de profondeur.
\end{itemize}
\end{attention}

\courssub{Bilan de l'activité}

Cette « radiographie du Globe » a permis de mobiliser l'ensemble des savoir-faire de la leçon dans une démarche complète d'investigation : mesure sur un document réel (sismogramme IRGM Yaoundé), comparaison critique à un modèle de référence mondial (calcul de résidus AK135), exploitation quantitative d'une phase réfléchie (PcP) pour localiser la discontinuité de Gutenberg, argumentation physique rigoureuse sur l'état liquide du noyau externe (zones d'ombre S et P, loi de Snell-Descartes), et construction d'un modèle en couches à l'échelle. Tu as ainsi reproduit, à ton niveau, le raisonnement historique de Mohorovičić (1909), Gutenberg (1913) et Lehmann (1936) : c'est bien l'analyse des temps de parcours des ondes sismiques — et jamais une observation directe — qui a révélé la structure interne de la Terre. Cette compétence d'exploitation de données expérimentales, de la mesure au modèle, est exactement celle attendue au Probatoire des séries C et TI.

\newpage

\coursec{Méthodes et exercices résolus}

\courssub{Les messagers du profond : ondes P et ondes S}

\begin{experience}[Observation d'un sismogramme : deux types d'arrivées]
\textbf{Contexte.} La station sismologique de Yaoundé (Cameroun) enregistre un séisme lointain. Le sismogramme (enregistrement du mouvement du sol en fonction du temps) montre deux groupes distincts de secousses séparés par un intervalle de calme.
\textbf{Matériel.} Sismogramme numérisé (composante verticale), règle, calculatrice.
\textbf{Protocole.}
\begin{enumerate}
\item Repérer la première arrivée brutale (début de l'enregistrement du signal).
\item Mesurer l'instant $t_1$ de cette première arrivée.
\item Repérer la seconde arrivée, plus ample, souvent plus destructrice.
\item Mesurer l'instant $t_2$ de cette seconde arrivée.
\item Calculer l'écart $\Delta t = t_2 - t_1$.
\end{enumerate}
\textbf{Observations.} Quel que soit le séisme, la première arrivée est toujours la plus rapide. L'écart $\Delta t$ augmente avec la distance entre la station et l'épicentre.
\textbf{Interprétation.} Deux types d'ondes de volume se propagent à l'intérieur de la Terre : les \cle{ondes P} (Primaires, longitudinales, compression-dilatation) et les \cle{ondes S} (Secondaires, transversales, cisaillement). La différence de vitesse explique l'ordre d'arrivée.
\end{experience}

\begin{definition}[Ondes de volume : P et S]
\begin{itemize}
\item \textbf{Onde P (Primaire) :} onde de \cle{compression-dilatation} (longitudinale). Le mouvement des particules est \textbf{parallèle} à la direction de propagation. Elle se propage dans les \textbf{solides et les liquides}. Vitesse $v_P = \sqrt{\dfrac{K + \frac{4}{3}\mu}{\rho}}$ ($K$ module de compressibilité, $\mu$ module de rigidité, $\rho$ densité).
\item \textbf{Onde S (Secondaire) :} onde de \cle{cisaillement} (transversale). Le mouvement des particules est \textbf{perpendiculaire} à la direction de propagation. Elle ne se propage que dans les \textbf{solides} ($\mu > 0$). Vitesse $v_S = \sqrt{\dfrac{\mu}{\rho}}$.
\item \textbf{Conséquence fondamentale :} dans un même milieu, $v_P > v_S$ (généralement $v_P \approx \sqrt{3}\,v_S \approx 1{,}73\,v_S$ pour un solide poissonien, environ $1{,}8$ à $2$ fois plus rapide dans le manteau).
\end{itemize}
\end{definition}

\begin{propriete}[Propriétés comparées des ondes P et S]
\begin{center}
\begin{tabularx}{\linewidth}{|l|X|X|}\hline
\textbf{Propriété} & \textbf{Onde P} & \textbf{Onde S} \\ \hline
Nature & Compression-dilatation (longitudinale) & Cisaillement (transversale) \\
Milieu de propagation & Solides \textbf{et} Liquides & \textbf{Uniquement Solides} \\
Vitesse (manteau supérieur) & $\approx \SI{8}{km/s}$ & $\approx \SI{4,6}{km/s}$ \\
Rapport $v_P/v_S$ & $\approx 1{,}73$ (théorique) à $1{,}8-2{,}0$ (réel) & -- \\
Sensibilité aux fluides & Ralentissement marqué si liquide & \textbf{Arrêt total} si liquide \\ \hline
\end{tabularx}
\end{center}
\textbf{Application au Probatoire :} La disparition des ondes S est \emph{la} preuve directe de la présence d'un milieu liquide sur le trajet.
\end{propriete}

\begin{methode}[Reconnaître et distinguer les ondes P et S sur un sismogramme]
\begin{enumerate}
\item \textbf{Rappeler les propriétés fondamentales} : les ondes \cle{P} (primaires) sont des ondes de \cle{compression-dilatation} (longitudinales) : elles se propagent dans les solides \textbf{et} les liquides, avec une vitesse élevée ($v_P \approx 6$ à \SI{13}{km/s} selon la profondeur). Les ondes \cle{S} (secondaires) sont des ondes de \cle{cisaillement} (transversales) : elles ne se propagent que dans les \textbf{solides}, avec $v_S \approx 0{,}55\,v_P$.
\item \textbf{Exploiter l'ordre d'arrivée} : sur un sismogramme, l'onde arrivée en premier est toujours l'onde P (la plus rapide) ; l'onde S arrive ensuite. L'écart de temps $\Delta t = t_S - t_P$ augmente avec la distance à l'épicentre.
\item \textbf{Exploiter le comportement aux interfaces} : si une onde disparaît brutalement au-delà d'une certaine distance angulaire (zone d'ombre), c'est qu'elle a rencontré un milieu où elle ne peut pas se propager (onde S face à un milieu liquide) ou qu'elle a été fortement réfractée (onde P).
\item \textbf{Conclure sur l'état physique du milieu} : propagation des S $\Rightarrow$ milieu solide ; arrêt des S et ralentissement des P $\Rightarrow$ milieu liquide.
\end{enumerate}
\textbf{Réflexe Probatoire} : toute conclusion sur l'état d'une couche doit citer explicitement le comportement observé des ondes (disparition, ralentissement, réflexion) et la propriété des ondes qui l'explique.
\end{methode}

\begin{exoresolu}[Analyse d'un sismogramme enregistré à Yaoundé]
Un séisme se produit dans la fosse océanique de l'Atlantique. La station sismologique de Yaoundé, située à une distance angulaire de $65°$ de l'épicentre, enregistre deux arrivées successives : une première secousse à $t_1 = 10~\text{min}~12~\text{s}$, puis une seconde à $t_2 = 18~\text{min}~30~\text{s}$.
\begin{enumerate}
\item Identifier les ondes correspondant à chaque arrivée. Justifier.
\item Une station située à $150°$ de distance angulaire n'enregistre aucune de ces deux ondes directes. Proposer une explication.
\end{enumerate}
\tcblower
\textbf{Solution détaillée}

\textbf{1. Identification des ondes.}

Les ondes sismiques de volume émises par le foyer sont de deux types : les ondes P (longitudinales, compression-dilatation) et les ondes S (transversales, cisaillement). Or, dans un même milieu, la vitesse des ondes P est toujours supérieure à celle des ondes S ($v_P \approx 1{,}8\,v_S$ en moyenne dans le manteau). Pour un même trajet, l'onde la plus rapide arrive la première.

La première arrivée ($t_1 = 10~\text{min}~12~\text{s}$) correspond donc aux \textbf{ondes P}, et la seconde ($t_2 = 18~\text{min}~30~\text{s}$) aux \textbf{ondes S}. L'écart $\Delta t = t_2 - t_1 = 8~\text{min}~18~\text{s}$ est cohérent avec une station éloignée de $65°$.

\textbf{2. Absence d'ondes directes à $150°$.}

À $150°$ de distance angulaire, la station se trouve dans la \cle{zone d'ombre}. Les rayons sismiques qui pénètrent dans le noyau externe subissent une \textbf{réfraction} vers le bas à la discontinuité de Gutenberg (manteau--noyau, vers \SI{2900}{km} de profondeur) car la vitesse des ondes P chute brutalement dans le noyau externe liquide : les ondes P directes ne peuvent alors atteindre les stations situées entre environ $103°$ et $143°$ (zone d'ombre des P). De plus, les ondes S, ondes de cisaillement, \textbf{ne se propagent pas dans les liquides} : elles sont totalement arrêtées par le noyau externe liquide, ce qui crée une zone d'ombre des S pour toute distance supérieure à environ $103°$.

\textbf{Conclusion} : la station à $150°$ ne reçoit ni P ni S directes car le noyau externe liquide réfracte les ondes P (zone d'ombre des P) et bloque totalement les ondes S — preuve directe de l'état liquide du noyau externe.
\end{exoresolu}

\courssub{Le voyage des ondes : réfraction, réflexion et loi de Snell-Descartes en géophysique}

\begin{propriete}[Loi de Snell-Descartes en sismologie (Terre sphérique)]
Dans une Terre modélisée par des couches concentriques à symétrie sphérique, le \cle{paramètre de rayon} $p$ est constant le long d'un trajet de rayon :
\[ p = \dfrac{r \sin i}{v(r)} = \text{constante} \]
où $r$ est la distance au centre de la Terre, $i$ l'angle d'incidence (angle entre le rayon et la normale locale, i.e. le rayon terrestre), et $v(r)$ la vitesse locale de l'onde (P ou S).
\begin{itemize}
\item \textbf{Réfraction} : si $v$ augmente avec la profondeur ($dv/dr > 0$), $\sin i$ doit augmenter $\Rightarrow$ le rayon courbe vers le haut (remonte vers la surface).
\item \textbf{Réfraction} : si $v$ diminue avec la profondeur ($dv/dr < 0$), $\sin i$ diminue $\Rightarrow$ le rayon courbe vers le bas (penche vers le centre).
\item \textbf{Réflexion} : à une discontinuité, une partie de l'énergie est réfléchie (angle de réflexion = angle d'incidence) et une partie transmise (réfractée selon la loi de Snell).
\item \textbf{Distance angulaire $\Delta$} : $\Delta = 2 \int_{r_{min}}^{R} \dfrac{p\,v(r)}{r\sqrt{r^2 - p^2 v(r)^2}}\,dr$ (intégrale sur la branche descendante).
\end{itemize}
\textbf{Note Probatoire} : la forme intégrale n'est pas exigible, mais le principe « vitesse croissante $\to$ rayon courbe vers le haut » et « vitesse décroissante $\to$ rayon courbe vers le bas » est \textbf{exigible}. La loi des sinus $\frac{\sin i_1}{v_1} = \frac{\sin i_2}{v_2}$ à une interface plane est connue.
\end{propriete}

\begin{definition}[Discontinuité sismique et Zone d'ombre]
\begin{itemize}
\item \textbf{Discontinuité sismique :} surface de séparation entre deux milieux où au moins une des vitesses ($v_P$ ou $v_S$) varie brutalement (saut de vitesse). Elle se détecte par une rupture de pente sur la courbe hodochrone $t = f(\Delta)$ ou par l'apparition d'ondes réfléchies (ex. PcP, ScS).
\item \textbf{Zone d'ombre (shadow zone) :} intervalle de distances angulaires où aucune onde directe (P ou S) n'est enregistrée. Elle résulte de la réfraction forte (courbure des rayons) ou de l'arrêt d'une onde (S dans un liquide).
\end{itemize}
\end{definition}

\begin{savaistu}[Histoire des sciences : Les pionniers de la sismologie globale]
\begin{itemize}
\item \textbf{Andrija Mohorovičić (1909)} : découvre la discontinuité croûte-manteau (Moho) en Croatie en observant deux arrivées P distinctes.
\item \textbf{Beno Gutenberg (1913)} : identifie la discontinuité manteau-noyau à \SI{2900}{km} (discontinuité de Gutenberg) grâce à la zone d'ombre des ondes P.
\item \textbf{Inge Lehmann (1936)} : découvre la discontinuité noyau externe-noyau interne (discontinuité de Lehmann) à \SI{5150}{km} en analysant des arrivées P dans la zone d'ombre (ondes PKP, puis PKJKP).
\item \textbf{Modèles de référence :} \cle{PREM} (Preliminary Reference Earth Model, Dziewonski \& Anderson, 1981) et \cle{AK135} (Kennett et al., 1995) sont les modèles 1D standards utilisés aujourd'hui pour la localisation des séismes et l'imagerie tomographique.
\end{itemize}
\end{savaistu}

\courssub{Exploitation des courbes de temps de parcours : localiser les discontinuités}

\begin{methode}[Situer une discontinuité à partir des données sismiques]
\begin{enumerate}
\item \textbf{Organiser les données} : relever pour chaque station la distance angulaire $\Delta$ (en degrés) et les temps d'arrivée $t_P$ et $t_S$.
\item \textbf{Tracer ou exploiter les courbes hodochrones} : $t = f(\Delta)$. Une \textbf{rupture de pente}, un \textbf{décalage brutal} de la courbe ou une \textbf{zone sans données} (zone d'ombre) traduisent l'existence d'une discontinuité.
\item \textbf{Exploiter les ondes réfléchies} : une onde réfléchie sur une discontinuité (ex. PcP, réfléchie sur le noyau) a un temps de parcours $t = \dfrac{2h}{v}$ pour un trajet vertical aller-retour, d'où la profondeur $h = \dfrac{v \times t}{2}$.
\item \textbf{Calculer une vitesse moyenne} si nécessaire : $v = \dfrac{d}{t}$ avec $d$ la longueur du trajet (attention : pour un trajet en profondeur, $d$ n'est pas la distance en surface !).
\item \textbf{Conclure} en nommant la discontinuité (Mohorovičić/Moho vers 5--70 km, Gutenberg vers \SI{2900}{km}, Lehmann vers \SI{5150}{km}) et les couches qu'elle sépare.
\end{enumerate}
\textbf{Formules à retenir} : $v = \dfrac{d}{t}$ ; profondeur d'une interface par réflexion verticale : $h = \dfrac{v\,t}{2}$ ; distance épicentrale par la méthode des trois stations (triangulation).
\end{methode}

\begin{exoresolu}[Localisation de la discontinuité de Gutenberg]
Une onde P émise par un séisme peu profond se réfléchit verticalement sur la surface du noyau (onde notée PcP) et revient à la station située quasiment à la verticale du foyer. Le temps aller-retour mesuré est $t = 16~\text{min}~40~\text{s}$. On admet que la vitesse moyenne des ondes P dans le manteau est $v = \SI{11}{km/s}$.
\begin{enumerate}
\item Calculer la profondeur de la discontinuité manteau--noyau.
\item Nommer cette discontinuité et préciser la nature des milieux qu'elle sépare.
\end{enumerate}
\tcblower
\textbf{Solution détaillée}

\textbf{1. Calcul de la profondeur.}

Le temps mesuré correspond à un trajet \textbf{aller-retour} vertical : descente de l'onde de la surface jusqu'à la discontinuité (distance $h$), réflexion, puis remontée (distance $h$). La distance totale parcourue est donc $d = 2h$.

La relation fondamentale est $v = \dfrac{d}{t}$, soit $d = v \times t$, donc :
\[ h = \dfrac{v \times t}{2} \]

Conversion du temps en secondes :
\[ t = 16~\text{min}~40~\text{s} = 16 \times 60 + 40 = 1000~\text{s} \]

Application numérique :
\begin{align*}
h &= \dfrac{11 \times 1000}{2} \\
h &= \dfrac{11000}{2} \\
h &= 5500~\text{km}
\end{align*}

\textit{Remarque} : la valeur obtenue est surestimée par rapport à la valeur réelle (\SI{2900}{km}) car la vitesse moyenne réelle des P dans l'ensemble du trajet manteau est inférieure à \SI{11}{km/s} (qui est la vitesse du manteau inférieur). Avec $v_{moy} \approx \SI{5,8}{km/s}$ sur tout le trajet, on obtiendrait $h = \dfrac{5{,}8 \times 1000}{2} = 2900$ km, valeur de référence. L'exercice illustre l'importance du choix de la vitesse moyenne.

\textbf{2. Identification.}

Cette discontinuité située vers \SI{2900}{km} de profondeur est la \textbf{discontinuité de Gutenberg}. Elle sépare :
\begin{itemize}
\item le \textbf{manteau inférieur} (solide, silicaté), dans lequel les ondes S se propagent ;
\item le \textbf{noyau externe} (liquide, métallique, fer-nickel), qui bloque les ondes S et ralentit fortement les ondes P.
\end{itemize}

\textbf{Conclusion} : la discontinuité de Gutenberg, mise en évidence par la réflexion des ondes PcP et la zone d'ombre des S, marque la limite manteau--noyau à environ \SI{2900}{km} de profondeur.
\end{exoresolu}

\courssub{Le noyau terrestre : la preuve par les ondes (liquide externe, solide interne)}

\begin{methode}[Déduire l'état physique d'une couche à partir des données sismiques]
\begin{enumerate}
\item \textbf{Identifier l'observation} : disparition d'une onde, chute ou saut de vitesse, réflexion, réfraction, zone d'ombre.
\item \textbf{Associer chaque observation à une cause physique} :
\begin{itemize}
\item les ondes S disparaissent $\Rightarrow$ milieu \textbf{liquide} (un liquide ne transmet pas le cisaillement) ;
\item la vitesse des P chute brutalement $\Rightarrow$ passage à un milieu moins rigide (solide $\to$ liquide) ;
\item la vitesse des P augmente brutalement $\Rightarrow$ passage à un milieu plus rigide ou plus dense (liquide $\to$ solide, ou changement de composition) ;
\item les ondes S réapparaissent $\Rightarrow$ retour à un milieu \textbf{solide}.
\end{itemize}
\item \textbf{Nommer la discontinuité} correspondante et les couches concernées.
\item \textbf{Rédiger la conclusion en chaîne causale} : observation sismique $\to$ propriété des ondes $\to$ état physique de la couche.
\end{enumerate}
\textbf{Piège classique} : le manteau est \textbf{solide} (les S s'y propagent) même s'il est ductile sur de longues durées ; ne jamais écrire « manteau liquide ».
\end{methode}

\begin{exoresolu}[Preuve de la solidité de la graine (noyau interne)]
Le document ci-dessous résume les observations sismiques à travers la Terre :
\begin{itemize}
\item les ondes S disparaissent à toutes les stations situées au-delà de $103°$ de distance angulaire ;
\item les ondes P sont fortement ralenties entre $103°$ et $143°$ (zone d'ombre des P) ;
\item au-delà de $143°$, on enregistre des ondes P ayant traversé le centre de la Terre, et des ondes issues de la conversion P $\to$ S $\to$ P (ondes PKJKP) dont l'existence prouve qu'une onde S a traversé la partie la plus profonde du globe.
\end{itemize}
Déduire de ces observations la structure et l'état physique du noyau terrestre.
\tcblower
\textbf{Solution détaillée}

\textbf{Étape 1 : le noyau externe est liquide.}

Observation : les ondes S disparaissent au-delà de $103°$. Or les ondes S sont des ondes de cisaillement qui ne peuvent se propager que dans les milieux solides (un liquide n'oppose aucune résistance au cisaillement). Leur disparition totale prouve que les rayons sismiques rencontrent, à partir de la discontinuité de Gutenberg (\SI{2900}{km}), un milieu \textbf{liquide} : le \textbf{noyau externe}. La chute de vitesse des P à cette même profondeur (d'où la zone d'ombre des P entre $103°$ et $143°$) confirme la baisse de rigidité du milieu.

\textbf{Étape 2 : le noyau interne (graine) est solide.}

Observation : l'existence d'ondes PKJKP signifie qu'une onde P entrée dans le noyau s'est convertie en onde S (notée J dans la graine) à l'intérieur du noyau, puis reconvertie en P à la sortie. Or une onde S ne peut exister que dans un \textbf{solide}. La partie centrale du noyau, appelée \textbf{graine} ou noyau interne, est donc \textbf{solide}. De plus, la vitesse des P augmente à la discontinuité de Lehmann (vers \SI{5150}{km}), ce qui traduit une rigidité accrue, cohérente avec un état solide (cristallisation du fer sous l'effet des pressions colossales, malgré des températures de l'ordre de \SI{5000}{\celsius}).

\textbf{Conclusion} : le noyau terrestre est double : un \textbf{noyau externe liquide} (de 2900 à 5150 km), qui bloque les S et ralentit les P, entourant une \textbf{graine solide} (de 5150 km au centre, rayon $\approx$ \SI{1220}{km}), dont la solidité est démontrée par la propagation d'ondes S de conversion (PKJKP) et l'accélération des ondes P.
\end{exoresolu}

\courssub{Modèle de référence : Structure en couches de la Terre (Modèle PREM / AK135)}

\begin{aretenir}[Structure interne de la Terre (valeurs de référence PREM/AK135)]
\begin{center}
\begin{tabularx}{\linewidth}{|l|c|c|X|X|}\hline
\textbf{Couche} & \textbf{Profondeur (km)} & \textbf{Épaisseur (km)} & \textbf{État physique} & \textbf{Composition dominante / Remarques} \\ \hline
Croûte continentale & 0 -- 30 à 70 & 30 à 70 & Solide & Granites, gneiss, roches sédimentaires ($v_P \approx 6$ km/s) \\
Croûte océanique & 0 -- 5 à 10 & 5 à 10 & Solide & Basalts, gabbros ($v_P \approx 7$ km/s) \\
\rowcolor{popblueL} \textbf{Discontinuité de Mohorovičić (Moho)} & \textbf{5--70} & -- & \textbf{Saut $v_P$, $v_S$} & \textbf{Limite croûte / manteau} \\ \hline
Manteau supérieur & 30/10 -- 660 & $\approx$ 650 & Solide (ductile) & Péridotites (olivine, pyroxènes) ; Zone de basse vitesse (asthénosphère) vers 100-200 km \\
Zone de transition & 410 -- 660 & 250 & Solide & Transitions de phase olivine $\to$ wadsleyite $\to$ ringwoodite \\
Manteau inférieur & 660 -- 2900 & $\approx$ 2240 & Solide (ductile) & Bridgmanite, ferropericlase ($v_P$ 10--13 km/s) \\
\rowcolor{poporangeL} \textbf{Discontinuité de Gutenberg} & \textbf{2900} & -- & \textbf{Chute $v_P$, arrêt $v_S$} & \textbf{Limite manteau / noyau externe} \\ \hline
Noyau externe & 2900 -- 5150 & $\approx$ 2250 & \textbf{Liquide} & Alliage Fe-Ni (+ éléments légers S, O, Si) ; $v_P \approx 8$ km/s, $v_S = 0$ \\
\rowcolor{popgreenL} \textbf{Discontinuité de Lehmann} & \textbf{5150} & -- & \textbf{Saut $v_P$, apparition $v_S$} & \textbf{Limite noyau externe / graine} \\ \hline
Graine (Noyau interne) & 5150 -- 6371 & $\approx$ 1220 & \textbf{Solide} & Fe-Ni cristallisé (structure hcp) ; $v_P \approx 11$ km/s, $v_S \approx 3{,}5$ km/s \\ \hline
\end{tabularx}
\end{center}
\textbf{Repères} : Rayon terrestre $R = \SI{6371}{km}$. Masse $M = 5,97 \times 10^{24}$ kg. Masse du noyau $\approx 32\%$ de la masse terrestre.
\end{aretenir}

\begin{methode}[Construire et légender un modèle en couches à l'échelle]
\begin{enumerate}
\item \textbf{Choisir l'échelle} : le rayon terrestre est $R = \SI{6371}{km}$. Pour un schéma de rayon 6 cm, l'échelle est 1 cm $\leftrightarrow$ environ \SI{1000}{km} (précisément $6371/6 \approx 1062$ km/cm).
\item \textbf{Placer les discontinuités} à leurs profondeurs caractéristiques :
\begin{itemize}
\item Moho : 5 à 10 km (océanique), 30 à 70 km (continentale) — souvent non représentable à l'échelle ;
\item Gutenberg (manteau--noyau) : \SI{2900}{km} ;
\item Lehmann (noyau externe--graine) : \SI{5150}{km}.
\end{itemize}
\item \textbf{Calculer les rayons des interfaces} depuis le centre : $r = R - \text{profondeur}$. Ex. : surface du noyau : $6371 - 2900 = 3471$ km ; surface de la graine : $6371 - 5150 = 1221$ km.
\item \textbf{Dessiner les couches concentriques} et légender chaque élément par une ligne fine : croûte, manteau (supérieur/inférieur), noyau externe (liquide), graine (solide).
\item \textbf{Annoter l'état physique} et la composition dominante (silicates pour croûte et manteau ; fer-nickel pour le noyau).
\end{enumerate}
\end{methode}

\begin{exoresolu}[Construction d'un modèle à l'échelle]
On veut représenter la structure interne de la Terre par des cercles concentriques de rayon maximal 6 cm.
\begin{enumerate}
\item Calculer les rayons, sur le schéma, des interfaces manteau--noyau et noyau externe--graine.
\item Indiquer pourquoi la croûte ne peut pas être représentée fidèlement à cette échelle.
\end{enumerate}
\tcblower
\textbf{Solution détaillée}

\textbf{1. Calcul des rayons sur le schéma.}

Échelle : 6 cm représentent \SI{6371}{km}, soit 1 cm $\leftrightarrow \dfrac{6371}{6} \approx 1062$ km.

\textit{Interface manteau--noyau (Gutenberg)} : profondeur \SI{2900}{km}, donc rayon réel depuis le centre :
\[ r_1 = 6371 - 2900 = 3471~\text{km} \]
Sur le schéma :
\[ r_1' = \dfrac{3471}{1062} \approx 3{,}3~\text{cm} \]

\textit{Interface noyau externe--graine (Lehmann)} : profondeur \SI{5150}{km}, donc :
\[ r_2 = 6371 - 5150 = 1221~\text{km} \]
Sur le schéma :
\[ r_2' = \dfrac{1221}{1062} \approx 1{,}1~\text{cm} \]

Le schéma comporte donc trois cercles concentriques de rayons 6 cm (surface), 3,3 cm (noyau externe) et 1,1 cm (graine).

\textbf{2. Le problème de la croûte.}

La croûte continentale a une épaisseur moyenne de 30 à 40 km (jusqu'à 70 km sous les chaînes de montagnes), et la croûte océanique seulement 5 à 10 km. À l'échelle du schéma, 40 km représentent :
\[ \dfrac{40}{1062} \approx 0{,}04~\text{cm} = 0{,}4~\text{mm} \]
Cette épaisseur est inférieure à l'épaisseur du trait de crayon : la croûte est \textbf{trop fine pour être représentée à l'échelle} d'un globe de 6 cm de rayon. On la figure donc par le simple contour extérieur, en précisant dans la légende qu'elle n'est pas à l'échelle.

\textbf{Conclusion} : le modèle en couches comporte, du centre vers la surface : la graine solide (rayon 1,1 cm sur le schéma), le noyau externe liquide (jusqu'à 3,3 cm), le manteau solide (jusqu'à 6 cm) et la croûte, non représentable à l'échelle. Cette disproportion illustre la finesse extrême de la croûte, comparable à la peau d'une pomme.
\end{exoresolu}

\courssub{Au-delà du modèle 1D : Hétérogénéités, Anisotropie et Dynamique (Complément Probatoire)}

\begin{methode}[Interpréter un écart au modèle de référence (PREM/AK135)]
\begin{enumerate}
\item \textbf{Comparer} la valeur mesurée (vitesse, temps de parcours) à la valeur prédite par le modèle de référence à la même profondeur.
\item \textbf{Identifier le sens de l'anomalie} : onde plus rapide que prévu $\Rightarrow$ matériau plus froid, plus dense ou plus rigide (ex. plaque plongeante en subduction) ; onde plus lente $\Rightarrow$ matériau plus chaud, partiellement fondu ou moins dense (ex. point chaud, rift).
\item \textbf{Proposer une interprétation géodynamique} : le modèle 1D (PREM, AK135) suppose une Terre en couches homogènes ; les écarts révèlent les \cle{hétérogénéités} latérales liées à la dynamique interne (convection mantellique, panaches, subduction).
\item \textbf{Rester mesuré} : une anomalie de vitesse seule ne prouve pas une composition ; elle renseigne d'abord sur la température et l'état physique.
\end{enumerate}
\end{methode}

\begin{exoresolu}[Anomalie de vitesse sous la dorsale océanique]
Sous une dorsale océanique, à \SI{100}{km} de profondeur, les ondes P se propagent à \SI{7,6}{km/s} alors que le modèle PREM prédit \SI{8,1}{km/s} à cette profondeur. Interpréter cette anomalie.
\tcblower
\textbf{Solution détaillée}

\textbf{Étape 1 : quantifier l'anomalie.}
\[ \Delta v = 8{,}1 - 7{,}6 = 0{,}5~\text{km/s}, \quad \text{soit} \quad \dfrac{0{,}5}{8{,}1} \approx 6\% \text{ de ralentissement.} \]

\textbf{Étape 2 : interpréter physiquement.} La vitesse des ondes sismiques diminue quand la rigidité du milieu diminue, c'est-à-dire quand la température augmente ou quand une fraction du matériau est fondue. Un ralentissement de 6 \% sous la dorsale traduit donc un manteau \textbf{anormalement chaud}, avec une \textbf{fusion partielle} (quelques pourcents de liquide silicaté).

\textbf{Étape 3 : interprétation géodynamique.} Sous les dorsales, le manteau chaud remonte par convection (divergence des plaques) ; la décompression provoque la fusion partielle des péridotites, produisant le magma basaltique qui construit la croûte océanique. L'anomalie de vitesse cartographie donc directement la \textbf{zone d'accrétion océanique}.

\textbf{Conclusion} : l'écart au modèle PREM révèle une hétérogénéité thermique et pétrologique du manteau : sous la dorsale, le manteau chaud partiellement fondu ralentit les ondes P, confirmant que le modèle 1D n'est qu'une moyenne d'une Terre dynamique et hétérogène.
\end{exoresolu}

\begin{savaistu}[Tomographie sismique et ligne volcanique du Cameroun]
La \cle{tomographie sismique} (inversion des temps de parcours de milliers de séismes enregistrés par des réseaux mondiaux et locaux) permet d'imager la structure 3D du manteau. Sous le Cameroun, la \textbf{ligne volcanique du Cameroun} (Mont Cameroun, Manengouba, Bamboutos, etc.) est associée à une anomalie de basse vitesse dans le manteau supérieur, signature d'un \textbf{panache mantellique} (point chaud) ou d'une remontée de manteau chaud liée à la fracturation lithosphérique. Le lac Nyos (catastrophe de 1986) est un maar volcanique sur cette ligne ; la surveillance sismique (réseau du Centre de Recherche en Géophysique de Dschang) utilise ces principes pour monitorer l'activité magmatique et les risques de dégazage.
\end{savaistu}

\courssub{Méthode 3 : Localiser un épicentre par triangulation (méthode des trois stations)}

\begin{methode}[Déterminer la position d'un épicentre]
\begin{enumerate}
\item \textbf{Mesurer} sur chaque sismogramme l'écart $\Delta t = t_S - t_P$ entre les arrivées des ondes S et P.
\item \textbf{Convertir} cet écart en distance épicentrale $d$, soit à l'aide d'un abaque fourni (courbe $d = f(\Delta t)$), soit par le calcul : $d = \dfrac{v_P \times v_S}{v_P - v_S} \times \Delta t$.
\item \textbf{Tracer} pour chaque station un cercle centré sur la station, de rayon $d$ (à l'échelle de la carte).
\item \textbf{Lire} l'intersection des trois cercles : c'est l'épicentre. Deux stations donnent deux points possibles ; il faut \textbf{trois stations minimum} pour une localisation unique.
\end{enumerate}
\textbf{Réflexe} : vérifier la cohérence des unités (km, s) et justifier pourquoi trois stations sont nécessaires.
\end{methode}

\begin{exoresolu}[Séisme dans la région de la faille de Foumban]
Trois stations sismiques enregistrent un séisme en pays bamiléké. Les écarts mesurés sont : $\Delta t_A = 40$ s (station A), $\Delta t_B = 25$ s (station B), $\Delta t_C = 60$ s (station C). On donne $v_P = \SI{6}{km/s}$ et $v_S = \SI{3,5}{km/s}$ dans la croûte.
\begin{enumerate}
\item Calculer la distance épicentrale pour chaque station.
\item Expliquer comment on obtient la position de l'épicentre et pourquoi trois stations sont nécessaires.
\end{enumerate}
\tcblower
\textbf{Solution détaillée}

\textbf{1. Calcul des distances épicentrales.}

Soit $d$ la distance foyer--station (on suppose le foyer proche de la surface). Les temps de parcours sont :
\[ t_P = \dfrac{d}{v_P} \qquad \text{et} \qquad t_S = \dfrac{d}{v_S} \]
L'écart mesuré s'écrit :
\[ \Delta t = t_S - t_P = \dfrac{d}{v_S} - \dfrac{d}{v_P} = d \times \dfrac{v_P - v_S}{v_P \times v_S} \]
D'où la formule de la distance épicentrale :
\[ d = \dfrac{v_P \times v_S}{v_P - v_S} \times \Delta t \]

Calcul du coefficient :
\[ \dfrac{v_P \times v_S}{v_P - v_S} = \dfrac{6 \times 3{,}5}{6 - 3{,}5} = \dfrac{21}{2{,}5} = 8{,}4~\text{km/s} \]

Application numérique pour chaque station :
\begin{align*}
d_A &= 8{,}4 \times 40 = 336~\text{km} \\
d_B &= 8{,}4 \times 25 = 210~\text{km} \\
d_C &= 8{,}4 \times 60 = 504~\text{km}
\end{align*}

\textbf{2. Localisation de l'épicentre.}

Sur une carte à l'échelle, on trace trois cercles : centrés sur A, B et C, de rayons respectifs 336 km, 210 km et 504 km. Le point d'intersection commun des trois cercles est l'\textbf{épicentre}.

\textit{Justification du nombre de stations} : avec une seule station, l'épicentre peut se trouver n'importe où sur un cercle (infinité de positions). Avec deux stations, les deux cercles se coupent en \textbf{deux points} : ambiguïté. Seule une troisième station lève l'ambiguïté : l'intersection des trois cercles est unique.

\textbf{Conclusion} : l'épicentre du séisme se situe au point d'intersection des cercles de rayon 336 km (A), 210 km (B) et 504 km (C) ; la station B étant la plus proche, c'est elle qui a enregistré le plus faible écart S--P.
\end{exoresolu}

\courssub{Synthèse et Vigilance Probatoire}

\begin{attention}[Les 5 erreurs qui coûtent des points au Probatoire]
\begin{enumerate}
\item \textbf{Oublier le facteur 2 dans les calculs de réflexion} : le temps mesuré pour une onde réfléchie (PcP, ScS) correspond à un \textbf{aller-retour}. Écrire $h = v \times t$ au lieu de $h = \dfrac{v \times t}{2}$ double la profondeur et fausse tout l'exercice. Toujours préciser le trajet dans la rédaction.
\item \textbf{Confondre distance épicentrale et profondeur du trajet} : dans $v = d/t$, $d$ est la longueur du \textbf{trajet de l'onde} (à travers la Terre), pas la distance mesurée en surface entre l'épicentre et la station. Pour les grandes distances angulaires, le trajet est une corde ou un arc dans le globe.
\item \textbf{Écrire « le manteau est liquide »} : c'est faux. Les ondes S traversent le manteau, il est donc \textbf{solide} (ductile sur de longues échelles de temps, mais solide vis-à-vis des ondes sismiques). Seul le \textbf{noyau externe} est liquide. Confondre manteau et noyau externe est l'erreur la plus sanctionnée.
\item \textbf{Conclure sans citer la propriété des ondes} : écrire « les S disparaissent donc le noyau est liquide » sans rappeler que \emph{les ondes S sont des ondes de cisaillement qui ne se propagent que dans les solides} fait perdre le point de justification. Toute conclusion doit suivre la chaîne : observation $\to$ propriété des ondes $\to$ état du milieu.
\item \textbf{Négliger les unités et les conversions} : mélanger minutes et secondes (convertir : $16~\text{min}~40~\text{s} = 1000$ s), oublier de convertir les degrés de distance angulaire en kilomètres quand c'est demandé ($1° \approx 111$ km en surface), ou donner une profondeur en mètres au lieu de kilomètres. Toujours encadrer le résultat avec son unité et vérifier son ordre de grandeur (une profondeur terrestre ne dépasse jamais \SI{6371}{km}).
\end{enumerate}
\end{attention}

\begin{exercice}[Entraînement Probatoire : Structure du globe et données sismiques]
\begin{enumerate}
\item \textbf{Doc. 1 :} Un sismogramme enregistré à une distance angulaire de $120°$ montre une arrivée P à $t_P = 15$ min 20 s. Aucune onde S directe n'est enregistrée. On observe cependant une phase PKP (P traversant le noyau) arrivant à $t_{PKP} = 18$ min 10 s.
\begin{itemize}
\item[a)] Expliquer l'absence d'onde S directe.
\item[b)] Que prouve l'existence de la phase PKP sur la structure du noyau ?
\end{itemize}
\item \textbf{Doc. 2 :} Une onde ScS (S réfléchie sur le noyau) est enregistrée à la verticale du foyer. Le temps aller-retour est $t = 52$ min. La vitesse moyenne des ondes S dans le manteau est $v_S = \SI{6,2}{km/s}$.
\begin{itemize}
\item[a)] Calculer la profondeur de la discontinuité de Gutenberg.
\item[b)] Comparer à la valeur admise (\SI{2900}{km}) et commenter.
\end{itemize}
\item \textbf{Doc. 3 :} Schéma simplifié d'un rayon P traversant la Terre (source en surface, récepteur en surface). Le rayon pénètre dans le noyau externe.
\begin{itemize}
\item[a)] Dessiner qualitativement la courbure du rayon dans le manteau (vitesse croissante avec la profondeur) et dans le noyau externe (vitesse décroissante brutale à l'entrée).
\item[b)] Justifier la courbure par la loi de Snell-Descartes (paramètre de rayon $p = cste$).
\end{itemize}
\end{enumerate}
\end{exercice}

\begin{corrige}[Exercice NCT09]
\textbf{1. Doc. 1}
\begin{itemize}
\item[a)] Distance $120° > 103°$. Les rayons S directs pénètrent dans le noyau externe. Or les ondes S sont des ondes de cisaillement qui \textbf{ne se propagent pas dans les liquides}. Le noyau externe étant liquide, il bloque totalement les ondes S directes $\Rightarrow$ zone d'ombre des S pour $\Delta > 103°$.
\item[b)] L'onde PKP traverse le noyau (K = noyau en notation allemande \emph{Kern}). Son existence prouve que les ondes P traversent le noyau. Le fait qu'elle arrive après la P directe (qui a été réfractée) et qu'on observe des branches PKP$_{ab}$, PKP$_{bc}$, PKP$_{df}$ (diffractées) confirme la structure : manteau solide $\to$ noyau externe liquide (ralentissement P) $\to$ graine solide (accélération P). La réfraction dans le noyau liquide crée la zone d'ombre des P ($103°-143°$).
\end{itemize}

\textbf{2. Doc. 2}
\begin{itemize}
\item[a)] Trajet aller-retour vertical : $d = 2h$. $t = 52$ min $= 3120$ s. $v_S = 6,2$ km/s.
\[ h = \dfrac{v_S \times t}{2} = \dfrac{6{,}2 \times 3120}{2} = \dfrac{19344}{2} = 9672~\text{km} \]
\item[b)] Valeur trouvée (9672 km) $\gg$ valeur réelle (2900 km). L'erreur vient de la vitesse utilisée : $v_S = 6,2$ km/s est une vitesse typique du \textbf{manteau inférieur} (profond). La vitesse \textbf{moyenne} sur tout le trajet manteau (supérieur + transition + inférieur) est bien plus faible ($\approx 4,5$ à $5$ km/s). Avec $v_{S,moy} \approx 4,6$ km/s : $h = \frac{4,6 \times 3120}{2} \approx 7176$ km (toujours surestimé car le trajet n'est pas parfaitement vertical pour un séisme peu profond, mais l'ordre de grandeur se rapproche). \textbf{Leçon :} il faut utiliser une vitesse \emph{moyenne pondérée par le trajet}, pas une vitesse locale.
\end{itemize}

\textbf{3. Doc. 3}
\begin{itemize}
\item[a)] Schéma attendu : Rayon partant de la surface, courbant \textbf{vers le haut} (concave vers le haut) dans le manteau car $v_P$ augmente avec la profondeur ($dv/dr > 0 \Rightarrow di/dr > 0$). À l'entrée du noyau externe (Gutenberg), $v_P$ chute brutalement ($\approx 13 \to 8$ km/s). Le rayon se réfracte \textbf{vers la normale} (courbe \textbf{vers le bas}, convexe vers le haut) car $v$ diminue. Il traverse le noyau, se réfracte à nouveau à la sortie (accélération, courbe vers le haut), et remonte vers la station.
\item[b)] Loi de Snell-Descartes sphérique : $p = \frac{r \sin i}{v} = cste$.
Dans le manteau : $v \uparrow \Rightarrow \sin i \uparrow \Rightarrow i \uparrow$ (le rayon s'éloigne de la normale, courbe vers le haut).
Dans le noyau externe : $v \downarrow$ brutalement $\Rightarrow \sin i \downarrow \Rightarrow i \downarrow$ (le rayon se rapproche de la normale, courbe vers le bas).
C'est cette forte courbure vers le bas dans le noyau qui crée la zone d'ombre des P ($103°-143°$).
\end{itemize}
\end{corrige}

\newpage

\coursec{Exercices}

\courssub{Je vérifie mes connaissances}

\begin{exercice}[QCM : Propriétés des ondes P et S]
Parmi les affirmations suivantes, laquelle est \textbf{fausse} concernant les ondes sismiques de volume ?
\begin{enumerate}
    \item Les ondes P (primaires) sont des ondes de compression-dilatation ; leur vitesse de propagation dépend du module d'incompressibilité et de la densité du milieu.
    \item Les ondes S (secondaires) sont des ondes de cisaillement ; elles ne se propagent pas dans les fluides (liquides et gaz) car le module de cisaillement y est nul.
    \item La vitesse des ondes P est toujours supérieure à celle des ondes S dans un même milieu solide isotrope ($v_P > v_S$).
    \item Au passage d'une discontinuité de premier ordre (saut de densité et de modules élastiques), les ondes P et S subissent uniquement une réflexion, sans réfraction.
\end{enumerate}
\end{exercice}

\begin{exercice}[Vrai ou Faux : Comportement aux discontinuités]
Pour chaque affirmation, répondre \textbf{Vrai} ou \textbf{Faux} et \textbf{justifier} brièvement en invoquant la loi de Snell-Descartes ou les conditions de continuité des déplacements/contraintes.
\begin{enumerate}
    \item Lorsqu'une onde P incidente arrive sur la discontinuité Mohorovičić (Moho), elle génère uniquement une onde P réfléchie et une onde P réfractée.
    \item Le paramètre de rayon $p = \frac{\sin i}{v}$ est constant le long du trajet d'une onde sismique dans un milieu sphériquement stratifié.
    \item L'existence d'une \emph{zone d'ombre} pour les ondes S entre $103^\circ$ et $180^\circ$ de distance épicentrale est la preuve directe que le noyau externe est liquide.
    \item La discontinuité de Lehmann (frontière noyau externe / noyau interne) se manifeste par une augmentation brutale de la vitesse des ondes P ($v_P$) et la réapparition des ondes S.
\end{enumerate}
\end{exercice}

\begin{exercice}[Définitions et notions clés]
Donner une définition précise et rigoureuse des termes suivants :
\begin{enumerate}
    \item \textbf{Temps de parcours} (ou \emph{travel time}) d'une phase sismique.
    \item \textbf{Distance épicentrale} $\Delta$ (en degrés ou en km).
    \item \textbf{Paramètre de rayon} $p$ (ou \emph{ray parameter}) et son unité au SI.
    \item \textbf{Discontinuité de premier ordre} vs \textbf{discontinuité de second ordre} (donner un exemple de chacune dans la structure terrestre).
    \item \textbf{Phase PKIKP} : décrire le chemin complet de l'onde dans le globe.
\end{enumerate}
\end{exercice}

\begin{exercice}[Calcul direct : Vitesse et profondeur]
On considère un modèle simplifié de croûte continentale à deux couches horizontales sur un demi-espace (manteau supérieur).
\begin{center}
\begin{tabular}{|c|c|c|c|}\hline
Couche & Épaisseur (km) & $v_P$ (km/s) & $v_S$ (km/s) \\ \hline
Sédiments & 2 & 3,0 & 1,5 \\
Croûte cristalline & 30 & 6,2 & 3,6 \\
Manteau (demi-espace) & $\infty$ & 8,1 & 4,6 \\ \hline
\end{tabular}
\end{center}
Une onde P verticale (incidence normale) partie de la surface traverse ces couches.
\begin{enumerate}
    \item Calculer le temps de parcours total pour atteindre le Moho (base de la croûte).
    \item Calculer le temps de parcours pour atteindre une profondeur de 50 km (dans le manteau).
    \item En déduire l'expression du temps réduit $t^* = t - \frac{\Delta}{v_{\text{ref}}}$ pour un rayon vertical ($\Delta = 0$) et discuter de son intérêt pour mettre en évidence les discontinuités.
\end{enumerate}
\end{exercice}

\courssub{Je m'entraîne}

\begin{exercice}[Exploitation de courbes de temps de parcours : Discontinuité de Gutenberg]
On donne ci-dessous des temps de parcours $t$ (en secondes) observés pour les phases \textbf{P} (onde directe dans le manteau) et \textbf{PcP} (onde P réfléchie sur la frontière noyau-manteau) en fonction de la distance épicentrale $\Delta$ (en degrés).

\begin{center}
\begin{tabular}{|c|cccccccc|}\hline
$\Delta$ ($^\circ$) & 0 & 20 & 40 & 60 & 80 & 100 & 120 & 140 \\ \hline
$t_P$ (s) & 0 & 380 & 720 & 1020 & 1290 & 1530 & 1740 & 1920 \\
$t_{PcP}$ (s) & 660 & 720 & 900 & 1140 & 1400 & 1660 & 1900 & 2100 \\ \hline
\end{tabular}
\end{center}
\begin{enumerate}
    \item Tracer sur un même repère orthonormé les courbes $t_P(\Delta)$ et $t_{PcP}(\Delta)$.
    \item À quelle distance $\Delta$ les deux courbes se coupent-elles ? Interpréter physiquement ce point de rencontre.
    \item En utilisant la méthode du \textbf{temps réduit} $t_{red} = t - \frac{\Delta}{v_{\text{ref}}}$ avec $v_{\text{ref}} = 11,1$ km/s (vitesse moyenne approximative dans le manteau inférieur), tracer $t_{red}(\Delta)$ pour la phase PcP. En déduire la profondeur de la discontinuité de Gutenberg (rayon du noyau $R_c$). On prendra le rayon terrestre $R_T = 6371$ km.
    \item Rappeler la relation liant le paramètre de rayon $p$, la distance $\Delta$ et le temps de parcours $t$ : $p = \frac{dt}{d\Delta}$ (en s/rad ou s/$^\circ$). Estimer $p$ pour la phase PcP à $\Delta = 60^\circ$.
\end{enumerate}
\end{exercice}

\begin{exercice}[Analyse d'un sismogramme théorique : Station YAO (Yaoundé)]
Un séisme de foyer superficiel ($h \approx 10$ km) se produit à une distance épicentrale $\Delta = 75^\circ$ de la station sismologique de Yaoundé (YAO). On enregistre les composantes verticale (Z) et horizontales (N, E). Les arrivées principales observées sur la composante Z sont (temps après l'origine du séisme) :
\begin{itemize}
    \item $t_1 = 10$ min 10 s : impulsion compressive (premier mouvement vers le haut), amplitude forte.
    \item $t_2 = 13$ min 45 s : impulsion dilatative, amplitude moyenne.
    \item $t_3 = 19$ min 20 s : impulsion compressive, amplitude faible.
    \item $t_4 = 22$ min 10 s : début d'ondes de grande amplitude, période plus longue ($\approx 15$ s).
\end{itemize}
\begin{enumerate}
    \item Identifier les phases sismiques correspondantes parmi : P, pP, sP, PcP, PKP, PP, S, ScS, SS, Rayleigh (R). Justifier chaque identification par le temps relatif, la polarité (compressive/dilatative) et la nature de l'onde (volume vs surface).
    \item Calculer le temps de parcours théorique de l'onde P pour $\Delta = 75^\circ$ en utilisant la table de Jeffreys-Bullen simplifiée ($t_P \approx 10$ min 10 s). Vérifier la cohérence avec $t_1$.
    \item Expliquer pourquoi la phase \textbf{pP} (réflexion en surface au point de départ) arrive après P et a une polarité opposée si le mécanisme au foyer est une faille normale.
    \item Pourquoi les ondes S (phase $t_4$) arrivent-elles après les ondes P ? Rapporter le rapport $v_P/v_S$ typique du manteau supérieur.
\end{enumerate}
\end{exercice}

\begin{exercice}[Modélisation du trajet d'onde : Réfraction dans le noyau]
On considère un rayon d'onde P qui pénètre dans le noyau externe liquide. Au point d'entrée (discontinuité de Gutenberg, rayon $r_c = 3480$ km), l'angle d'incidence dans le manteau est $i_m = 30^\circ$. La vitesse P dans le manteau juste au-dessus est $v_{P,m} = 13,7$ km/s. La vitesse P dans le noyau externe juste en dessous est $v_{P,c} = 8,1$ km/s.
\begin{enumerate}
    \item Calculer le paramètre de rayon $p$ (en s/km ou s/rad). Rappeler sa conservation.
    \item Calculer l'angle de réfraction $i_c$ dans le noyau externe (loi de Snell-Descartes).
    \item L'onde P se propage-t-elle plus "verticalement" ou plus "horizontalement" dans le noyau qu'au manteau ? Justifier.
    \item Tracer le chemin du rayon dans une coupe méridienne (quart de cercle Terre). Indiquer les angles $i_m$ et $i_c$.
    \item Que devient une onde S incidente sur cette même discontinuité ? Pourquoi ?
\end{enumerate}
\end{exercice}

\begin{exercice}[Preuve de la liquidité du noyau externe : Zones d'ombre]
Les tables de temps de parcours (Jeffreys-Bullen ou modèle PREM) montrent que :
\begin{itemize}
    \item Les ondes P sont enregistrées jusqu'à $\Delta = 103^\circ$, puis disparaissent (zone d'ombre P) jusqu'à $\Delta \approx 143^\circ$, puis réapparaissent (phases PKP, PKIKP).
    \item Les ondes S ne sont \textbf{jamais} enregistrées au-delà de $\Delta = 103^\circ$ (zone d'ombre S totale).
\end{itemize}
\begin{enumerate}
    \item Expliquer géométriquement, à l'aide d'un schéma de rayons dans une Terre sphérique à noyau liquide, pourquoi il existe une zone d'ombre pour les ondes P. Calculer l'angle critique d'incidence $i_c^{\text{crit}}$ à la frontière noyau-manteau pour la réfraction P $\to$ P (noyau) sachant $v_{P,\text{manteau}} \approx 13,7$ km/s et $v_{P,\text{noyau}} \approx 8,1$ km/s.
    \item Montrer que la distance épicentrale minimale de la zone d'ombre P ($\Delta_{\min} \approx 103^\circ$) correspond aux rayons tangents au noyau (p = $R_c / v_{P,\text{manteau}}$).
    \item Pourquoi la zone d'ombre des ondes S est-elle \textbf{totale} au-delà de $103^\circ$ (aucune phase SKS, SKKS n'arrive directement par le noyau) ? En déduire l'état physique du noyau externe.
    \item Citer la phase qui permet néanmoins de "voir" le noyau interne solide (PKIKP) et expliquer pourquoi elle traverse le noyau externe liquide.
\end{enumerate}
\end{exercice}

\courssub{Je me prépare au Probatoire}

\begin{exercice}[Sujet type Probatoire : Détermination de la profondeur du noyau par la méthode des temps réduits]
\textbf{Contexte :} Lors d'une campagne sismologique au Cameroun, la station de \textbf{Yaoundé (YAO)} enregistre un séisme télésismique ($\Delta = 85^\circ$). On exploite les temps de parcours des phases \textbf{P} (manteau) et \textbf{PcP} (réflexion noyau-manteau) pour plusieurs distances.

\textbf{Données :} Temps de parcours $t$ (en secondes) pour les phases P et PcP.
\begin{center}
\begin{tabular}{|c|ccccccc|}\hline
$\Delta$ ($^\circ$) & 30 & 40 & 50 & 60 & 70 & 80 & 90 \\ \hline
$t_P$ (s) & 530 & 720 & 910 & 1090 & 1260 & 1420 & 1570 \\
$t_{PcP}$ (s) & 850 & 950 & 1090 & 1260 & 1450 & 1650 & 1860 \\ \hline
\end{tabular}
\end{center}
\textbf{Rappels :} Rayon terrestre moyen $R_T = 6371$ km. Vitesse de référence pour le temps réduit : $v_{\text{ref}} = 11,1$ km/s. Relation $\Delta$ (rad) = $\Delta^\circ \times \pi/180$. Distance angulaire en radians : $x = \Delta \times \pi/180$.

\begin{enumerate}
    \item \textbf{Représentation graphique} : Tracer, sur un même graphique (papier millimétré ou logiciel), les courbes $t_P(\Delta)$ et $t_{PcP}(\Delta)$. Commenter la forme des courbes et leur écart.
    \item \textbf{Méthode des temps réduits} :
    \begin{enumerate}
        \item Calculer le temps réduit $t_{red} = t - \frac{R_T \cdot x}{v_{\text{ref}}}$ pour la phase PcP à chaque distance. Compléter le tableau.
        \item Tracer la courbe $t_{red}(\Delta)$ pour PcP. Quelle allure présente cette courbe au voisinage de $\Delta = 0^\circ$ (extrapolation) ?
        \item Montrer que pour une réflexion sur une interface horizontale (approximation locale) ou sphérique (noyau), le temps réduit à distance nulle $t_{red}(0)$ est égal au temps de parcours aller-retour vertical dans le manteau au-dessus du réflecteur : $t_{red}(0) = 2 \int_0^{R_T - R_c} \frac{dr}{v_P(r)}$.
    \end{enumerate}
    \item \textbf{Calcul de la profondeur du noyau} :
    En supposant une vitesse moyenne $\bar{v}_P = 10,5$ km/s dans le manteau au-dessus du noyau (de la surface à $R_c$), utiliser la valeur extrapolée $t_{red}(0) \approx 660$ s pour calculer le rayon du noyau $R_c$ et sa profondeur $z_c = R_T - R_c$.
    \item \textbf{Validation par la zone d'ombre} : La zone d'ombre P commence à $\Delta \approx 103^\circ$. Vérifier que cette valeur est cohérente avec le $R_c$ trouvé, en utilisant la relation $\sin i_c^{\text{crit}} = \frac{v_{P,\text{noyau}}}{v_{P,\text{manteau}}}$ et $\Delta_{\min} = 2 \left( \frac{\pi}{2} - i_c^{\text{crit}} \right)$ (en radians) pour un modèle à deux couches à vitesse constante. Prendre $v_{P,\text{manteau}} = 13,7$ km/s, $v_{P,\text{noyau}} = 8,1$ km/s.
    \item \textbf{Conclusion} : Rédiger une conclusion structurée : "Les données sismiques permettent de localiser la discontinuité de Gutenberg à une profondeur de ... km, confirmant l'existence d'un noyau de rayon ... km. Le saut de vitesse P (chute brutale) et l'absence d'ondes S transmises prouvent que le noyau externe est ..."
\end{enumerate}
\end{exercice}

\begin{exercice}[Sujet type Probatoire : Analyse du modèle PREM et argumentation sur l'état du noyau]
\textbf{Document :} Extrait du modèle de référence PREM (Preliminary Reference Earth Model) donnant les vitesses $v_P(r)$ et $v_S(r)$ en fonction du rayon $r$ (km) au voisinage du noyau.

\begin{center}
\begin{tabular}{|c|c|c|}\hline
Profondeur (km) / Rayon $r$ (km) & $v_P$ (km/s) & $v_S$ (km/s) \\ \hline
2885 (Base manteau) & 13,72 & 7,27 \\
2891 (Sommet noyau externe) & 8,06 & \textbf{0,00} \\
4000 (Milieu noyau externe) & 8,70 & 0,00 \\
5150 (Base noyau externe) & 10,35 & 0,00 \\
5150 (Sommet noyau interne) & 11,00 & 3,50 \\
6371 (Centre) & 11,26 & 3,67 \\ \hline
\end{tabular}
\end{center}

\begin{enumerate}
    \item \textbf{Analyse des discontinuités} :
    \begin{enumerate}
        \item Identifier les deux discontinuités majeures visibles dans ce tableau. Donner leur nom usuel et leur profondeur.
        \item Calculer le contraste relatif de vitesse $\frac{\Delta v_P}{v_P}$ pour chacune. Laquelle est la plus marquée pour $v_P$ ?
        \item Commenter l'évolution de $v_S$ : que signifie $v_S = 0$ dans le noyau externe ? Quelle propriété physique du milieu cela implique-t-il (module de cisaillement $\mu$) ?
    \end{enumerate}
    \item \textbf{Interprétation pétrologique et physique} :
    \begin{enumerate}
        \item La chute brutale de $v_P$ à 2891 km (de 13,7 à 8,1 km/s) s'accompagne d'une augmentation de la densité (de $\approx 5,5$ à $\approx 9,9$ g/cm$^3$). Rappeler la formule de la vitesse des ondes P : $v_P = \sqrt{\frac{K + \frac{4}{3}\mu}{\rho}}$. Expliquer comment $v_P$ peut diminuer alors que $\rho$ augmente fortement. En déduire l'évolution du module d'incompressibilité $K$.
        \item Le noyau externe est liquide (fer-nickel fondu + éléments légers). Le noyau interne est solide (fer-nickel cristallisé). Justifier la réapparition de $v_S > 0$ et le saut de $v_P$ à 5150 km par le changement d'état physique.
    \end{enumerate}
    \item \textbf{Anisotropie du noyau interne (Complément)} : Les ondes PKIKP (traversant le noyau interne) arrivent plus tôt quand elles parcourent l'axe polaire que le plan équatorial (différence $\approx 3-4$ s). Proposer une hypothèse structurale (texture cristalline) pour expliquer cette anisotropie sismique.
    \item \textbf{Synthèse} : Rédiger un paragraphe argumenté (10-15 lignes) résumant comment la sismologie (vitesses P, S, zones d'ombre, phases réfléchies/réfractées) a permis de construire le modèle actuel : Croûte - Manteau - Noyau externe liquide - Noyau interne solide.
\end{enumerate}
\end{exercice}

\begin{exercice}[Sujet type Probatoire : Comparaison Terre-Lune et champ magnétique]
\textbf{Contexte :} Les missions \textbf{Apollo} (1969-1972) ont déposé des sismomètres sur la Lune (réseau ALSEP). Les données lunaires révèlent une structure interne très différente de la Terre.

\textbf{Données comparées :}
\begin{center}
\begin{tabular}{|l|c|c|}\hline
Paramètre & Terre & Lune \\ \hline
Rayon moyen (km) & 6371 & 1737 \\
Masse ($10^{24}$ kg) & 5,97 & 0,073 \\
Densité moyenne (g/cm$^3$) & 5,51 & 3,34 \\
Croûte : épaisseur moyenne (km) & 30 (cont.) / 7 (oc.) & 40-60 (face visible) / $\sim$ 80 (face cachée) \\
Manteau : épaisseur (km) & $\sim$ 2850 & $\sim$ 1000 \\
Noyau : rayon (km) & 3480 (total) & $\sim$ 350-400 (total) \\
Noyau : état & Externe liquide / Interne solide & Probablement \textbf{entièrement liquide} ou petit noyau solide central \\
Champ magnétique global actuel & Oui (dynamo active) & Non (vestiges crustaux uniquement) \\ \hline
\end{tabular}
\end{center}

\begin{enumerate}
    \item \textbf{Structure en couches} :
    \begin{enumerate}
        \item Calculer le rapport $\frac{\text{Rayon noyau}}{\text{Rayon planète}}$ pour la Terre et la Lune. Commenter la différence.
        \item La croûte lunaire est proportionnellement plus épaisse que la croûte terrestre. Relier cette observation à l'histoire thermique précoce (océan de lave, cristallisation flottante de l'anorthosite) et à l'absence de tectonique des plaques.
    \end{enumerate}
    \item \textbf{Noyau et champ magnétique} :
    \begin{enumerate}
        \item Rappeler les \textbf{trois conditions} nécessaires au fonctionnement d'une \emph{dynamo planétaire} (champ magnétique auto-entretenu).
        \item En utilisant le tableau, expliquer pourquoi la Lune n'a \textbf{plus} de champ magnétique global aujourd'hui, alors qu'elle en a possédé un il y a 3-4 Ga (magnétisation des roches crustales).
        \item Discuter : La petite taille du noyau lunaire et son état probablement liquide (mais non en convection vigoureuse) sont-ils compatibles avec l'arrêt de la dynamo ?
    \end{enumerate}
    \item \textbf{Sismologie lunaire vs terrestre} :
    Les séismes lunaires (moonquakes) sont de faible magnitude ($M < 5$), profonds (700-1000 km) et liés aux marées terrestres. Les ondes S lunaires s'atténuent très peu ($Q_S$ très élevé $\approx 10^4$) comparé à la Terre ($Q_S \approx 100-300$).
    \begin{enumerate}
        \item Expliquer la faible atténuation lunaire par l'absence d'eau et la température plus basse du manteau lunaire.
        \item Pourquoi la sismologie lunaire a-t-elle plus de mal à préciser la structure du noyau que la sismologie terrestre ?
    \end{enumerate}
    \item \textbf{Conclusion} : En 15 lignes maximum, synthétiser les différences majeures de structure et d'évolution entre la Terre et la Lune, en insistant sur le rôle de la \textbf{taille}, de la \textbf{chaleur interne} et de la \textbf{présence d'eau} pour la dynamique planétaire (tectonique, dynamo, vie).
\end{enumerate}
\end{exercice}

\newpage

\coursec{Corrigés des exercices}

\courssub{Exercices et problèmes corrigés}

\begin{corrige}[Exercice 1 --- QCM : Propriétés des ondes P et S]
L'affirmation \textbf{fausse} est l'affirmation \textbf{4}.

\textbf{Analyse détaillée :}
\begin{enumerate}
    \item \textbf{Vrai.} Les ondes P sont des ondes de compression-dilatation (ondes longitudinales) ; leur vitesse s'écrit $v_P = \sqrt{\dfrac{K + \frac{4}{3}\mu}{\rho}}$, qui dépend bien du module d'incompressibilité $K$, du module de cisaillement $\mu$ et de la densité $\rho$.
    \item \textbf{Vrai.} Les ondes S sont des ondes de cisaillement (ondes transversales) de vitesse $v_S = \sqrt{\mu/\rho}$. Dans un fluide parfait, $\mu = 0$, donc $v_S = 0$ : les ondes S ne s'y propagent pas.
    \item \textbf{Vrai.} Comme $K > 0$ et $\mu > 0$ dans un solide, $v_P = \sqrt{\dfrac{K + \frac{4}{3}\mu}{\rho}} > \sqrt{\dfrac{\mu}{\rho}} = v_S$. Le rapport typique est $v_P/v_S \approx \sqrt{3} \approx 1,73$ pour un solide de Poisson.
    \item \textbf{Faux.} À une discontinuité de premier ordre, une onde incidente subit \textbf{à la fois} une réflexion \textbf{et} une réfraction (transmission), avec en général \textbf{conversion d'onde} : une onde P incidente génère une P réfléchie, une S réfléchie (PS), une P transmise et une S transmise (si le milieu récepteur est solide). Dire « uniquement une réflexion » est donc incorrect.
\end{enumerate}
\end{corrige}

\begin{attention}[Point de vigilance]
Au Probatoire, ne jamais oublier la \cle{conversion d'onde} (P $\to$ S et S $\to$ P) aux interfaces : c'est elle qui explique l'existence de phases comme PS, SP, ou les conversions à la frontière noyau-manteau.
\end{attention}

\begin{corrige}[Exercice 2 --- Vrai ou Faux : Comportement aux discontinuités]
\begin{enumerate}
    \item \textbf{Faux.} Par les conditions de continuité du déplacement et de la contrainte à l'interface, une onde P incidente sur le Moho génère \textbf{quatre} ondes : P réfléchie, S réfléchie (conversion), P réfractée et S réfractée. La conversion P $\to$ S est d'autant plus importante que l'incidence est oblique ; elle n'est nulle qu'à incidence strictement normale.
    \item \textbf{Vrai.} Dans un milieu sphériquement stratifié (vitesse ne dépendant que du rayon $r$), la généralisation sphérique de la loi de Snell-Descartes donne l'invariant de Benndorf :
    \[ p = \frac{r \sin i}{v(r)} = \text{constante le long du rayon}. \]
    (Dans un milieu plan stratifié horizontalement, $p = \sin i / v$.) C'est le \cle{paramètre de rayon}, conservé sur tout le trajet.
    \item \textbf{Vrai.} Au-delà de $\Delta \approx 103^\circ$, tout rayon direct devrait traverser le noyau externe ; or aucune onde S n'est observée. Comme les ondes S ne se propagent pas dans un fluide ($\mu = 0$), l'absence totale de S au-delà de $103^\circ$ démontre directement que le noyau externe est \textbf{liquide}.
    \item \textbf{Vrai.} La discontinuité de Lehmann (vers 5150 km de profondeur) correspond au passage noyau externe liquide $\to$ noyau interne solide : $\mu$ redevient non nul, donc $v_S$ réapparaît ($\approx 3,5$ km/s) et $v_P$ augmente brutalement (de $\approx 10,4$ à $\approx 11,0$ km/s), car le terme $\frac{4}{3}\mu$ s'ajoute au numérateur de $v_P$.
\end{enumerate}
\end{corrige}

\begin{corrige}[Exercice 3 --- Définitions et notions clés]
\begin{enumerate}
    \item \textbf{Temps de parcours} : durée écoulée entre l'émission de l'onde au foyer (instant origine) et son arrivée à la station, pour une phase donnée. Il s'exprime en secondes (ou minutes) et dépend de la distance épicentrale, de la profondeur du foyer et du modèle de vitesses.
    \item \textbf{Distance épicentrale} $\Delta$ : angle au centre de la Terre entre l'épicentre et la station (en degrés), ou longueur d'arc correspondante à la surface (en km ; $1^\circ \approx 111,19$ km).
    \item \textbf{Paramètre de rayon} $p$ : invariant du rayon dans un milieu stratifié, $p = \dfrac{r \sin i}{v}$ (Terre sphérique) ou $p = \dfrac{\sin i}{v}$ (Terre plate). Unité SI : $\mathrm{s/m}$ (souvent exprimé en s/km, s/rad ou s/$^\circ$). Il vaut aussi $p = \dfrac{dt}{d\Delta}$ : c'est la pente de la courbe de temps de parcours.
    \item \textbf{Discontinuité de premier ordre} : saut brutal (discontinuité mathématique) de la vitesse, de la densité ou des modules élastiques. Exemple : le \textbf{Moho} (croûte/manteau), la discontinuité de \textbf{Gutenberg} (manteau/noyau, 2891 km). \textbf{Discontinuité de second ordre} : changement brutal du \emph{gradient} de vitesse (la vitesse elle-même reste continue). Exemple : les discontinuités de \textbf{410 km} et \textbf{660 km} (transitions de phase minéralogiques dans le manteau).
    \item \textbf{Phase PKIKP} : onde P qui traverse successivement le \textbf{manteau} (P), le \textbf{noyau externe} (K, de l'allemand \emph{Kern}), le \textbf{noyau interne} (I), puis à nouveau le noyau externe (K) et le manteau (P). C'est la phase qui prouve l'existence et la solidité du noyau interne.
\end{enumerate}
\end{corrige}

\begin{corrige}[Exercice 4 --- Calcul direct : Vitesse et profondeur]
Pour une propagation verticale, le temps de parcours dans une couche d'épaisseur $e$ traversée à la vitesse $v$ est $t = e/v$.

\textbf{1. Temps pour atteindre le Moho} (base de la croûte, profondeur $2 + 30 = 32$ km) :
\[ t_{\text{Moho}} = \frac{2}{3{,}0} + \frac{30}{6{,}2} = 0{,}667 + 4{,}839 \approx 5{,}51 \ \text{s}. \]

\textbf{2. Temps pour atteindre 50 km de profondeur} : il faut ajouter la traversée de $50 - 32 = 18$ km de manteau à $8,1$ km/s :
\[ t_{50} = 5{,}51 + \frac{18}{8{,}1} = 5{,}51 + 2{,}22 \approx 7{,}73 \ \text{s}. \]

\textbf{3. Temps réduit pour un rayon vertical.} Pour $\Delta = 0$, le terme de réduction est nul :
\[ t^* = t - \frac{\Delta}{v_{\text{ref}}} = t \quad (\Delta = 0). \]
\textbf{Intérêt du temps réduit :} en soustrayant le terme linéaire $\Delta / v_{\text{ref}}$, on « aplatit » les courbes de temps de parcours ; les branches associées aux différentes couches ou réflexions deviennent quasi horizontales et les \textbf{ruptures de pente} (discontinuités) ou les \textbf{branches de réflexion} apparaissent beaucoup plus clairement. C'est l'outil standard pour repérer et situer les interfaces.
\end{corrige}

\begin{corrige}[Exercice 5 --- Exploitation de courbes de temps de parcours : Discontinuité de Gutenberg]
\textbf{1. Tracé.} On porte $\Delta$ en abscisse (0 à $140^\circ$) et $t$ en ordonnée (0 à 2200 s). La courbe $t_P(\Delta)$ est croissante et \textbf{concave vers le bas} (la pente $p = dt/d\Delta$ diminue quand $\Delta$ augmente, car les rayons profonds sont plus rapides). La courbe $t_{PcP}(\Delta)$ est croissante, située au-dessus pour $\Delta$ faible, avec une pente qui augmente (branche de réflexion).

\textbf{2. Intersection.} Les deux courbes se coupent entre $\Delta = 0^\circ$ et $20^\circ$ : à $\Delta = 0$, $t_{PcP} = 660 > t_P = 0$ ; à $\Delta = 20^\circ$, $t_P = 380 < t_{PcP} = 720$. En comparant les pentes, l'écart se réduit continûment ; par interpolation grossière, l'intersection graphique se situe vers $\Delta \approx 15$--$20^\circ$ (zone où $t_P \approx t_{PcP} \approx 600$--700 s). \textbf{Interprétation physique :} à cette distance, l'onde directe P et l'onde réfléchie PcP arrivent \emph{en même temps} : les deux trajets (direct dans le manteau et aller-retour jusqu'au noyau) sont équivalents en durée. Au-delà, la phase directe P arrive avant PcP.

\textbf{3. Temps réduit et rayon du noyau.} Conversion : $1^\circ$ d'arc à la surface vaut $R_T \pi/180 = 6371 \times 0{,}017453 \approx 111{,}2$ km. Le terme de réduction est $\dfrac{111{,}2 \times \Delta^\circ}{11{,}1} \approx 10{,}02 \times \Delta^\circ$ (en secondes).

\begin{center}
\begin{tabular}{|c|cccccccc|}\hline
$\Delta$ ($^\circ$) & 0 & 20 & 40 & 60 & 80 & 100 & 120 & 140 \\ \hline
$\Delta / v_{\text{ref}}$ (s) & 0 & 200 & 400 & 601 & 801 & 1002 & 1202 & 1403 \\
$t_{PcP}$ (s) & 660 & 720 & 900 & 1140 & 1400 & 1660 & 1900 & 2100 \\ \hline
$t_{red}$ (s) & 660 & 520 & 500 & 539 & 599 & 658 & 698 & 697 \\ \hline
\end{tabular}
\end{center}

La courbe $t_{red}(\Delta)$ présente un \textbf{minimum} vers $\Delta \approx 40^\circ$ puis remonte : c'est l'allure typique d'une branche de réflexion réduite. L'extrapolation à $\Delta = 0$ donne $t_{red}(0) \approx 660$ s, égal au \textbf{temps aller-retour vertical} dans le manteau jusqu'au noyau.

En supposant une vitesse moyenne $\bar v_P \approx 10{,}5$ km/s dans le manteau :
\[ 2\,\frac{R_T - R_c}{\bar v_P} = 660 \ \text{s} \quad\Rightarrow\quad R_T - R_c = \frac{660 \times 10{,}5}{2} = 3465 \ \text{km}. \]
\[ R_c = 6371 - 3465 \approx 2906 \ \text{km} \quad\text{soit une profondeur } z_c \approx 3465 \ \text{km}. \]
Ordre de grandeur correct (valeur réelle : $R_c = 3480$ km, $z_c = 2891$ km) ; l'écart vient de l'approximation de vitesse constante.

\textbf{4. Paramètre de rayon.} $p = \dfrac{dt}{d\Delta}$. À $\Delta = 60^\circ$, on estime la pente de PcP par différence finie :
\[ p \approx \frac{t(80) - t(40)}{80 - 40} = \frac{1400 - 900}{40} = 12{,}5 \ \text{s}/^\circ. \]
Soit $p \approx 12{,}5 \times \dfrac{180}{\pi} \approx 716$ s/rad, ou encore $p \approx \dfrac{12{,}5}{111{,}2} \approx 0{,}112$ s/km.
\end{corrige}

\begin{corrige}[Exercice 6 --- Analyse d'un sismogramme théorique : Station YAO]
\textbf{1. Identification des phases.}

\begin{center}
\begin{tabularx}{\linewidth}{|c|l|X|}\hline
\rowcolor{popblueL} Arrivée & Phase & Justification \\ \hline
$t_1 = 10$ min 10 s & \textbf{P} & Première arrivée (les P sont les plus rapides), impulsion \emph{compressive} (premier mouvement vers le haut) cohérente avec une P directe à $\Delta = 75^\circ$. \\ \hline
$t_2 = 13$ min 45 s & \textbf{pP} & Arrive $\approx 3$ min 35 s après P : c'est la P réfléchie en surface au-dessus du foyer ; le retard correspond au double du temps de montée depuis le foyer ($h \approx 10$ km) plus trajet de profondeur ; polarité \emph{dilatative} (inversion de phase par réflexion en surface libre). \\ \hline
$t_3 = 19$ min 20 s & \textbf{PP} & Onde P réfléchie une fois à la surface à mi-parcours ; arrive nettement après P, amplitude faible (double trajet + réflexion), polarité compressive. \\ \hline
$t_4 = 22$ min 10 s & \textbf{S} puis train d'\textbf{ondes de Rayleigh} & Les ondes S (volume, $v_S < v_P$) arrivent après les P ; les grandes amplitudes à longue période ($\approx 15$ s) en fin de sismogramme sont caractéristiques des \textbf{ondes de surface} (Rayleigh sur la composante Z). \\ \hline
\end{tabularx}
\end{center}

(On écarte PcP et PKP : à $\Delta = 75^\circ < 103^\circ$, on est avant la zone d'ombre, mais PcP à $75^\circ$ arriverait vers 22--23 min et se confondrait avec le train S ; ScS arrive encore plus tard. La réponse attendue au niveau Première est P, pP, PP, S/Rayleigh.)

\textbf{2. Cohérence de $t_1$.} La table de Jeffreys-Bullen simplifiée donne $t_P(75^\circ) \approx 10$ min 10 s, identique à $t_1$ : l'identification de la première arrivée comme phase P est cohérente.

\textbf{3. Retard et polarité de pP.} La phase pP monte d'abord du foyer jusqu'à la surface (distance $\approx h$), s'y réfléchit, puis suit un trajet quasi identique à P. Le retard $\Delta t \approx 2h / v_P \times (\text{facteur d'incidence})$ explique l'arrivée tardive. La réflexion sur la \textbf{surface libre} (impédance nulle) introduit un \textbf{déphasage de $\pi$} (inversion de polarité) : pour un mécanisme de faille normale, si P arrive compressive à la station, pP arrive dilatative.

\textbf{4. Retard des ondes S.} Dans le manteau supérieur, $v_P \approx 8$ km/s et $v_S \approx 4,5$ km/s, soit $v_P / v_S \approx 1{,}7$--$1{,}8$ ($\approx \sqrt{3}$). Les ondes S, plus lentes, parcourent le même trajet en un temps environ $1{,}8$ fois plus long : elles arrivent donc systématiquement après les P, avec un écart S--P qui croît avec la distance (principe de la localisation des séismes).
\end{corrige}

\begin{corrige}[Exercice 7 --- Modélisation du trajet d'onde : Réfraction dans le noyau]
\textbf{1. Paramètre de rayon.} À la frontière noyau-manteau ($r = r_c = 3480$ km) :
\[ p = \frac{r_c \sin i_m}{v_{P,m}} = \frac{3480 \times \sin 30^\circ}{13{,}7} = \frac{3480 \times 0{,}5}{13{,}7} = \frac{1740}{13{,}7} \approx 127 \ \text{s}. \]
(unité : s en coordonnées sphériques, i.e. s/rad). $p$ est \textbf{conservé} le long de tout le trajet du rayon (invariant de Benndorf), y compris lors des réfractions et réflexions.

\textbf{2. Angle de réfraction.} Loi de Snell-Descartes à l'interface (le facteur $r_c$ est identique des deux côtés) :
\[ \frac{\sin i_m}{v_{P,m}} = \frac{\sin i_c}{v_{P,c}} \quad\Rightarrow\quad \sin i_c = \frac{v_{P,c}}{v_{P,m}} \sin i_m = \frac{8{,}1}{13{,}7} \times 0{,}5 \approx 0{,}296. \]
\[ i_c \approx \arcsin(0{,}296) \approx 17{,}2^\circ. \]

\textbf{3. Direction de propagation.} $i_c < i_m$ : le rayon se \textbf{rapproche de la normale} à l'interface, donc se propage plus \textbf{verticalement} (plus radialement) dans le noyau que dans le manteau. C'est la conséquence de la \emph{chute} de vitesse ($8{,}1 < 13{,}7$ km/s) : le noyau externe est un milieu « réfringent » vers le bas, ce qui courbe les rayons et crée la zone d'ombre des P.

\textbf{4. Schéma du trajet.} La figure ci-dessous illustre la réfraction à la frontière noyau-manteau.
\end{corrige}

\begin{popfigure}
\begin{center}
\begin{tikzpicture}[scale=1.1]
% surface terrestre (quart de cercle)
\draw[thick, popblue] (4.5,0) arc (0:180:4.5);
\draw[thick, poporange] (2.45,0) arc (0:180:2.45);
% rayon incident dans le manteau
\draw[very thick, popgreen, -{Stealth}] (-4.2,1.55) -- (-1.70,2.45);
% rayon réfracté dans le noyau
\draw[very thick, popgreen, -{Stealth}] (-1.70,2.45) -- (0.4,-1.6);
% normale (radiale)
\draw[dashed, popmuted] (0,0) -- (-2.12,3.06);
\draw[dashed, popmuted] (0,0) -- (0.62,-2.48);
% angles
\draw[popdark] (-2.35,2.62) arc (125:160:0.75);
\node[popdark, font=\small] at (-2.75,2.35) {$i_m = 30^\circ$};
\draw[popdark] (-1.28,1.83) arc (235:262:0.85);
\node[popdark, font=\small] at (-0.75,1.35) {$i_c \approx 17^\circ$};
% labels
\node[popblue, font=\small] at (3.6,3.4) {Manteau ($v_P = 13{,}7$ km/s)};
\node[poporange, font=\small] at (0,-0.7) {Noyau externe};
\node[poporange, font=\small] at (0,-1.15) {($v_P = 8{,}1$ km/s)};
\node[popmuted, font=\small] at (-3.3,0.35) {Gutenberg ($r_c = 3480$ km)};
\node[popgreen, font=\small] at (-3.9,2.2) {P incidente};
\node[popgreen, font=\small] at (1.3,-1.9) {P réfractée (PKP)};
\fill (0,0) circle (1pt);
\end{tikzpicture}
\end{center}
\legende{Coupe méridienne : réfraction d'une onde P à la frontière noyau-manteau. La vitesse chute dans le noyau externe, donc le rayon se rapproche de la normale ($i_c < i_m$).}
\end{popfigure}

\begin{corrige}[Exercice 7 (suite) --- Onde S incidente]
\textbf{5. Onde S incidente.} Le noyau externe est liquide : $\mu = 0$ donc $v_S = 0$. Il n'existe \textbf{aucune onde S transmise} ; l'onde S incidente est intégralement réfléchie (ou convertie en P transmise, phase SKP impossible : en réalité une S incidente génère une S réfléchie et une P réfractée --- conversion S $\to$ P --- mais pas de S dans le noyau). C'est la cause de la zone d'ombre totale des S.
\end{corrige}

\begin{corrige}[Exercice 8 --- Preuve de la liquidité du noyau externe : Zones d'ombre]
\textbf{1. Origine géométrique de la zone d'ombre P et angle critique.} Comme $v_P$ \emph{diminue} à l'entrée du noyau, les rayons sont réfractés vers la normale (vers le bas) : ils « plongent » et ressortent plus loin qu'ils ne le feraient sans noyau. Les rayons de plus en plus profonds dans le manteau finissent par être tangents au noyau ; aucun rayon direct ne peut émerger entre $\Delta \approx 103^\circ$ et $\approx 143^\circ$.

L'angle critique de réfraction (transmission rasante dans le noyau) est donné par :
\[ \sin i_c^{\text{crit}} = \frac{v_{P,\text{noyau}}}{v_{P,\text{manteau}}} = \frac{8{,}1}{13{,}7} \approx 0{,}591 \quad\Rightarrow\quad i_c^{\text{crit}} \approx 36{,}2^\circ. \]

\textbf{2. Distance minimale de la zone d'ombre.} Pour un modèle à vitesse constante dans le manteau, le rayon limite est \textbf{tangent au noyau} : à son point le plus profond, $i = 90^\circ$, donc
\[ p = \frac{R_c \sin 90^\circ}{v_{P,\text{manteau}}} = \frac{R_c}{v_{P,\text{manteau}}}. \]
Géométriquement, un rayon tangent au noyau subit une déflexion totale
\[ \Delta_{\min} = 2\left(\frac{\pi}{2} - i_c^{\text{crit}}\right) = 2 \times (90^\circ - 36{,}2^\circ) = 2 \times 53{,}8^\circ \approx 107{,}6^\circ, \]
proche de la valeur observée $103^\circ$ (l'écart vient du gradient de vitesse réel dans le manteau, qui courbe les rayons). La correspondance confirme que la zone d'ombre est bordée par les rayons tangents au noyau.

\textbf{3. Zone d'ombre S totale.} Tout trajet direct vers $\Delta > 103^\circ$ traverse obligatoirement le noyau externe. Or les ondes S ne se propagent pas dans un liquide ($\mu = 0 \Rightarrow v_S = 0$) : aucune phase S transmise (ni SKS au sens S-dans-le-noyau) ne peut exister. La zone d'ombre S est donc \textbf{totale} de $103^\circ$ à $180^\circ$. \textbf{Conclusion : le noyau externe est liquide.} (La phase SKS existe bien, mais le « K » est une portion en \emph{onde P} dans le noyau : S est convertie en P à l'entrée, puis reconvertie en S à la sortie.)

\textbf{4. Phase PKIKP.} La phase \textbf{PKIKP} est une onde P qui traverse le manteau (P), le noyau externe (K), le noyau interne (I), puis ressort (K puis P). Elle traverse le noyau externe liquide \emph{sous forme d'onde de compression P}, qui, elle, se propage dans les fluides ($v_P = \sqrt{K/\rho} \neq 0$ car $K \neq 0$). Son observation (et son avance systématique) prouve l'existence du \textbf{noyau interne solide}.
\end{corrige}

\begin{corrige}[Exercice 9 --- Sujet type Probatoire : Détermination de la profondeur du noyau]
\textbf{Barème indicatif (sur 20) :} Q1 : 3 pts ; Q2 : 6 pts (a : 3, b : 1, c : 2) ; Q3 : 4 pts ; Q4 : 4 pts ; Q5 : 3 pts.

\textbf{1. Représentation graphique.} Les deux courbes sont croissantes et concaves vers le bas. $t_{PcP}$ est toujours supérieur à $t_P$ sur l'intervalle $[30^\circ, 90^\circ]$ ; l'écart $t_{PcP} - t_P$ \emph{diminue} d'abord jusqu'à un minimum vers $60^\circ$ (170 s), puis \emph{augmente} (valeurs : 320, 230, 180, 170, 190, 230, 290 s). Ce minimum traduit la géométrie de la branche de réflexion.

\textbf{2. Méthode des temps réduits.}
\textbf{(a)} Conversion : $x = \Delta \pi/180$ ; distance d'arc : $R_T x = 111{,}19 \times \Delta^\circ$ km ; terme de réduction : $\dfrac{111{,}19\,\Delta^\circ}{11{,}1} \approx 10{,}02\,\Delta^\circ$ s.

\begin{center}
\begin{tabular}{|c|ccccccc|}\hline
$\Delta$ ($^\circ$) & 30 & 40 & 50 & 60 & 70 & 80 & 90 \\ \hline
$R_T x / v_{\text{ref}}$ (s) & 300{,}5 & 400{,}7 & 500{,}9 & 601{,}1 & 701{,}3 & 801{,}4 & 901{,}6 \\
$t_{PcP}$ (s) & 850 & 950 & 1090 & 1260 & 1450 & 1650 & 1860 \\ \hline
$t_{red}$ (s) & 549{,}5 & 549{,}3 & 589{,}1 & 658{,}9 & 748{,}7 & 848{,}6 & 958{,}4 \\ \hline
\end{tabular}
\end{center}

\textbf{(b)} La courbe $t_{red}(\Delta)$ passe par un \textbf{minimum} vers $\Delta \approx 35$--$40^\circ$ ($t_{red} \approx 549$ s) puis croît fortement. Extrapolée vers $\Delta = 0$, elle tend vers une valeur $t_{red}(0) \approx 660$ s (donnée de l'énoncé, cohérente avec l'ordonnée à l'origine de la branche de réflexion).

\textbf{(c)} Pour une réflexion, le temps réduit s'écrit $t_{red} = t - p\Delta$ avec $p = dt/d\Delta$. Quand $\Delta \to 0$, le rayon devient vertical ($p \to 0$) et le trajet est un aller-retour radial de la surface au réflecteur de rayon $R_c$ :
\[ t_{red}(0) = t(0) = 2 \int_{R_c}^{R_T} \frac{dr}{v_P(r)} = 2 \int_0^{R_T - R_c} \frac{dz}{v_P(z)}. \]
C'est bien le temps aller-retour vertical dans le manteau au-dessus du noyau. \textbf{Point de vigilance :} citer la condition $\Delta \to 0 \Rightarrow$ incidence normale $\Rightarrow$ trajet radial.

\textbf{3. Profondeur du noyau.} Avec $\bar v_P = 10{,}5$ km/s et $t_{red}(0) = 660$ s :
\[ R_T - R_c = \frac{\bar v_P \times t_{red}(0)}{2} = \frac{10{,}5 \times 660}{2} = 3465 \ \text{km}. \]
\[ R_c = 6371 - 3465 = 2906 \ \text{km} \ ; \qquad z_c = R_T - R_c \approx 3465 \ \text{km}. \]
(Valeur de référence : $R_c = 3480$ km, $z_c = 2891$ km ; l'écart provient de la vitesse moyenne supposée constante.)

\textbf{4. Validation par la zone d'ombre.}
\[ \sin i_c^{\text{crit}} = \frac{8{,}1}{13{,}7} = 0{,}591 \Rightarrow i_c^{\text{crit}} \approx 36{,}2^\circ. \]
\[ \Delta_{\min} = 2\left(\frac{\pi}{2} - i_c^{\text{crit}}\right) = 2 \times (90 - 36{,}2) = 107{,}6^\circ \approx 103^\circ \text{ (ordre de grandeur).} \]
La cohérence est bonne (à quelques degrés près, dus au gradient réel de vitesse dans le manteau) : le modèle à deux couches avec un noyau de rayon $\approx 2900$--3480 km reproduit bien le début de la zone d'ombre.

\textbf{5. Conclusion rédigée.} « Les données sismiques permettent de localiser la discontinuité de Gutenberg à une profondeur de \textbf{$\approx 2900$ km} (3465 km avec nos approximations), confirmant l'existence d'un noyau de rayon \textbf{$\approx 3480$ km}. Le saut de vitesse P (chute brutale de 13,7 à 8,1 km/s) et l'absence d'ondes S transmises prouvent que le noyau externe est \textbf{liquide}. »
\end{corrige}

\begin{attention}[Points de vigilance au Probatoire]
\begin{itemize}
    \item Convertir les degrés en radians avant tout calcul d'arc ($x = \Delta \pi/180$).
    \item Justifier chaque étape : loi de Snell-Descartes citée, conditions d'application (milieu stratifié, vitesse constante) explicitées.
    \item Ne pas confondre rayon du noyau $R_c$ et profondeur $z_c = R_T - R_c$.
\end{itemize}
\end{attention}

\begin{corrige}[Exercice 10 --- Sujet type Probatoire : Analyse du modèle PREM]
\textbf{Barème indicatif (sur 20) :} Q1 : 6 pts ; Q2 : 6 pts ; Q3 : 3 pts ; Q4 : 5 pts.

\textbf{1. Analyse des discontinuités.}
\textbf{(a)} Deux discontinuités majeures :
\begin{itemize}
    \item \textbf{Discontinuité de Gutenberg} (manteau / noyau externe) à \textbf{2891 km} de profondeur ;
    \item \textbf{Discontinuité de Lehmann} (noyau externe / noyau interne) à \textbf{5150 km} de profondeur.
\end{itemize}
\textbf{(b)} Contrastes relatifs :
\[ \text{Gutenberg : } \frac{\Delta v_P}{v_P} = \frac{8{,}06 - 13{,}72}{13{,}72} = \frac{-5{,}66}{13{,}72} \approx -41\%. \]
\[ \text{Lehmann : } \frac{\Delta v_P}{v_P} = \frac{11{,}00 - 10{,}35}{10{,}35} \approx +6{,}3\%. \]
La discontinuité de \textbf{Gutenberg} est de loin la plus marquée pour $v_P$ (chute de plus de 40 \%).
\textbf{(c)} $v_S = 0$ dans tout le noyau externe signifie que les ondes de cisaillement ne s'y propagent pas. Or $v_S = \sqrt{\mu/\rho}$, donc $v_S = 0 \Rightarrow \mu = 0$ : le \textbf{module de cisaillement est nul}, propriété caractéristique d'un \textbf{fluide} (liquide) : le noyau externe est liquide. À 5150 km, $v_S$ redevient non nulle ($3,5$ km/s) : le noyau interne est \textbf{solide}.

\textbf{2. Interprétation pétrologique et physique.}
\textbf{(a)} On a $v_P = \sqrt{\dfrac{K + \frac{4}{3}\mu}{\rho}}$. Au passage manteau $\to$ noyau externe, $\mu$ s'annule (fusion) : le numérateur passe de $K + \frac{4}{3}\mu$ à $K$ seul. Même si $\rho$ augmente fortement (5,5 $\to$ 9,9 g/cm$^3$) et même si $K$ augmente (le fer liquide sous forte pression est très incompressible), la \textbf{perte du terme de cisaillement} $\frac{4}{3}\mu$ domine : $v_P$ chute. On en déduit que $K$ augmente à travers la discontinuité, mais moins vite que ne disparaît la contribution de $\mu$ ; le rapport $(K + \frac{4}{3}\mu)/\rho$ diminue globalement.
\textbf{(b)} À 5150 km, le fer-nickel \textbf{cristallise} (la pression croissante impose la solidification malgré la température élevée) : le milieu redevient solide, donc $\mu > 0$, d'où la réapparition de $v_S \approx 3,5$ km/s et le saut de $v_P$ (ajout du terme $\frac{4}{3}\mu$ au numérateur). C'est la preuve sismologique du \textbf{noyau interne solide} (Lehmann, 1936).

\textbf{3. Anisotropie du noyau interne.} Hypothèse structurale : les cristaux de fer (structure hexagonale compacte, $\varepsilon$-Fe) du noyau interne présentent une \textbf{orientation préférentielle} (texture cristalline) avec l'axe cristallographique rapide aligné selon l'axe de rotation terrestre. La vitesse des ondes P étant plus élevée le long de l'axe cristallin, les ondes PKIKP polaires arrivent 3 à 4 s plus tôt que les ondes équatoriales : c'est une \textbf{anisotropie sismique} d'origine cristalline (acquisition possible par convection interne ou croissance préférentielle de la graine).

\textbf{4. Synthèse (paragraphe argumenté).} La sismologie repose sur l'exploitation des temps de parcours des ondes P et S émises par les séismes et enregistrées par un réseau mondial de stations. La loi de Snell-Descartes et la conservation du paramètre de rayon permettent de remonter des courbes $t(\Delta)$ au profil de vitesses $v(r)$. Les réflexions (PcP, ScS) et les temps réduits localisent les discontinuités : Moho, Gutenberg (2891 km), Lehmann (5150 km). La zone d'ombre totale des ondes S au-delà de $103^\circ$ prouve que le noyau externe est liquide ($\mu = 0$), tandis que la phase PKIKP et la réapparition de $v_S$ à 5150 km démontrent l'existence d'un noyau interne solide. La chute de $v_P$ à la frontière noyau-manteau, malgré l'augmentation de densité, s'explique par l'annulation du module de cisaillement. Ainsi, l'ensemble des données sismiques (vitesses, zones d'ombre, phases réfléchies et réfractées) a permis de construire le modèle en couches : croûte, manteau solide, noyau externe liquide, noyau interne solide --- modèle que précisent aujourd'hui les modèles de référence comme PREM.
\end{corrige}

\begin{corrige}[Exercice 11 --- Sujet type Probatoire : Comparaison Terre-Lune et champ magnétique]
\textbf{Barème indicatif (sur 20) :} Q1 : 5 pts ; Q2 : 6 pts ; Q3 : 4 pts ; Q4 : 5 pts.

\textbf{1. Structure en couches.}
\textbf{(a)} Rapports rayon du noyau / rayon de la planète :
\[ \text{Terre : } \frac{3480}{6371} \approx 0{,}55 \ ; \qquad \text{Lune : } \frac{375}{1737} \approx 0{,}22. \]
Le noyau terrestre occupe plus de la moitié du rayon planétaire (et $\approx 16\%$ du volume, soit $(0{,}55)^3 \approx 0{,}166$), alors que le noyau lunaire est proportionnellement très petit ($\approx 1\%$ du volume). Cela explique la \textbf{densité moyenne} beaucoup plus faible de la Lune (3,34 vs 5,51 g/cm$^3$) : la Lune est pauvre en fer, ce qui soutient l'hypothèse de sa formation par \textbf{impact géant} (arrachement de matériau du manteau terrestre, déjà différencié).
\textbf{(b)} La croûte lunaire épaisse (40--80 km, soit jusqu'à $\approx 4,5\%$ du rayon contre $\approx 0,5\%$ pour la Terre) s'explique par la cristallisation précoce d'un \textbf{océan de magma} global : le plagioclase (anorthosite), moins dense, a \textbf{flotté} et formé une croûte épaisse et primitive. L'absence de \textbf{tectonique des plaques} (petite taille $\Rightarrow$ refroidissement rapide, lithosphère rigide unique) a préservé cette croûte jusqu'à aujourd'hui, contrairement à la Terre où la croûte est recyclée en permanence.

\textbf{2. Noyau et champ magnétique.}
\textbf{(a)} Les trois conditions d'une \textbf{dynamo planétaire} :
\begin{enumerate}
    \item une région de \textbf{fluide conducteur} d'électricité (fer liquide) ;
    \item des \textbf{mouvements de convection} vigoureux dans ce fluide (source d'énergie : thermique ou compositionnelle) ;
    \item une \textbf{rotation} suffisamment rapide de la planète (organisation des écoulements par la force de Coriolis).
\end{enumerate}
\textbf{(b)} La Lune remplit la condition 1 (petit noyau de fer) mais plus les conditions 2 et 3 : sa \textbf{petite taille} a entraîné un \textbf{refroidissement rapide} (rapport surface/volume élevé), figeant la convection du noyau ; sa \textbf{rotation lente} (27,3 jours) est défavorable. La dynamo s'est donc \textbf{arrêtée} il y a environ 3--4 Ga ; seule subsiste la \textbf{magnétisation rémanente} des roches crustales, témoin du champ ancien.
\textbf{(c)} Oui, c'est compatible : un noyau liquide \emph{seul} ne suffit pas --- sans convection vigoureuse ni rotation rapide, les mouvements du fer liquide sont incapables d'auto-entretenir un champ magnétique. La petite taille du noyau lunaire réduit en outre la réserve d'énergie thermique et la puissance d'une éventuelle dynamo.

\textbf{3. Sismologie lunaire vs terrestre.}
\textbf{(a)} L'atténuation des ondes S est liée aux défauts, à l'eau dissoute et à la température (proximité du point de fusion). Le manteau lunaire est \textbf{sec} (pas d'eau) et \textbf{froid} (loin du solidus) : les mécanismes de relaxation anélastique sont quasi inexistants, d'où $Q_S \approx 10^4$ (très faible atténuation --- la Lune « sonne comme une cloche »). Sur Terre, l'eau et la chaleur du manteau abaissent $Q_S$ à 100--300.
\textbf{(b)} Difficultés lunaires : réseau de seulement \textbf{4 stations} (géométrie limitée), séismes \textbf{peu nombreux et faibles} ($M < 5$), sources profondes mal réparties, durée de fonctionnement limitée (1969--1977). Il y a donc peu de rayons échantillonnant le noyau, d'où une grande incertitude sur sa taille et son état --- contrairement à la Terre, couverte par des milliers de stations et des milliers de séismes.

\textbf{4. Conclusion (synthèse).} Terre et Lune illustrent le rôle fondamental de la \textbf{taille} dans l'évolution planétaire. La Terre, massive, a conservé une chaleur interne importante (radioactivité, chaleur primordiale), qui entretient la convection du manteau (tectonique des plaques) et celle du noyau externe liquide (dynamo magnétique). La Lune, dix fois plus petite en rayon, s'est refroidie rapidement : sa lithosphère unique et épaisse interdit toute tectonique, sa croûte anorthositique primitive est figée, et sa dynamo s'est éteinte il y a plus de 3 Ga. L'\textbf{eau}, présente sur Terre et absente sur la Lune, lubrifie la tectonique (abaissement du solidus, hydratation du manteau) et rend possible l'atténuation sismique observée. Ainsi, taille, chaleur interne et eau contrôlent conjointement la dynamique interne d'une planète, son champ magnétique protecteur et, in fine, son habitabilité.
\end{corrige}

\newpage

\coursec{Fiche bilan}

\begin{popfigure}
\centering
\begin{center}\tcbox[colback=poporange,colframe=poporange,coltext=white,arc=9pt,boxsep=4pt,left=6pt,right=6pt]{\bfseries\small Structure interne de la Terre par Sismologie}\end{center}
\vspace{-2pt}
\begin{tcbraster}[raster columns=3,raster equal height=rows,raster column skip=6pt,raster row skip=6pt,enhanced,blankest]
\begin{tcolorbox}[enhanced,colback=poporange,colframe=poporange,coltext=white,boxrule=0.9pt,arc=3pt,left=4pt,right=4pt,top=3pt,bottom=3pt,boxsep=1pt,halign=left]\bfseries\scriptsize 1. Ondes sismiques P et S\end{tcolorbox}
\begin{tcolorbox}[enhanced,colback=popgreen,colframe=popgreen,coltext=white,boxrule=0.9pt,arc=3pt,left=4pt,right=4pt,top=3pt,bottom=3pt,boxsep=1pt,halign=left]\bfseries\scriptsize 2. Comportement aux discontinuités\end{tcolorbox}
\begin{tcolorbox}[enhanced,colback=poppurple,colframe=poppurple,coltext=white,boxrule=0.9pt,arc=3pt,left=4pt,right=4pt,top=3pt,bottom=3pt,boxsep=1pt,halign=left]\bfseries\scriptsize 3. Courbes de temps de parcours\end{tcolorbox}
\begin{tcolorbox}[enhanced,colback=poppink,colframe=poppink,coltext=white,boxrule=0.9pt,arc=3pt,left=4pt,right=4pt,top=3pt,bottom=3pt,boxsep=1pt,halign=left]\bfseries\scriptsize 4. Preuve du noyau liquide / solide\end{tcolorbox}
\begin{tcolorbox}[enhanced,colback=popgold,colframe=popgold,coltext=white,boxrule=0.9pt,arc=3pt,left=4pt,right=4pt,top=3pt,bottom=3pt,boxsep=1pt,halign=left]\bfseries\scriptsize 5. Modèle de référence (PREM)\end{tcolorbox}
\begin{tcolorbox}[enhanced,colback=popblue,colframe=popblue,coltext=white,boxrule=0.9pt,arc=3pt,left=4pt,right=4pt,top=3pt,bottom=3pt,boxsep=1pt,halign=left]\bfseries\scriptsize 6. Limites : hétérogénéités\end{tcolorbox}
\end{tcbraster}

\legende{Carte mentale de la leçon NCT09 : des ondes sismiques au modèle structuré de la Terre.}
\end{popfigure}

\begin{bilan}

\courssub{Définitions clés à maîtriser}

\begin{itemize}
    \item \cle{Onde P (Primaire)} : Onde de compression-dilatation (longitudinale), se propage dans les solides \textbf{et} les liquides. Vitesse $v_P = \sqrt{\frac{K + \frac{4}{3}\mu}{\rho}}$.
    \item \cle{Onde S (Secondaire)} : Onde de cisaillement (transversale), se propage \textbf{uniquement dans les solides} ($\mu > 0$). Vitesse $v_S = \sqrt{\frac{\mu}{\rho}}$.
    \item \cle{Discontinuité sismique} : Surface de séparation entre deux milieux où la vitesse des ondes change brutalement (changement de $\rho$, $K$, $\mu$).
    \item \cle{Moho (Mohorovičić)} : Discontinuité Croûte / Manteau (saut de $v_P$ de $\sim 6,5$ à $\sim 8,0$ km/s).
    \item \cle{Discontinuité de Gutenberg} : Limite Manteau / Noyau externe (disparition brutale des ondes S, chute de $v_P$).
    \item \cle{Discontinuité de Lehmann} : Limite Noyau externe (liquide) / Noyau interne (solide) (réapparition des ondes S, augmentation de $v_P$).
    \item \cle{Temps de parcours} $t(d)$ : Temps mis par une onde pour aller du foyer à une station à distance épicentrale $d$ (ou $\Delta$).
    \item \cle{Modèle PREM / AK135} : Modèles de référence 1D (vitesse, densité, pression en fonction de la profondeur).
\end{itemize}

\courssub{Formules et principes fondamentaux (à connaître par cœur)}

\begin{center}
\begin{tabularx}{\linewidth}{|l|X|X|}\hline
\rowcolor{popblueL}
\textbf{Concept} & \textbf{Expression mathématique} & \textbf{Signification / Condition d'application} \\ \hline
Loi de Snell-Descartes (sismique) & $\dfrac{\sin i_1}{v_1} = \dfrac{\sin i_2}{v_2} = p$ (cste) & $i$ : angle d'incidence (par rapport à la normale), $v$ : vitesse dans le milieu, $p$ : paramètre de rayon. \\ \hline
Réflexion totale & $\sin i_c = \dfrac{v_1}{v_2}$ ($v_2 > v_1$) & Au-delà de l'angle critique $i_c$, l'onde est totalement réfléchie. \\ \hline
Temps de parcours (rayon) & $t = \displaystyle\int_{\text{chemin}} \dfrac{ds}{v(s)}$ & Intégrale le long de la trajectoire courbe du rayon. \\ \hline
Distance épicentrale & $\Delta = \displaystyle\int \dfrac{p\,v(r)}{\sqrt{1-p^2v^2(r)}}\,\dfrac{dr}{r}$ & Relation $\Delta(p)$ pour un modèle sphérique stratifié. \\ \hline
Vitesses P et S & $v_P = \sqrt{\dfrac{K+\frac{4}{3}\mu}{\rho}},\quad v_S = \sqrt{\dfrac{\mu}{\rho}}$ & $K$ module d'incompressibilité, $\mu$ module de cisaillement, $\rho$ densité. \\ \hline
Critère liquide / solide & $v_S = 0 \iff \mu = 0$ & L'absence d'ondes S traversant une couche prouve son état liquide. \\ \hline
\end{tabularx}
\end{center}

\courssub{Méthodes express (savoir-faire probatoire)}

\begin{enumerate}
    \item \textbf{Lire une courbe $t(\Delta)$ :} Identifier les branches directes (P, S), les branches réfléchies (PP, SS, PcP, ScS) et les \emph{discontinuités de pente} (changement de vitesse) ou \emph{cassures} (nouvelle phase arrivée).
    \item \textbf{Localiser une discontinuité (Méthode de Herglotz-Wiechert / inversion) :} Mesurer le saut de temps $\Delta t$ ou la distance critique $\Delta_c$ où une phase réfléchie (ex: PcP) arrive. Utiliser la loi de Snell pour remonter au rayon $r$ de la discontinuité : $r = R_{\oplus} \sin i_c$ (pour une couche superficielle) ou inversion intégrale pour la profondeur.
    \item \textbf{Prouver l'état du noyau :}
    \begin{itemize}
        \item Zone d'ombre des ondes S au-delà de $\sim 103^\circ$ $\rightarrow$ Noyau externe \textbf{liquide} ($\mu=0$).
        \item Zone d'ombre des ondes P entre $\sim 103^\circ$ et $\sim 143^\circ$ $\rightarrow$ Réfraction forte dans le noyau (chute de $v_P$).
        \item Arrivée d'ondes PKIKP (traversant le centre) et PKJKP (réfléchi sur graine) $\rightarrow$ Noyau interne \textbf{solide} ($\mu>0$).
    \end{itemize}
    \item \textbf{Construire un modèle 1D simplifié :} Tracer les cercles concentriques (Croûte 0-30 km, Manteau 30-2900 km, Noyau externe 2900-5150 km, Noyau interne 5150-6371 km) en annotant vitesses $v_P, v_S$, densité $\rho$, état physique.
\end{enumerate}

\end{bilan}

\begin{aretenir}[Checklist avant le Probatoire]
\begin{itemize}
    \item $\square$ Je sais définir une onde P et une onde S, donner leurs vitesses respectives et dire dans quels milieux elles se propagent.
    \item $\square$ J'applique la loi de Snell-Descartes sous forme $\frac{\sin i}{v} = cste$ pour expliquer la courbure des rayons dans un milieu à gradient de vitesse.
    \item $\square$ J'identifie sur une courbe $t(\Delta)$ les branches P, S, PcP, ScS, PKP, PKIKP et j'explique leur origine (réflexion, transmission, conversion).
    \item $\square$ J'explique pourquoi l'existence d'une \emph{zone d'ombre des ondes S} prouve que le noyau externe est liquide.
    \item $\square$ J'explique pourquoi l'existence d'une \emph{zone d'ombre des ondes P} (103°-143°) prouve une forte diminution de $v_P$ à l'entrée du noyau.
    \item $\square$ Je cite les trois discontinuités majeures (Moho, Gutenberg, Lehmann) avec leurs profondeurs approximatives et le saut de vitesse associé.
    \item $\square$ Je décris la structure en couches de la Terre (Croûte, Manteau, Noyau externe, Noyau interne) : épaisseur, état physique, composition dominante, variation de $v_P, v_S, \rho$.
    \item $\square$ Je sais calculer la distance épicentrale $\Delta$ (en degrés) à partir du temps de parcours d'une phase de référence (ex: P ou S) en utilisant les tables de temps de parcours (Jeffreys-Bullen ou AK135).
    \item $\square$ Je distingue le modèle 1D (PREM) de la réalité 3D : hétérogénéités latérales (panaches, dalles), anisotropie (manteau supérieur, noyau interne), discontinuité $D''$.
    \item $\square$ Je relie la convection mantellique (dynamique) aux anomalies de vitesse sismique (tomographie) : chaud = lent, froid = rapide.
    \item $\square$ Je sais rédiger une démonstration structurée : Observation (courbes) $\rightarrow$ Hypothèse (modèle) $\rightarrow$ Prédiction (calcul de courbes théoriques) $\rightarrow$ Comparaison/Validation.
\end{itemize}
\end{aretenir}

\findecours

\end{document}
